% Leçon 14 — Grammaire : Le compl. de résultat ; fraction avec 分之、百分之、三分之 — Chinois 1ère A — Cours signé OnBuch+
% Programme officiel MINESEC. Compiler avec : tectonic ZHA14-grammaire-le-compl-de-resultat-fraction-avec.tex
\newcommand{\DOCMATIERE}{Chinois}
\newcommand{\DOCNIVEAU}{1ère A}
\newcommand{\DOCMODULE}{Module 2 : Environnement et santé}
\newcommand{\DOCLECON}{ZHA14}
\newcommand{\DOCTITRE}{Leçon 14 — Grammaire : Le compl. de résultat ; fraction avec 分之、百分之、三分之}
% OnBuch+ — préambule « cours signés » ENRICHI (compilé avec tectonic / XeLaTeX)
% Système visuel « Pop » d'OnBuch+ : fond crème (jamais blanc), bordures encre
% épaisses, ombres « dures » décalées, titres Archivo Black en capitales.
% Version MATHÉMATIQUES : graphes (pgfplots/tikz), boîtes pédagogiques
% (définition, propriété, méthode, attention, le savais-tu, exercice résolu,
% exercice, TP, bilan), page de garde et signature OnBuch+ ancrée en bas.
%
% Macros attendues dans main.tex AVANT \input{preamble} :
%   \DOCMATIERE  \DOCNIVEAU  \DOCMODULE  \DOCLECON (numéro)  \DOCTITRE
\documentclass[11pt]{article}

\usepackage{fontspec}
\usepackage[french]{babel} % français = langue principale ; le chinois via \zh{..} / textebox (xeCJK)
\usepackage{xeCJK}
\usepackage{pifont} % \ding{51} \ding{55} utilisés par les modèles
\usepackage[a4paper,margin=2cm,top=3.4cm,bottom=2.8cm,headheight=1.4cm,footskip=1.3cm]{geometry}
\usepackage{amsmath,amssymb,amsfonts}
\usepackage{mathtools} % \xmapsto, \coloneqq... souvent utilisés par les modèles
\usepackage{cancel} % simplifications barrées (bilans de forces)
% \abs{z} (module, valeur absolue) et \norm{u} (norme), utilisés par les modèles
\providecommand{\abs}{}\let\abs\relax\DeclarePairedDelimiter{\abs}{\lvert}{\rvert}
\providecommand{\norm}{}\let\norm\relax\DeclarePairedDelimiter{\norm}{\lVert}{\rVert}
\usepackage[table]{xcolor} % [table] : fournit aussi \rowcolors (alternance de lignes)
\usepackage{pagecolor}
\usepackage{tikz}
\usetikzlibrary{arrows.meta,positioning,calc,shapes.geometric,shapes.misc,shapes.arrows,decorations.pathmorphing,decorations.pathreplacing,decorations.markings,patterns,shadows,mindmap,trees,fit,backgrounds,matrix,intersections,angles,quotes}
\usepackage{pgfplots}
\usepackage[version=4]{mhchem} % équations nucléaires \ce{^{235}_{92}U}
\usepackage[european,siunitx]{circuitikz} % schémas électriques
\pgfplotsset{compat=1.18}
\usepgfplotslibrary{fillbetween,groupplots} % aires entre courbes (calcul intégral)
\usepackage{enumitem}
\usepackage{fancyhdr}
% [most] charge toutes les bibliothèques tcolorbox (listings, minted, raster,
% documentation...) et gonfle inutilement la mémoire TeX sur un document long
% (~300 boîtes) : on ne charge que ce qui est réellement utilisé.
\usepackage[skins,breakable]{tcolorbox}
\usepackage{graphicx}
\usepackage{colortbl}
\usepackage{booktabs}
\usepackage{tabularx}
\usepackage{diagbox} % case d'en-tête diagonale des tableaux à double entrée
% Colonnes extensibles centrée / gauche / droite (C, L, R), souvent utilisées par les modèles
\newcolumntype{C}{>{\centering\arraybackslash}X}
\newcolumntype{L}{>{\raggedright\arraybackslash}X}
\newcolumntype{R}{>{\raggedleft\arraybackslash}X}
\usepackage{array}
\usepackage{multicol}
\usepackage{multirow}
\usepackage{siunitx}
% parse-numbers=false : accepte aussi \SI{2,5 \times 10^{-3}}{...}
\sisetup{locale=FR,output-decimal-marker={,},group-minimum-digits=4,parse-numbers=false}
\DeclareSIUnit\ohm{\ensuremath{\symup{\Omega}}} % U+2126 absent de la police
\DeclareSIUnit\division{div} % réglages d'oscilloscope (ms/div, V/div)
\DeclareSIUnit\tr{tr} % tours (vitesse de rotation, ex. \SI{1500}{\tr\per\minute})
\DeclareSIUnit\molar{M} % concentration molaire (ex. \SI{3}{\molar}, \SI{1}{\micro\molar})
\DeclareSIUnit\an{an} % année (le modèle confond parfois avec \year, un compteur TeX) — ex. \SI{5}{\kilo\meter\per\an}
\DeclareSIUnit\FCFA{FCFA} % franc CFA (ex. \SI{500}{\FCFA\per\kilo\gram})
\DeclareSIUnit\degreeBrix{\textdegree Brix} % teneur en sucre d'un sirop/jus (réfractométrie)
\DeclareSIUnit\month{mois} % le modèle utilise parfois \month, un compteur TeX, comme unité de durée
\usepackage{qrcode}
% \degree : commande fréquemment utilisée par les modèles pour un angle ou une
% température en dehors de \SI{}{} (non fournie par les paquets chargés).
\providecommand{\degree}{^{\circ}}
% \qed (fin de démonstration), utilisé par les modèles sans amsthm
\providecommand{\qed}{\hfill$\square$}
% \barre : double barre des valeurs interdites dans un tableau de variations (tkz-tab)
\providecommand{\barre}{\ensuremath{\|}}
% environnement proof (amsthm non chargé) utilisé par les modèles
\ifdefined\proof\else\newenvironment{proof}[1][Démonstration]{\par\noindent\textit{#1.}\ }{\qed\par}\fi
\usepackage{lastpage}
\usepackage{hyperref}
\hypersetup{hidelinks}

% ============================================================
% Palette « Pop » OnBuch+ — mêmes hex que la classe P (plus_ui.dart)
% ============================================================
\definecolor{poporange}{HTML}{F59321}
\definecolor{popdark}{HTML}{E07A0C}
\definecolor{popcream}{HTML}{FFF6E8}
\definecolor{popink}{HTML}{1C1714}
\definecolor{poppurple}{HTML}{7A5AE0}
\definecolor{popgreen}{HTML}{2E9E6A}
\definecolor{poppink}{HTML}{E0457B}
\definecolor{popblue}{HTML}{2F7DE1}
\definecolor{popgold}{HTML}{C98A12}
\definecolor{popmuted}{HTML}{8A7F76}
\definecolor{popline}{HTML}{EADFD2}
\definecolor{popgray}{HTML}{8A7F76}
\colorlet{popgrey}{popgray}
% Alias tolérants (le modèle nomme parfois les couleurs différemment)
\colorlet{popwhite}{white}
\colorlet{popblack}{popink}
\colorlet{popred}{poppink}
\colorlet{popyellow}{popgold}
% Teintes claires (fonds de tableaux/encadrés) — le modèle nomme parfois
% les couleurs avec un suffixe L (ex. popblueL) sans les avoir définies.
\colorlet{poporangeL}{poporange!20}
\colorlet{popdarkL}{popdark!20}
\colorlet{poppurpleL}{poppurple!20}
\colorlet{popgreenL}{popgreen!20}
\colorlet{poppinkL}{poppink!20}
\colorlet{popblueL}{popblue!20}
\colorlet{popgoldL}{popgold!20}
\colorlet{popredL}{popred!20}
\colorlet{popgrayL}{popgray!20}
% Teintes claires pour fonds de tableaux / zones
\colorlet{popblueL}{popblue!10!white}
\colorlet{poporangeL}{poporange!14!white}
\colorlet{popgreenL}{popgreen!12!white}
\colorlet{poppurpleL}{poppurple!12!white}
\colorlet{poppinkL}{poppink!12!white}
\colorlet{popgoldL}{popgold!14!white}

\pagecolor{popcream}

% ============================================================
% Typographie
% ============================================================
\defaultfontfeatures{Ligatures=TeX}
\setmainfont{PlusJakartaSans-Medium.ttf}[Path=../../../fonts/,BoldFont=PlusJakartaSans-Bold.ttf]
% Symboles, API et emojis absents de la police principale : repli sur DejaVu Sans (caractères actifs)
\newfontface\symfont{DejaVuSans.ttf}[Path=../../../fonts/]
\makeatletter
\newcommand\symactive[1]{\catcode#1=13 \begingroup\lccode`\~=#1 \lowercase{\endgroup\def~}{{\symfont\char#1}}}
\newcommand\symblank[1]{\catcode#1=13 \begingroup\lccode`\~=#1 \lowercase{\endgroup\def~}{}}
\newcommand\symhyph[1]{\catcode#1=13 \begingroup\lccode`\~=#1 \lowercase{\endgroup\def~}{\mbox{-}}}
\newcount\symc
\newcommand\symrange[2]{\symc=#1 \@whilenum\symc<#2 \do{\expandafter\symactive\expandafter{\the\symc}\advance\symc by 1 }}
\newcommand\symblankrange[2]{\symc=#1 \@whilenum\symc<#2 \do{\expandafter\symblank\expandafter{\the\symc}\advance\symc by 1 }}
\makeatother
\symrange{"0250}{"02FF}\symactive{"2713}\symactive{"2714}\symactive{"2717}\symactive{"2718}\symactive{"2610}\symactive{"2611}\symactive{"2612}
\symrange{"2600}{"27BF}
\catcode"2705=13 \begingroup\lccode`\~="2705 \lowercase{\endgroup\def~}{{\symfont\char"2713}}\symblank{"26BD}\symblank{"26A1}
\symrange{"2460}{"257F}\symactive{"223C}\symactive{"2248}\symactive{"2260}
\symhyph{"2011}\symhyph{"2E1A}\symblankrange{"1F300}{"1FAFF}\symblank{"FE0F}
% Caractères chinois : WenQuanYi Zen Hei (sous-ensemble GB2312 + pinyin), gras/italique simulés
\setCJKmainfont{WenQuanYiZenHei-subset.ttf}[Path=../../../fonts/,AutoFakeBold=3,AutoFakeSlant=0.2]
\xeCJKsetup{CJKspace=false}
% Pinyin : voyelles à caron (ǎ ǐ ǒ ǔ ǖ ǘ ǚ ǜ) absentes de la police principale -> caractères actifs
\newfontface\pyfont{WenQuanYiZenHei-subset.ttf}[Path=../../../fonts/]
\newcommand\pyactive[1]{\catcode#1=13 \begingroup\lccode`\~=#1 \lowercase{\endgroup\def~}{{\pyfont\char#1}}}
\pyactive{"01CD}\pyactive{"01CE}\pyactive{"01CF}\pyactive{"01D0}\pyactive{"01D1}\pyactive{"01D2}\pyactive{"01D3}\pyactive{"01D4}\pyactive{"01D5}\pyactive{"01D6}\pyactive{"01D7}\pyactive{"01D8}\pyactive{"01D9}\pyactive{"01DA}\pyactive{"01DB}\pyactive{"01DC}
% Maths en sans-serif (Fira Math) avec chiffres et lettres droites de Plus
% Jakarta, pour que les formules se fondent dans le texte.
\usepackage[math-style=ISO,bold-style=ISO]{unicode-math}
\setmathfont{FiraMath-Regular.otf}[Path=../../../fonts/]
\setmathfont{PlusJakartaSans-Medium.ttf}[Path=../../../fonts/,range={up/num,up/{Latin,latin},\mathup}]
\usepackage{icomma} % virgule décimale sans espace en mode math
% Caractères mathématiques Unicode absents des polices de texte (Plus Jakarta,
% Archivo Black) : rendus par leur équivalent mathématique, en texte comme en
% formule — sinon ils s'affichent en carré vide (ex. « Arithmétique dans ℤ »).
\usepackage{newunicodechar}
\newunicodechar{ℝ}{\ensuremath{\mathbb{R}}}
\newunicodechar{ℤ}{\ensuremath{\mathbb{Z}}}
\newunicodechar{ℂ}{\ensuremath{\mathbb{C}}}
\newunicodechar{ℕ}{\ensuremath{\mathbb{N}}}
\newunicodechar{ℚ}{\ensuremath{\mathbb{Q}}}
\newunicodechar{𝔻}{\ensuremath{\mathbb{D}}}
\newunicodechar{α}{\ensuremath{\alpha}}
\newunicodechar{β}{\ensuremath{\beta}}
\newunicodechar{γ}{\ensuremath{\gamma}}
\newunicodechar{δ}{\ensuremath{\delta}}
\newunicodechar{ε}{\ensuremath{\varepsilon}}
\newunicodechar{θ}{\ensuremath{\theta}}
\newunicodechar{λ}{\ensuremath{\lambda}}
\newunicodechar{μ}{\ensuremath{\mu}}
\newunicodechar{σ}{\ensuremath{\sigma}}
\newunicodechar{τ}{\ensuremath{\tau}}
\newunicodechar{φ}{\ensuremath{\varphi}}
\newunicodechar{ψ}{\ensuremath{\psi}}
\newunicodechar{ω}{\ensuremath{\omega}}
\newunicodechar{Σ}{\ensuremath{\Sigma}}
% majuscules produites par \MakeUppercase dans les titres (α-aminés -> Α)
\newunicodechar{Α}{\ensuremath{\alpha}}
\newunicodechar{Β}{\ensuremath{\beta}}
\newunicodechar{Γ}{\ensuremath{\Gamma}}
\newunicodechar{Δ}{\ensuremath{\Delta}}
\newunicodechar{Ε}{\ensuremath{\varepsilon}}
\newunicodechar{Θ}{\ensuremath{\Theta}}
\newunicodechar{Λ}{\ensuremath{\Lambda}}
\newunicodechar{Μ}{\ensuremath{\mu}}
\newunicodechar{Τ}{\ensuremath{\tau}}
\newunicodechar{Φ}{\ensuremath{\Phi}}
\newunicodechar{Ψ}{\ensuremath{\Psi}}
\newunicodechar{Ω}{\ensuremath{\Omega}}
\newunicodechar{↦}{\ensuremath{\mapsto}}
\newunicodechar{∈}{\ensuremath{\in}}
\newunicodechar{∉}{\ensuremath{\notin}}
\newunicodechar{∧}{\ensuremath{\wedge}}
\newunicodechar{⇔}{\ensuremath{\Leftrightarrow}}
\newunicodechar{⟺}{\ensuremath{\Longleftrightarrow}}
\newunicodechar{⇒}{\ensuremath{\Rightarrow}}
\newunicodechar{⟹}{\ensuremath{\Longrightarrow}}
\newunicodechar{∘}{\ensuremath{\circ}}
\newunicodechar{∪}{\ensuremath{\cup}}
\newunicodechar{∩}{\ensuremath{\cap}}
\newunicodechar{≡}{\ensuremath{\equiv}}
\newunicodechar{⊂}{\ensuremath{\subset}}
\newunicodechar{⊥}{\ensuremath{\perp}}
\newunicodechar{∃}{\ensuremath{\exists}}
\newunicodechar{∀}{\ensuremath{\forall}}
\newunicodechar{∛}{\ensuremath{\sqrt[3]{\,}}}
\newunicodechar{∜}{\ensuremath{\sqrt[4]{\,}}}
\newunicodechar{⃗}{\ensuremath{{}^{\to}}}
\newunicodechar{ⁿ}{\ifmmode^{n}\else\textsuperscript{\ensuremath{n}}\fi}
\newunicodechar{ˣ}{\ifmmode^{x}\else\textsuperscript{\ensuremath{x}}\fi}
\newunicodechar{ᵇ}{\ifmmode^{b}\else\textsuperscript{\ensuremath{b}}\fi}
\newunicodechar{ʳ}{\ifmmode^{r}\else\textsuperscript{\ensuremath{r}}\fi}
\newunicodechar{ᵘ}{\ifmmode^{u}\else\textsuperscript{\ensuremath{u}}\fi}
\newunicodechar{ʸ}{\ifmmode^{y}\else\textsuperscript{\ensuremath{y}}\fi}
\newunicodechar{ᵃ}{\ifmmode^{a}\else\textsuperscript{\ensuremath{a}}\fi}
\newunicodechar{ᵉ}{\ifmmode^{e}\else\textsuperscript{\ensuremath{e}}\fi}
\newunicodechar{ᵏ}{\ifmmode^{k}\else\textsuperscript{\ensuremath{k}}\fi}
\newunicodechar{ᵖ}{\ifmmode^{p}\else\textsuperscript{\ensuremath{p}}\fi}
\newunicodechar{ᴺ}{\ifmmode^{N}\else\textsuperscript{\ensuremath{N}}\fi}
\newunicodechar{ᶜ}{\ifmmode^{c}\else\textsuperscript{\ensuremath{c}}\fi}
\newunicodechar{ʰ}{\ifmmode^{h}\else\textsuperscript{\ensuremath{h}}\fi}
\newunicodechar{ᵠ}{\ifmmode^{q}\else\textsuperscript{\ensuremath{q}}\fi}
\newunicodechar{ₙ}{\ifmmode_{n}\else\textsubscript{\ensuremath{n}}\fi}
\newunicodechar{ₐ}{\ifmmode_{a}\else\textsubscript{\ensuremath{a}}\fi}
\newunicodechar{ᵢ}{\ifmmode_{i}\else\textsubscript{\ensuremath{i}}\fi}
\newunicodechar{ᵣ}{\ifmmode_{r}\else\textsubscript{\ensuremath{r}}\fi}
\newunicodechar{ₖ}{\ifmmode_{k}\else\textsubscript{\ensuremath{k}}\fi}
\newunicodechar{ᵦ}{\ifmmode_{\beta}\else\textsubscript{\ensuremath{\beta}}\fi}
\newunicodechar{ₓ}{\ifmmode_{x}\else\textsubscript{\ensuremath{x}}\fi}
\newfontfamily\popdisplay{ArchivoBlack-Regular.ttf}[Path=../../../fonts/]
\newfontfamily\poplogo{SpaceGrotesk-Bold.ttf}[Path=../../../fonts/]
\newfontfamily\popsemi{PlusJakartaSans-SemiBold.ttf}[Path=../../../fonts/]
\newfontfamily\popxbold{PlusJakartaSans-ExtraBold.ttf}[Path=../../../fonts/]
\newfontfamily\popxboldls{PlusJakartaSans-ExtraBold.ttf}[Path=../../../fonts/,LetterSpace=3.0]

\setlist[enumerate]{leftmargin=1.7em,itemsep=5pt,topsep=4pt}
\setlist[itemize]{leftmargin=1.7em,itemsep=4pt,topsep=4pt,label={\color{poporange}\textbullet}}
\setlength{\parindent}{0pt}
\setlength{\parskip}{5pt}
\renewcommand{\arraystretch}{1.25}

% pgfplots : style Pop par défaut
\pgfplotsset{
  every axis/.append style={axis line style={popink,line width=1pt},tick style={popink},
    grid style={popline,line width=0.5pt},label style={font=\small},tick label style={font=\footnotesize}},
  popcurve/.style={poporange,line width=1.8pt,no markers,smooth},
}
% Même style disponible dans un \draw TikZ simple (hors pgfplots)
\tikzset{popcurve/.style={poporange,line width=1.8pt}}
\tikzset{popgrid/.style={popline,very thin}}
\colorlet{popcurve}{poporange} % le nom du style est parfois utilisé comme couleur

% ============================================================
% Pill / squiggle
% ============================================================
\newcommand{\poppill}[3]{%
  \tikz[baseline=-0.55ex]\node[draw=popink,line width=1.2pt,rounded corners=2.6pt,%
    fill=#2,inner xsep=6.5pt,inner ysep=3pt]{%
    \popxboldls\fontsize{7}{7}\selectfont\color{#3}\MakeUppercase{#1}};%
}
\newcommand{\popsquiggle}[1]{%
  \tikz[baseline=-0.4ex]{
    \draw[#1,line width=2.6pt,line cap=round]
      plot [smooth] coordinates {(0,0.09) (0.34,0.01) (0.68,0.21) (1.02,0.01) (1.36,0.10)};
  }%
}
% Mot-clé surligné (marqueur orange sous le texte)
\newcommand{\cle}[1]{{\popxbold\color{popdark}#1}}

% ============================================================
% En-tête / pied
% ============================================================
\pagestyle{fancy}
\fancyhf{}
\renewcommand{\headrulewidth}{0pt}
\renewcommand{\footrulewidth}{0pt}
\lhead{\raisebox{-0.15ex}{\poplogo\fontsize{13}{13}\selectfont\color{popink}OnBuch\color{poporange}+}}
\rhead{\poppill{\DOCMATIERE\ \textbullet\ \DOCNIVEAU\ \textbullet\ Leçon \DOCLECON}{white}{popink}}
\lfoot{\popxbold\fontsize{7.6}{7.6}\selectfont\color{popmuted}\MakeUppercase{Cours signé OnBuch+}\ \textbullet\ {\popsemi\DOCTITRE}}
\rfoot{\tikz[baseline=-0.6ex]\node[rounded corners=8.5pt,fill=popink,minimum height=17pt,minimum width=17pt,inner xsep=5pt,inner sep=0]{\popxbold\fontsize{8}{8}\selectfont\color{popcream}\thepage\,/\,\pageref*{LastPage}};}

% ============================================================
% Titres
% ============================================================
\newcommand{\chaptitre}[2]{%
  \begin{tcolorbox}[enhanced,colback=poporange,colframe=popink,boxrule=2.6pt,arc=11pt,
      drop shadow=popink,left=15pt,right=15pt,top=12pt,bottom=11pt,before skip=3pt,after skip=18pt]
    \poppill{#2}{popink}{popcream}\par\vspace{9pt}
    {\raggedright\hyphenpenalty=10000\exhyphenpenalty=10000\popdisplay\fontsize{22}{24}\selectfont\color{popcream}\MakeUppercase{#1}\par}
  \end{tcolorbox}
}
% Section numérotée : pastille orange avec le numéro + titre Archivo
\newcounter{popsec}
\newcommand{\coursec}[1]{%
  \stepcounter{popsec}\setcounter{popsub}{0}%
  \par\vspace{16pt}\noindent
  \begin{minipage}{\linewidth}
  \tikz[baseline=-0.7ex]\node[draw=popink,line width=1.6pt,fill=poporange,rounded corners=4pt,minimum size=22pt,inner sep=0]{\popdisplay\fontsize{12}{12}\selectfont\color{popcream}\thepopsec};%
  \hspace{8pt}{\popdisplay\fontsize{15}{17}\selectfont\color{popink}\MakeUppercase{#1}}\par
  \vspace{4pt}\noindent{\color{poporange}\rule{36pt}{3.6pt}}
  \end{minipage}\par\nopagebreak\vspace{8pt}
}
\newcounter{popsub}
\newcommand{\courssub}[1]{%
  \stepcounter{popsub}%
  \par\vspace{9pt}\noindent{\popxbold\fontsize{11.5}{13}\selectfont\color{popink}{\color{poporange}\thepopsec.\thepopsub}\hspace{6pt}#1}\par\nopagebreak\vspace{4pt}
}

% ============================================================
% Boîtes pédagogiques — même recette : fond blanc, bordure épaisse colorée,
% ombre dure encre, badge plein. Titre optionnel [#1].
% ============================================================
\tcbset{popbase/.style={breakable,enhanced,colback=white,boxrule=2.3pt,arc=9pt,
  shadow={3pt}{-3pt}{0pt}{popink},left=14pt,right=14pt,top=11pt,bottom=11pt,before skip=11pt,after skip=15pt,
  fontupper=\popsemi\fontsize{10.6}{14.5}\selectfont\color{popink},
  fontlower=\popsemi\fontsize{10.6}{14.5}\selectfont\color{popink}}}
% Coche dessinée (le glyphe ✓ n'existe pas dans les polices du document)
\newcommand{\popcheck}[1]{\tikz[baseline=-0.5ex]\draw[#1,line width=1.6pt,line cap=round,line join=round](-0.13,0.0)--(-0.04,-0.09)--(0.13,0.1);}
\providecommand{\coche}{\popcheck{popgreen}}
\providecommand{\resultat}[1]{\par\smallskip\noindent\fcolorbox{popgreen}{popgreenL}{\parbox{\dimexpr\linewidth-2\fboxsep-2\fboxrule\relax}{\textbf{Résultat.} #1}}\par\smallskip}
\newcommand{\popboxhead}[3]{\poppill{#1}{#2}{white}\ifx\relax#3\relax\else\hspace{6pt}{\popxbold\fontsize{10.6}{12}\selectfont\color{popink}#3}\fi\par\vspace{7pt}}

\newtcolorbox{aretenir}[1][]{popbase,colframe=popgreen,before upper={\popboxhead{À retenir}{popgreen}{#1}}}
% Texte en chinois (document de lecture, dialogue, extrait) : cadre
% d'étude. Un titre optionnel (source, auteur) s'affiche dans l'en-tête.
\newtcolorbox{textebox}[1][]{popbase,colframe=popblue,before upper={\popboxhead{文本 · Texte}{popblue}{#1}},after upper={}}
% Marques ✓ / ✗ (forme correcte / fautive) souvent utilisées par les modèles
\providecommand{\cross}{\textcolor{poppink}{\ensuremath{\times}}}
\providecommand{\xmark}{\cross}\providecommand{\crossmark}{\cross}\providecommand{\cmark}{\ensuremath{\checkmark}}
% Ligne de réponse / justification à compléter (inventée par les modèles)
\providecommand{\justification}[1]{\par\noindent\textit{Justification~:} \dotfill\par}
% \textipa (package tipa non chargé) : la police ne couvre pas l'API -> simple passage
\providecommand{\textipa}[1]{#1}
% Soulignement (ulem non chargé) : \uline, \ul
\providecommand{\uline}[1]{\underline{#1}}\providecommand{\ul}[1]{\underline{#1}}
% \glossaire{\item ...} inventée par les modèles : liste de glossaire
\providecommand{\glossaire}[1]{\begin{itemize}#1\end{itemize}}
% \alert (beamer) inventée par les modèles : texte mis en évidence
\providecommand{\alert}[1]{\textbf{\textcolor{poppink}{#1}}}
% variantes ASCII d'ordinaux écrites en macro
\providecommand{\iere}{\textsuperscript{re}}\providecommand{\ieme}{\textsuperscript{e}}\providecommand{\ere}{\textsuperscript{re}}\providecommand{\eme}{\textsuperscript{e}}\providecommand{\er}{\textsuperscript{er}}\providecommand{\ie}{\textsuperscript{e}}
% Ordinaux français écrits en macro par les modèles : 1\ière, 3\ième
\newcommand{\ière}{\textsuperscript{re}}\newcommand{\ième}{\textsuperscript{e}}
% \exemple (inventée par les modèles) : étiquette « Exemple : »
\providecommand{\exemple}{\textbf{Exemple~:}\ }
% \square utilisable aussi hors mode math (cases à cocher)
\AtBeginDocument{\let\squareMath\square\renewcommand{\square}{\ensuremath{\squareMath}}}
% Mot ou phrase en chinois à l'intérieur d'un paragraphe français
\newcommand{\zh}[1]{\ifmmode\text{#1}\else{#1}\fi}
\let\eng\zh \let\esp\zh \let\all\zh % alias tolérés
% flèches d'intonation inventées par les modèles
\providecommand{\textsearrow}{}\renewcommand{\textsearrow}{\ensuremath{\searrow}}\providecommand{\textnearrow}{}\renewcommand{\textnearrow}{\ensuremath{\nearrow}}
% \fr{..} : mot français (inventé par les modèles)
\providecommand{\fr}[1]{\textit{#1}}
% zhbox : encadré de texte chinois inventé par les modèles
\newenvironment{zhbox}{\begin{quote}}{\end{quote}}
% \vs : « versus » inventé par les modèles
\providecommand{\vs}{\textit{vs}\ }
% \blank : case à compléter inventée par les modèles
\providecommand{\blank}[1][]{\underline{\hspace{1.6em}}}
\newtcolorbox{exemplebox}[1][]{popbase,colframe=poporange,before upper={\popboxhead{Exemple}{poporange}{#1}}}
\newtcolorbox{definition}[1][]{popbase,colframe=popblue,before upper={\popboxhead{Définition}{popblue}{#1}}}
\newtcolorbox{propriete}[1][]{popbase,colframe=poppurple,before upper={\popboxhead{Propriété}{poppurple}{#1}}}
\newtcolorbox{methode}[1][]{popbase,colframe=popgold,before upper={\popboxhead{Méthode}{popgold}{#1}}}
\newtcolorbox{attention}[1][]{popbase,colframe=poppink,before upper={\popboxhead{Attention}{poppink}{#1}}}
\newtcolorbox{experience}[1][]{popbase,colframe=popblue,colback=popblueL,before upper={\popboxhead{Expérience}{popblue}{#1}}}
% « Le savais-tu ? » : carte violette pleine (contexte, culture, Cameroun)
\newtcolorbox{savaistu}[1][]{popbase,colback=poppurple,colframe=popink,
  fontupper=\popsemi\fontsize{10.4}{14.2}\selectfont\color{white},
  before upper={\poppill{Le savais-tu\,?}{white}{poppurple}\ifx\relax#1\relax\else\hspace{6pt}{\popxbold\color{white}#1}\fi\par\vspace{7pt}}}
% Exercice résolu : énoncé en haut, \tcblower puis la solution sur fond crème
\newcounter{popexo}\newcounter{popexr}
\newtcolorbox{exoresolu}[1][]{popbase,colframe=poporange,colbacklower=popcream,
  segmentation style={popink,line width=1.2pt,dashed},
  before upper={\stepcounter{popexr}\popboxhead{Exercice résolu \thepopexr}{poporange}{#1}},
  before lower={\poppill{Solution}{popink}{popcream}\par\vspace{6pt}}}
% Exercice (non résolu en place) — la correction est dans la partie Corrigés
\newtcolorbox{exercice}[1][]{popbase,colframe=popink,
  before upper={\stepcounter{popexo}\popboxhead{Exercice \thepopexo}{popink}{#1}}}
% Corrigé
\newtcolorbox{corrige}[1][]{popbase,colframe=popgreen,colback=popgreenL,
  before upper={\popboxhead{Corrigé}{popgreen}{#1}}}
% Objectifs / prérequis / bilan
\newtcolorbox{objectifs}{popbase,colframe=popink,colback=poporangeL,
  before upper={\popboxhead{Objectifs de la leçon}{popink}{}}}
\newtcolorbox{prerequis}{popbase,colframe=popink,colback=popblueL,
  before upper={\popboxhead{Prérequis}{popblue}{}}}
\newtcolorbox{bilan}{popbase,colframe=popink,colback=white,boxrule=2.8pt,
  before upper={\popboxhead{Fiche bilan}{poporange}{L'essentiel à connaître pour l'examen}}}
% Figure encadrée avec légende
\newtcolorbox{popfigure}[1][]{enhanced,breakable=false,colback=white,colframe=popline,boxrule=1.4pt,arc=8pt,
  left=10pt,right=10pt,top=10pt,bottom=8pt,before skip=10pt,after skip=12pt,halign=center,
  lower separated=false,halign lower=center,
  fontlower=\popsemi\fontsize{9}{11.5}\selectfont\color{popmuted},
  #1}
\newcommand{\legende}[1]{\par\vspace{6pt}{\popsemi\fontsize{9}{11.5}\selectfont\color{popmuted}#1}}

% ============================================================
% Signature OnBuch+
% ============================================================
\newcommand{\poplogoinline}[1]{{\poplogo\fontsize{#1}{#1}\selectfont\color{popink}OnBuch\color{poporange}+}}
\newcommand{\signatureonbuch}{%
  \begin{tcolorbox}[enhanced,colback=popink,colframe=popink,boxrule=2.2pt,arc=10pt,
      drop shadow=poporange,left=14pt,right=14pt,top=10pt,bottom=10pt,before skip=0pt,after skip=0pt]
    \tikz[baseline=-0.7ex]\node[circle,fill=popgreen,draw=popcream,line width=1.3pt,minimum size=22pt,inner sep=0]{\popcheck{white}};%
    \hspace{9pt}\begin{minipage}[c]{0.62\linewidth}
      {\popxbold\fontsize{12}{13}\selectfont\color{popcream}Signé\ \textbullet\ {\poplogo OnBuch\color{poporange}+}}\par\vspace{2pt}
      {\raggedright\popsemi\fontsize{8}{10.5}\selectfont\color{popline}\MakeUppercase{Équipe pédagogique OnBuch+}\\\MakeUppercase{Conforme au programme officiel MINESEC}\par}
    \end{minipage}\hfill
    {\poplogo\fontsize{22}{22}\selectfont\color{popcream}OnBuch\color{poporange}+}
  \end{tcolorbox}%
}

% ============================================================
% Page de garde
% #1 titre · #2 matière · #3 niveau · #4 module · #5 n° leçon · #6 durée ·
% #7 description / objectifs (texte ou itemize)
% ============================================================
\newcommand{\pagegarde}[7]{%
  \thispagestyle{empty}
  \newgeometry{a4paper,margin=1.7cm,top=1.6cm,bottom=1.4cm}%
  \begin{tikzpicture}[remember picture,overlay]
    \fill[poporange!18] ([xshift=-3.2cm,yshift=-3.6cm]current page.north east) circle (4.6cm);
    \fill[poppurple!14] ([xshift=2.4cm,yshift=7.6cm]current page.south west) circle (3.8cm);
  \end{tikzpicture}%
  {\poplogo\fontsize{30}{30}\selectfont\color{popink}OnBuch\color{poporange}+}\hfill
  \raisebox{0.6ex}{\poppill{Cours signé \textbullet\ Premium}{popink}{popcream}}\par
  \vspace{1pt}\popsquiggle{poppurple}\par
  \vspace{8pt}
  \begin{tcolorbox}[enhanced,colback=poporange,colframe=popink,boxrule=2.8pt,arc=13pt,
      drop shadow=popink,left=18pt,right=18pt,top=12pt,bottom=13pt,after skip=10pt]
    \poppill{#2 \textbullet\ #3}{popink}{popcream}\hspace{5pt}\poppill{#4}{white}{popink}\par\vspace{8pt}
    {\popdisplay\fontsize{14}{14}\selectfont\color{popink}LEÇON #5}\par\vspace{4pt}
    {\raggedright\hyphenpenalty=10000\exhyphenpenalty=10000\popdisplay\fontsize{25}{27}\selectfont\color{popcream}\MakeUppercase{#1}\par}
    \vspace{8pt}
    {\popxbold\fontsize{9.5}{9.5}\selectfont\color{popink}\MakeUppercase{Durée conseillée : #6}}
  \end{tcolorbox}
  \begin{tcolorbox}[enhanced,colback=white,colframe=popink,boxrule=2.2pt,arc=10pt,
      drop shadow=popink,left=16pt,right=16pt,top=10pt,bottom=10pt,after skip=8pt]
    \poppill{Dans cette leçon}{poppurple}{white}\par\vspace{5pt}
    \popsemi\fontsize{9.3}{12.6}\selectfont\color{popink}#7
  \end{tcolorbox}
  \noindent\begin{minipage}[t]{0.487\linewidth}
    \begin{tcolorbox}[enhanced,colback=popgreen,colframe=popink,boxrule=2.2pt,arc=10pt,
        drop shadow=popink,left=11pt,right=11pt,top=10pt,bottom=10pt,height=3.1cm,valign=center]
      \tcbox[colback=white,colframe=popink,boxrule=1.3pt,arc=5pt,boxsep=0pt,nobeforeafter,
        left=3pt,right=3pt,top=3pt,bottom=3pt,valign=center]{\qrcode[height=1.9cm]{https://play.google.com/store/apps/details?id=com.luvvix.onbuch&hl=fr}}%
      \hspace{8pt}\begin{minipage}[c]{0.5\linewidth}
        {\popdisplay\fontsize{11}{12.5}\selectfont\color{white}\MakeUppercase{Télécharge OnBuch}}\par\vspace{4pt}
        {\raggedright\popsemi\fontsize{8.4}{11}\selectfont\color{white}Gratuit \textbullet\ Google Play\\Tuteur IA, annales, concours, résultats\par}
      \end{minipage}
    \end{tcolorbox}
  \end{minipage}\hfill
  \begin{minipage}[t]{0.487\linewidth}
    \begin{tcolorbox}[enhanced,colback=poppurple,colframe=popink,boxrule=2.2pt,arc=10pt,
        drop shadow=popink,left=11pt,right=11pt,top=10pt,bottom=10pt,height=3.1cm,valign=center]
      \tikz[baseline=-0.7ex]\node[circle,fill=white,draw=popink,line width=1.3pt,minimum size=1.6cm,inner sep=0]{\popdisplay\fontsize{24}{24}\selectfont\color{poppurple}+};%
      \hspace{8pt}\begin{minipage}[c]{0.6\linewidth}
        {\popdisplay\fontsize{11}{12.5}\selectfont\color{white}\MakeUppercase{Passe à OnBuch+}}\par\vspace{4pt}
        {\raggedright\popsemi\fontsize{8.4}{11}\selectfont\color{white}Tous les cours signés, Tuteur IA illimité, fiches PDF — onglet OnBuch+ de l'app\par}
      \end{minipage}
    \end{tcolorbox}
  \end{minipage}
  \vspace{8pt}
  \popcontenu
  \vfill
  \signatureonbuch
  \restoregeometry
  \newpage
}

% Rangée « Ce cours contient » (page de garde)
\newcommand{\poptile}[3]{% #1 couleur #2 gros texte #3 légende
  \begin{tcolorbox}[enhanced,colback=#1,colframe=popink,boxrule=2pt,arc=9pt,shadow={2.5pt}{-2.5pt}{0pt}{popink},
      left=6pt,right=6pt,top=9pt,bottom=9pt,halign=center,valign=center,height=1.75cm,width=\linewidth,nobeforeafter]
    {\popdisplay\fontsize{12.5}{13}\selectfont\color{white}\MakeUppercase{#2}}\par\vspace{3pt}
    {\centering\popsemi\fontsize{7.8}{9.5}\selectfont\color{white}#3\par}
  \end{tcolorbox}}
\newcommand{\popcontenu}{%
  \par\noindent{\popxboldls\fontsize{7.5}{7.5}\selectfont\color{popmuted}CE COURS SIGNÉ CONTIENT}\par\vspace{7pt}
  \noindent\begin{minipage}[t]{0.235\linewidth}\poptile{popblue}{Cours}{complet, illustré, pas à pas}\end{minipage}\hfill
  \begin{minipage}[t]{0.235\linewidth}\poptile{poporange}{Méthodes}{et exercices résolus}\end{minipage}\hfill
  \begin{minipage}[t]{0.235\linewidth}\poptile{poppink}{Exercices}{type examen + corrigés}\end{minipage}\hfill
  \begin{minipage}[t]{0.235\linewidth}\poptile{popgreen}{Fiche bilan}{l'essentiel à retenir}\end{minipage}\par}

% Sommaire
\newtcolorbox{sommaire}{enhanced,colback=white,colframe=popink,boxrule=2.2pt,arc=10pt,
  drop shadow=popink,left=18pt,right=18pt,top=13pt,bottom=13pt,before skip=6pt,after skip=20pt,
  before upper={\poppill{Sommaire}{poppurple}{white}\par\vspace{9pt}\popsemi\fontsize{10.6}{14.5}\selectfont\color{popink}}}

% Signature de fin de cours — poussée tout en bas de la dernière page
\newcommand{\findecours}{%
  \par\vfill\nopagebreak
  \begin{center}\popsquiggle{poporange}\end{center}\vspace{4pt}
  \signatureonbuch
}
\AtBeginDocument{\def\llbracket{\mathopen{[\mkern-3.5mu[}}\def\rrbracket{\mathclose{]\mkern-3.5mu]}}} % intervalles d'entiers (glyphe ⟦ absent de la police)
% Glyphes absents de Fira Math : redéfinis à partir de glyphes disponibles
\AtBeginDocument{%
  \def\setminus{\mathbin{\backslash}}\let\smallsetminus\setminus
  \def\vdots{\mathord{\vbox{\baselineskip3.2pt\lineskiplimit0pt\kern4pt\hbox{.}\hbox{.}\hbox{.}}}}%
  \def\ddots{\mathinner{\mkern1mu\raise6.5pt\hbox{.}\mkern2mu\raise3.6pt\hbox{.}\mkern2mu\raise0.7pt\hbox{.}\mkern1mu}}%
  \def\triangle{\mathord{\scalebox{0.9}{$\symup{\Delta}$}}}%
  \def\nabla{\mathord{\raisebox{\depth}{\scalebox{1}[-1]{$\symup{\Delta}$}}}}%
  \def\checkmark{\popcheck{popdark}}%
  \def\star{\mathbin{\ast}}\def\bigstar{\mathord{\ast}}%
  \def\diamond{\mathbin{\scalebox{0.6}{$\Diamond$}}}%
  \def\bigcup{\mathop{\mathchoice{\raisebox{-0.3em}{\scalebox{1.7}{$\cup$}}}{\scalebox{1.25}{$\cup$}}{\cup}{\cup}}\displaylimits}%
  \def\bigcap{\mathop{\mathchoice{\raisebox{-0.3em}{\scalebox{1.7}{$\cap$}}}{\scalebox{1.25}{$\cap$}}{\cap}{\cap}}\displaylimits}%
  \def\lesseqqgtr{\mathrel{\lessgtr}}\def\bigoplus{\mathop{\oplus}\displaylimits}%
}
\providecommand{\ssi}{si et seulement si} % abréviation fréquente
\AtBeginDocument{\providecommand{\wideparen}{\overparen}} % arc de cercle

%% FIGFIX-BEGIN v5 — correctifs globaux des figures (docs/onbuch-generation/tools/figfix)
\usepackage{adjustbox}\usepackage{etoolbox}\tcbuselibrary{raster}
% toute figure TikZ/pgfplots plus large que la ligne est réduite à la largeur disponible
\BeforeBeginEnvironment{tikzpicture}{\begin{adjustbox}{max width=\linewidth}}
\AfterEndEnvironment{tikzpicture}{\end{adjustbox}}
% légendes pgfplots lisibles par-dessus les courbes
\pgfplotsset{every axis/.append style={legend style={fill=white,fill opacity=0.92,text opacity=1,draw=black!15,inner sep=3pt}}}
% étiquettes d'arêtes (syntaxe edge["..."]) avec fond blanc
\tikzset{every edge quotes/.append style={fill=white,inner sep=1.2pt}}
%% FIGFIX-END
\begin{document}

\pagegarde{Leçon 14 — Grammaire : Le compl. de résultat ; fraction avec 分之、百分之、三分之}{Chinois}{1ère A}{Module 2 : Environnement et santé}{ZHA14}{2 h}{Cette leçon de grammaire présente deux outils essentiels du chinois niveau 1ère : le complément de résultat (V + 完/到/见/懂/错…) qui indique l'aboutissement d'une action, et l'expression des fractions et pourcentages avec 分之、百分之、三分之. Schémas, comparaisons avec le français et exercices de transformation permettent d'éviter les pièges classiques des francophones.
\vspace{4pt}
\begin{multicols}{2}\small
\begin{itemize}
\item À la fin de cette leçon, je sais reconnaître et analyser un complément de résultat dans une phrase
\item je sais construire la structure Sujet + Verbe + complément de résultat (+ complément d'objet)
\item je sais employer les compléments de résultat courants : 完、到、见、懂、对、错、好、清楚
\item je sais mettre le complément de résultat à la forme négative avec 没(有) et à la forme potentielle avec 得/不
\item je sais former et lire une fraction avec 分之 (三分之一, 五分之二)
\item je sais exprimer un pourcentage avec 百分之 (百分之八十)
\end{itemize}\end{multicols}}

\begin{sommaire}
\begin{multicols}{2}
\begin{enumerate}
\item Observation : découvrir le complément de résultat
\item Schéma et construction de la structure
\item Négation, question et forme potentielle
\item Les fractions avec 分之
\item Les pourcentages avec 百分之
\item Comparaison avec le français et synthèse
\item Activité d'intégration
\item Méthodes et exercices résolus
\item Exercices
\item Corrigés des exercices
\item Fiche bilan
\end{enumerate}
\end{multicols}
\end{sommaire}

\begin{objectifs}
\begin{itemize}
\item À la fin de cette leçon, je sais reconnaître et analyser un complément de résultat dans une phrase
\item je sais construire la structure Sujet + Verbe + complément de résultat (+ complément d'objet)
\item je sais employer les compléments de résultat courants : 完、到、见、懂、对、错、好、清楚
\item je sais mettre le complément de résultat à la forme négative avec 没(有) et à la forme potentielle avec 得/不
\item je sais former et lire une fraction avec 分之 (三分之一, 五分之二)
\item je sais exprimer un pourcentage avec 百分之 (百分之八十)
\item je sais comparer l'ordre des mots chinois et français dans les fractions et éviter les erreurs d'inversion
\item je sais utiliser ces structures dans des phrases et un court texte sur l'environnement et la santé
\end{itemize}
\end{objectifs}

\begin{prerequis}
\begin{itemize}
\item Structure de base de la phrase chinoise : Sujet + Verbe + Complément (Seconde)
\item Négation avec 不 et 没(有) (Seconde)
\item Les nombres et les classificateurs courants (Seconde)
\item Le particule 了 d'accomplissement (début de 1ère)
\item Vocabulaire du Module 2 : environnement et santé (病、水、医生、干净…)
\end{itemize}
\end{prerequis}

\newpage

\coursec{Observation : découvrir le complément de résultat}

Au marché Mokolo de Yaoundé, un élève de 1ère A a fini de vendre ses tomates : il annonce fièrement à son père \zh{我卖完了！} (\emph{Wǒ màiwán le !} — « J'ai tout vendu ! »). Plus tard, au lycée, son professeur de SVT explique que \zh{百分之七十} (\emph{bǎifēn zhī qīshí} — 70 \%) des maladies viennent de l'eau non potable, et que \zh{三分之一} (\emph{sān fēn zhī yī} — un tiers) des élèves de la classe sont déjà tombés malades cette année. Comment dire en chinois qu'une action est « terminée, réussie, bien faite » ? Et comment exprimer les fractions et les pourcentages que l'on rencontre partout : notes, statistiques de santé, sondages ? C'est toute la richesse de cette leçon de grammaire.

Dans cette première section, nous allons \textbf{observer} des phrases authentiques, remarquer ce qu'elles ont en commun, puis dégager la règle du \cle{complément de résultat}.

\courssub{Lecture guidée : six phrases à observer}

Lisons attentivement les six phrases suivantes. Dans chacune, le \textbf{verbe principal} et le mot qui le suit immédiatement sont soulignés.

\begin{enumerate}
\item \zh{我做完了作业。} \emph{Wǒ zuòwán le zuòyè.} — J'ai fini mes devoirs. (\underline{做} + \underline{完})
\item \zh{我看见他了。} \emph{Wǒ kànjiàn tā le.} — Je l'ai vu / aperçu. (\underline{看} + \underline{见})
\item \zh{我听懂了老师的话。} \emph{Wǒ tīngdǒng le lǎoshī de huà.} — J'ai compris ce que le professeur a dit. (\underline{听} + \underline{懂})
\item \zh{你找没找到你的书？} \emph{Nǐ zhǎo méi zhǎodào nǐ de shū ?} — As-tu retrouvé ton livre ? (\underline{找} + \underline{到})
\item \zh{他说错了。} \emph{Tā shuōcuò le.} — Il s'est trompé en parlant. (\underline{说} + \underline{错})
\item \zh{饭做好了。} \emph{Fàn zuòhǎo le.} — Le repas est prêt. (\underline{做} + \underline{好})
\end{enumerate}

\begin{methode}[Question-guide : que remarque-t-on ?]
\begin{enumerate}
\item Chaque phrase contient un verbe d'action : 做 (faire), 看 (regarder), 听 (écouter), 找 (chercher), 说 (dire).
\item Juste après ce verbe, \textbf{sans aucun mot entre les deux}, on trouve un second mot : 完, 见, 懂, 到, 错, 好.
\item Ce second mot n'exprime pas une nouvelle action : il exprime le \cle{résultat} de l'action — l'action \emph{aboutit} : finie (完), aperçue (见), comprise (懂), atteinte (到), ratée (错), bien terminée (好).
\item La particule \zh{了} se place \textbf{après} le bloc verbe + résultat, jamais entre les deux.
\end{enumerate}
Conclusion : le chinois colle au verbe un mot qui dit si l'action a \emph{réussi}, \emph{abouti} ou \emph{comment elle s'est terminée}. C'est le complément de résultat.
\end{methode}

\courssub{Un texte pour observer la structure en contexte}

\begin{textebox}[Après le cours de chinois]
昨天的作业我做完了。老师的话我听懂了，可是有两个字我写错了。\\
\emph{Zuótiān de zuòyè wǒ zuòwán le. Lǎoshī de huà wǒ tīngdǒng le, kěshì yǒu liǎng ge zì wǒ xiěcuò le.}\\
Les devoirs d'hier, je les ai finis. Ce que le professeur a dit, je l'ai compris, mais il y a deux caractères que j'ai mal écrits.
\end{textebox}

Repérage : dans ce court texte, on trouve trois compléments de résultat — \zh{做完} (finir de faire), \zh{听懂} (comprendre en écoutant), \zh{写错} (mal écrire). Remarque importante : le locuteur place l'objet (\zh{昨天的作业}, \zh{老师的话}) \emph{avant} le verbe, une tournure très fréquente en chinois pour mettre l'objet en relief. La structure Verbe + Résultat, elle, ne change pas.

\coursec{Schéma et construction de la structure}

\courssub{Définition et schéma}

\begin{definition}[Le complément de résultat (结果补语)]
Le \cle{complément de résultat} (en chinois \zh{结果补语}, \emph{jiéguǒ bǔyǔ}) est un \textbf{verbe ou un adjectif placé juste après le verbe principal}, sans aucun mot entre les deux, pour indiquer que l'action \textbf{aboutit à un résultat} : action terminée, réussie, obtenue, correcte ou ratée. Le verbe principal exprime \emph{l'action} ; le complément exprime \emph{l'issue de l'action}.
\end{definition}

\begin{propriete}[Schéma de construction]
Structure de base : \textbf{Sujet + Verbe + Résultat + (了) + Objet}
\begin{itemize}
\item Le verbe et son résultat forment un bloc inséparable : on ne peut rien insérer entre eux.
\item \zh{了} se place après le bloc Verbe + Résultat (et non après le verbe seul) : \zh{我做完了作业} et non *\zh{我做了完作业}.
\item L'objet vient en dernier : \zh{我听懂了老师的话} (j'ai compris + les paroles du professeur).
\end{itemize}
\end{propriete}

\begin{popfigure}
\begin{tikzpicture}[node distance=0.55cm, every node/.style={font=\small}]
\node[draw=popblue, fill=popblueL, rounded corners, minimum width=1.9cm, minimum height=0.85cm, align=center] (s) {Sujet\\ \zh{我}};
\node[draw=poporange, fill=poporangeL, rounded corners, minimum width=2.6cm, minimum height=0.85cm, align=center, right=of s] (v) {Verbe + Résultat\\ \zh{做完}};
\node[draw=poppurple, fill=poppurpleL, rounded corners, minimum width=1.1cm, minimum height=0.85cm, align=center, right=of v] (le) {(了)\\ \zh{了}};
\node[draw=popgreen, fill=popgreenL, rounded corners, minimum width=1.9cm, minimum height=0.85cm, align=center, right=of le] (o) {Objet\\ \zh{作业}};
\draw[-{Stealth}, thick, popblue] (s) -- (v);
\draw[-{Stealth}, thick, popblue] (v) -- (le);
\draw[-{Stealth}, thick, popblue] (le) -- (o);
\node[draw=poppink, fill=poppinkL, rounded corners, align=center, below=0.7cm of v, font=\footnotesize] (res) {résultat de l'action :\\ terminée, réussie, correcte ou ratée};
\draw[-{Stealth}, thick, poppink] (res) -- (v);
\end{tikzpicture}
\legende{Schéma de la structure : Sujet → Verbe + Résultat → (了) → Objet}
\end{popfigure}

\courssub{Inventaire des compléments de résultat courants}

Voici les huit compléments de résultat les plus fréquents au niveau HSK 2-3. Chacun se combine avec de nombreux verbes.

\begin{center}
\begin{tabularx}{\linewidth}{|X|X|X|X|}
\hline
\rowcolor{popblueL} Caractères & Pinyin & Nature & Traduction \\
\hline
完 & wán & verbe/compl. & finir, terminer \\
\hline
到 & dào & verbe/compl. & atteindre, obtenir, réussir à \\
\hline
见 & jiàn & verbe/compl. & apercevoir, percevoir \\
\hline
懂 & dǒng & verbe/compl. & comprendre \\
\hline
对 & duì & adjectif/compl. & juste, correct \\
\hline
错 & cuò & adjectif/compl. & faux, erroné \\
\hline
好 & hǎo & adjectif/compl. & bien, terminé, prêt \\
\hline
清楚 & qīngchu & adjectif/compl. & clair, clairement \\
\hline
\end{tabularx}
\end{center}

\begin{exemplebox}[Combinaisons fréquentes Verbe + Résultat]
\begin{itemize}
\item \zh{看完} \emph{kànwán} : finir de lire/regarder → \zh{我看完了这本书。} \emph{Wǒ kànwán le zhè běn shū.} — J'ai fini de lire ce livre.
\item \zh{买到} \emph{mǎidào} : réussir à acheter → \zh{她买到票了。} \emph{Tā mǎidào piào le.} — Elle a réussi à acheter le billet.
\item \zh{听见} \emph{tīngjiàn} : entendre → \zh{你听见了吗？} \emph{Nǐ tīngjiàn le ma ?} — Tu as entendu ?
\item \zh{看清楚} \emph{kàn qīngchu} : voir clairement → \zh{我没看清楚黑板。} \emph{Wǒ méi kàn qīngchu hēibǎn.} — Je n'ai pas vu clairement le tableau.
\item \zh{答对} \emph{dáduì} : répondre juste → \zh{他答对了。} \emph{Tā dáduì le.} — Il a répondu juste.
\item \zh{做好} \emph{zuòhǎo} : bien faire, achever → \zh{饭做好了。} \emph{Fàn zuòhǎo le.} — Le repas est prêt.
\end{itemize}
\end{exemplebox}

\begin{popfigure}
\begin{tikzpicture}[mindmap, grow cyclic, every node/.style={concept, font=\small},
root concept/.append style={concept color=popblue, font=\small\bfseries},
level 1/.append style={level distance=2.9cm, sibling angle=45},
level 1 concept/.append style={concept color=poporange!70}]
\node[root concept]{V + résultat}
child { node[align=center]{\zh{完} wán\\ finir} }
child { node[align=center]{\zh{到} dào\\ obtenir} }
child { node[align=center]{\zh{见} jiàn\\ apercevoir} }
child { node[align=center]{\zh{懂} dǒng\\ comprendre} }
child { node[align=center]{\zh{对} duì\\ juste} }
child { node[align=center]{\zh{错} cuò\\ faux} }
child { node[align=center]{\zh{好} hǎo\\ bien, prêt} }
child { node[align=center]{\zh{清楚} qīngchu\\ clairement} };
\end{tikzpicture}
\legende{Carte mentale : les huit compléments de résultat courants}
\end{popfigure}

\coursec{Négation, question et forme potentielle}

\courssub{La négation avec 没(有)}

\begin{propriete}[Nier un résultat]
Pour dire qu'une action \textbf{n'a pas abouti}, on place \zh{没} (ou \zh{没有}) \textbf{devant} le bloc Verbe + Résultat, et \zh{了} disparaît :
\textbf{Sujet + 没(有) + Verbe + Résultat + Objet}
\begin{itemize}
\item \zh{我没做完作业。} \emph{Wǒ méi zuòwán zuòyè.} — Je n'ai pas fini mes devoirs.
\item \zh{她没买到票。} \emph{Tā méi mǎidào piào.} — Elle n'a pas réussi à acheter le billet.
\item \zh{我没听清楚。} \emph{Wǒ méi tīng qīngchu.} — Je n'ai pas bien entendu.
\end{itemize}
Exception à retenir : on n'utilise \textbf{jamais} \zh{不} pour nier un résultat accompli ; \zh{不} + complément de résultat est réservé à la forme potentielle (voir ci-dessous).
\end{propriete}

\courssub{La question}

Deux formes de question sont possibles :
\begin{itemize}
\item \textbf{Question avec 吗} : \zh{你看见王老师了吗？} \emph{Nǐ kànjiàn Wáng lǎoshī le ma ?} — As-tu vu le professeur Wang ?
\item \textbf{Question affirmative-négative} : on répète le bloc en opposant la forme pleine et la forme niée : \zh{你找没找到你的书？} \emph{Nǐ zhǎo méi zhǎodào nǐ de shū ?} — As-tu retrouvé ton livre ? (schéma : V + 没 + V-Résultat).
\end{itemize}

\courssub{La forme potentielle : 得 / 不 entre verbe et résultat}

\begin{definition}[Le complément de résultat potentiel]
En insérant \zh{得} (possibilité) ou \zh{不} (impossibilité) \textbf{entre} le verbe et le résultat, on exprime que l'action \textbf{peut ou ne peut pas aboutir} :
\textbf{V + 得 + Résultat} = capable d'aboutir ; \textbf{V + 不 + Résultat} = incapable d'aboutir.
\end{definition}

\begin{exemplebox}[Forme potentielle : exemples]
\begin{itemize}
\item \zh{我听得懂。} \emph{Wǒ tīng de dǒng.} — Je peux comprendre (en écoutant).
\item \zh{我听不懂。} \emph{Wǒ tīng bu dǒng.} — Je ne comprends pas (en écoutant).
\item \zh{这个字你看得见吗？} \emph{Zhège zì nǐ kàn de jiàn ma ?} — Ce caractère, arrives-tu à le voir ?
\item \zh{黑板上的字我看不见。} \emph{Hēibǎn shang de zì wǒ kàn bu jiàn.} — Je n'arrive pas à voir les caractères au tableau.
\item \zh{今天的作业我做得完。} \emph{Jīntiān de zuòyè wǒ zuò de wán.} — Les devoirs d'aujourd'hui, je peux les finir.
\end{itemize}
\end{exemplebox}

\begin{attention}[得 : deux emplois à ne pas confondre]
Dans \zh{听得懂}, \zh{得} est infixé entre verbe et résultat et marque la \textbf{possibilité} ; dans \zh{他说得很好} (\emph{Tā shuō de hěn hǎo.} — Il parle très bien), \zh{得} introduit un \textbf{complément de degré} après le verbe. Même caractère, deux structures différentes : ne les confondez pas dans vos traductions.
\end{attention}

\begin{exoresolu}[Transformer à la forme potentielle]
Énoncé : mets les phrases suivantes à la forme potentielle négative (impossibilité) et traduis : (1) \zh{我听懂了。} (2) \zh{他买到票了。}
\tcblower
Solution détaillée :
\begin{enumerate}
\item \zh{我听不懂。} \emph{Wǒ tīng bu dǒng.} — « Je ne comprends pas (ce que j'entends). » On remplace \zh{了} par l'infixe \zh{不} entre 听 et 懂.
\item \zh{他买不到票。} \emph{Tā mǎi bu dào piào.} — « Il n'arrive pas à acheter de billet. » Même procédé : V + 不 + Résultat, et \zh{了} disparaît.
\end{enumerate}
Méthode : à la forme potentielle, \zh{了} n'apparaît jamais ; seul l'infixe \zh{得} ou \zh{不} s'insère dans le bloc verbal.
\end{exoresolu}

\coursec{Les fractions avec 分之}

\courssub{Observer et construire}

Au marché, le père de notre élève partage la recette : \zh{三分之一} pour lui, \zh{三分之二} pour ses enfants. Comment lit-on une fraction en chinois ?

\begin{propriete}[Schéma de la fraction]
Structure : \textbf{Dénominateur + 分之 + Numérateur} (littéralement « sur X parts, Y »).
\begin{itemize}
\item \zh{三分之一} \emph{sān fēn zhī yī} — un tiers (1/3).
\item \zh{三分之二} \emph{sān fēn zhī èr} — deux tiers (2/3).
\item \zh{四分之一} \emph{sì fēn zhī yī} — un quart (1/4).
\item \zh{五分之三} \emph{wǔ fēn zhī sān} — trois cinquièmes (3/5).
\item \zh{十分之七} \emph{shí fēn zhī qī} — sept dixièmes (7/10).
\end{itemize}
\end{propriete}

\begin{attention}[L'ordre inversé du français]
Le français dit « un \textbf{tiers} » (numérateur d'abord) ; le chinois dit \zh{三分之一}, littéralement « de trois parts, une » : le \textbf{dénominateur vient d'abord}. Erreur fréquente : *\zh{一分之三} pour « un tiers » — c'est fautif. Retenez : \zh{分之} = « sur ... parts », le nombre \emph{avant} \zh{分之} est le dénominateur.
\end{attention}

\begin{exemplebox}[Fractions en contexte]
\begin{itemize}
\item \zh{三分之一的学生喜欢足球。} \emph{Sān fēn zhī yī de xuésheng xǐhuan zúqiú.} — Un tiers des élèves aiment le football.
\item \zh{他把四分之三的钱给了妈妈。} \emph{Tā bǎ sì fēn zhī sān de qián gěi le māma.} — Il a donné les trois quarts de l'argent à sa mère.
\item \zh{我们班五分之二的同学住在雅温得。} \emph{Wǒmen bān wǔ fēn zhī èr de tóngxué zhù zài Yǎwēndé.} — Deux cinquièmes des élèves de notre classe habitent à Yaoundé.
\end{itemize}
Note : la fraction se comporte comme un déterminant ; elle est suivie de \zh{的} devant le nom : \zh{三分之一的学生}.
\end{exemplebox}

\coursec{Les pourcentages avec 百分之}

\courssub{Schéma et emplois}

\begin{propriete}[Schéma du pourcentage]
Structure : \textbf{百分之 + nombre} (littéralement « sur cent parts, N ») = N \%.
\begin{itemize}
\item \zh{百分之一} \emph{bǎi fēn zhī yī} — 1 \%.
\item \zh{百分之二十} \emph{bǎi fēn zhī èrshí} — 20 \%.
\item \zh{百分之五十} \emph{bǎi fēn zhī wǔshí} — 50 \%.
\item \zh{百分之七十五} \emph{bǎi fēn zhī qīshíwǔ} — 75 \%.
\item \zh{百分之百} \emph{bǎi fēn zhī bǎi} — 100 \%.
\end{itemize}
Le pourcentage n'est qu'un cas particulier de la fraction : le dénominateur est toujours \zh{百} (cent).
\end{propriete}

\begin{exemplebox}[Pourcentages en contexte : santé et environnement]
\begin{itemize}
\item \zh{百分之七十的疾病来自不干净的水。} \emph{Bǎi fēn zhī qīshí de jíbìng láizì bù gānjìng de shuǐ.} — 70 \% des maladies viennent de l'eau insalubre.
\item \zh{我们班百分之九十的学生通过了考试。} \emph{Wǒmen bān bǎi fēn zhī jiǔshí de xuésheng tōngguò le kǎoshì.} — 90 \% des élèves de notre classe ont réussi l'examen.
\item \zh{今天百分之五十的商店开门了。} \emph{Jīntiān bǎi fēn zhī wǔshí de shāngdiàn kāimén le.} — Aujourd'hui, 50 \% des magasins ont ouvert.
\end{itemize}
\end{exemplebox}

\begin{attention}[Pièges avec les pourcentages]
\begin{itemize}
\item \textbf{Ordre} : on dit \zh{百分之二十} et jamais *\zh{二十百分之} ; le symbole \% se lit \emph{après} le nombre en français, mais la structure chinoise place \zh{百分之} \emph{avant}.
\item \textbf{Pas de classificateur} entre \zh{百分之} et le nombre : *\zh{百分之个二十} est fautif.
\item \textbf{Devant un nom}, ajoutez \zh{的} : \zh{百分之七十的疾病} (70 \% des maladies).
\item Ne confondez pas \zh{分} (fēn, part) avec \zh{分钟} (fēnzhōng, minute) : homophones partiels, sens très différents.
\end{itemize}
\end{attention}

\begin{exoresolu}[Traduire des fractions et des pourcentages]
Énoncé : traduis en chinois (caractères + pinyin) : (1) « deux cinquièmes » ; (2) « 30 \% des élèves sont absents » ; (3) « un quart de l'eau est sale ».
\tcblower
Solution détaillée :
\begin{enumerate}
\item \zh{五分之二} \emph{wǔ fēn zhī èr} — dénominateur 5 d'abord, puis \zh{分之}, puis numérateur 2.
\item \zh{百分之三十的学生没来。} \emph{Bǎi fēn zhī sānshí de xuésheng méi lái.} — « 30 \% des élèves ne sont pas venus. » Pourcentage + \zh{的} + nom ; négation du résultat avec \zh{没}.
\item \zh{四分之一的水很脏。} \emph{Sì fēn zhī yī de shuǐ hěn zāng.} — « Un quart de l'eau est sale. » Fraction + \zh{的} + nom.
\end{enumerate}
Méthode : (1) identifier dénominateur et numérateur ; (2) écrire dénominateur + \zh{分之} + numérateur ; (3) ajouter \zh{的} devant le nom concerné.
\end{exoresolu}

\coursec{Comparaison avec le français et synthèse}

\courssub{Tableau comparatif des structures}

\begin{center}
\begin{tabularx}{\linewidth}{|X|X|X|}
\hline
\rowcolor{popblueL} Structure & Exemple en caractères + pinyin & Traduction \\
\hline
V + 完 (action finie) & \zh{我做完了作业。} Wǒ zuòwán le zuòyè. & J'ai fini mes devoirs. \\
\hline
V + 见 (perception aboutie) & \zh{我看见他了。} Wǒ kànjiàn tā le. & Je l'ai aperçu. \\
\hline
V + 懂 (compréhension obtenue) & \zh{我听懂了。} Wǒ tīngdǒng le. & J'ai compris (en écoutant). \\
\hline
V + 错 (action ratée) & \zh{他说错了。} Tā shuōcuò le. & Il s'est trompé en parlant. \\
\hline
Négation : 没 + V + R & \zh{我没做完作业。} Wǒ méi zuòwán zuòyè. & Je n'ai pas fini mes devoirs. \\
\hline
Potentiel : V + 不 + R & \zh{我听不懂。} Wǒ tīng bu dǒng. & Je ne comprends pas. \\
\hline
Fraction : dénom. + 分之 + num. & \zh{三分之二} sān fēn zhī èr & deux tiers \\
\hline
Pourcentage : 百分之 + N & \zh{百分之七十} bǎi fēn zhī qīshí & 70 \% \\
\hline
\end{tabularx}
\end{center}

En français, le résultat est exprimé par des mots séparés : un adverbe (« mal écrire »), une préposition (« finir \emph{de} lire »), un auxiliaire (« réussir \emph{à} acheter ») ou un verbe différent (« regarder » / « voir »). En chinois, tout se dit en un seul bloc compact : verbe + résultat collés. C'est une différence majeure : le francophone doit apprendre à \emph{fusionner} les deux idées en une seule unité verbale. Pour les fractions, c'est l'inverse : le français met le numérateur d'abord (« deux tiers »), le chinois met le dénominateur d'abord (\zh{三分之二}).

\begin{attention}[Bilan des pièges fréquents des francophones]
\begin{itemize}
\item \textbf{Ne séparez jamais} le verbe de son résultat : *\zh{我做完的作业} est fautif ; on dit \zh{我做完了作业} (了 après le bloc, objet à la fin).
\item \textbf{了 ne suit pas le verbe seul} : *\zh{我听了懂老师的话} est impossible. Le bloc \zh{听懂} est inséparable.
\item \textbf{Ne traduisez pas mot à mot} « réussir à » : \zh{买到} suffit, inutile d'ajouter un verbe comme « réussir ».
\item \textbf{Attention à 看 seul} : \zh{看} signifie « regarder » (l'action), pas « voir » (le résultat). Pour « je l'ai vu », il faut \zh{我看见他了}.
\item \textbf{Négation} : toujours \zh{没(有)}, jamais \zh{不}, devant un résultat accompli ; \zh{不} sert d'infixe dans la forme potentielle (\zh{听不懂}).
\item \textbf{Fractions} : dénominateur d'abord — \zh{三分之一} = 1/3, et non l'inverse.
\end{itemize}
\end{attention}

\begin{savaistu}[Pourquoi « voir » se dit 看见 ?]
En chinois, « voir » se dit \zh{看见}, littéralement « regarder-apercevoir » : \zh{看} est l'action de diriger les yeux, \zh{见} est le résultat — l'objet entre effectivement dans le champ visuel. On ne dit presque jamais \zh{看} seul pour « voir » au sens de résultat. De même, « entendre » se dit \zh{听见} (« écouter-percevoir »). Le chinois distingue donc nettement l'\emph{effort} (regarder, écouter) et le \emph{résultat} (voir, entendre), là où le français utilise deux verbes différents sans lien visible entre eux. Autre fait culturel : en Chine comme au Cameroun, les statistiques de santé publique (eau potable, vaccination) s'expriment en pourcentages — savoir lire \zh{百分之} est indispensable pour comprendre les informations sanitaires dans les deux pays.
\end{savaistu}

\begin{experience}[Jeu de rôle : au dispensaire]
Par groupes de deux : l'un joue un médecin du quartier, l'autre un élève journaliste. Le médecin annonce des statistiques avec \zh{百分之} et \zh{分之} (par exemple : \zh{百分之六十的病人喝水不干净的水。}), le journaliste reformule et pose des questions avec le complément de résultat (\zh{你听清楚了吗？}, \zh{我听懂了吗？}). Chaque binôme présente ensuite deux phrases à la classe.
\end{experience}

\begin{exoresolu}[Identifier le complément de résultat]
Énoncé : dans les phrases suivantes, souligne le verbe principal et entoure le complément de résultat, puis traduis : (1) \zh{我吃饱了。} (2) \zh{她找到了工作。} (3) \zh{你听清楚了吗？}
\tcblower
Solution détaillée :
\begin{enumerate}
\item \zh{我吃饱了。} \emph{Wǒ chībǎo le.} — verbe : \underline{吃} (manger) ; complément : 饱 (rassasié). Traduction : « J'ai mangé à ma faim / je suis rassasié. »
\item \zh{她找到了工作。} \emph{Tā zhǎodào le gōngzuò.} — verbe : \underline{找} (chercher) ; complément : 到 (atteindre, obtenir). Traduction : « Elle a trouvé un travail. »
\item \zh{你听清楚了吗？} \emph{Nǐ tīng qīngchu le ma ?} — verbe : \underline{听} (écouter) ; complément : 清楚 (clairement). Traduction : « As-tu bien entendu ? »
\end{enumerate}
Méthode : le verbe principal est toujours le \emph{premier} mot du bloc ; le complément de résultat est le mot \emph{immédiatement collé} après lui, avant \zh{了} et avant l'objet.
\end{exoresolu}

\begin{exercice}[À toi de jouer]
\begin{enumerate}
\item Traduis en français, puis repère le verbe et le complément de résultat : (a) \zh{作业写完了。} (b) \zh{这个字我读错了。} (c) \zh{你看见王老师了吗？}
\item Mets à la forme négative : (a) \zh{我做完了作业。} (b) \zh{他买到票了。}
\item Traduis en chinois : (a) « un quart » ; (b) « 80 \% des élèves ont compris » ; (c) « Je n'ai pas fini de manger. »
\end{enumerate}
\end{exercice}

\begin{corrige}[Exercice « À toi de jouer »]
\textbf{1.} (a) \zh{作业写完了。} \emph{Zuòyè xiěwán le.} — « Les devoirs sont finis. » Verbe : 写 ; résultat : 完. (b) \zh{这个字我读错了。} \emph{Zhège zì wǒ dúcuò le.} — « Ce caractère, je l'ai mal lu. » Verbe : 读 ; résultat : 错. (c) \zh{你看见王老师了吗？} \emph{Nǐ kànjiàn Wáng lǎoshī le ma ?} — « As-tu vu le professeur Wang ? » Verbe : 看 ; résultat : 见.\\
\textbf{2.} (a) \zh{我没做完作业。} \emph{Wǒ méi zuòwán zuòyè.} (b) \zh{他没买到票。} \emph{Tā méi mǎidào piào.} — \zh{没} devant le bloc, \zh{了} disparaît.\\
\textbf{3.} (a) \zh{四分之一} \emph{sì fēn zhī yī}. (b) \zh{百分之八十的学生听懂了。} \emph{Bǎi fēn zhī bāshí de xuésheng tīngdǒng le.} (c) \zh{我没吃完饭。} \emph{Wǒ méi chīwán fàn.}
\end{corrige}

\begin{aretenir}[L'essentiel de la leçon]
\begin{itemize}
\item Le complément de résultat est un verbe ou un adjectif \textbf{collé} après le verbe principal : il exprime l'\textbf{issue} de l'action.
\item Schéma : Sujet + Verbe + Résultat + (了) + Objet. Rien ne s'insère entre le verbe et son résultat ; \zh{了} vient après le bloc.
\item Négation : \zh{没(有)} + V + Résultat, sans \zh{了}. Potentiel : V + \zh{得}/\zh{不} + Résultat (\zh{听得懂} / \zh{听不懂}).
\item Les huit compléments à connaître : 完, 到, 见, 懂, 对, 错, 好, 清楚.
\item Fraction : \textbf{dénominateur + 分之 + numérateur} (\zh{三分之一} = 1/3). Pourcentage : \textbf{百分之 + nombre} (\zh{百分之七十} = 70 \%). Devant un nom, ajouter \zh{的}.
\item Le français dit « finir de », « réussir à », « mal » avec des mots séparés ; le chinois fusionne action et résultat en un seul bloc verbal.
\end{itemize}
\end{aretenir}

\coursec{Schéma et construction de la structure}

Dans la section précédente, nous avons \emph{observé} des phrases comme \zh{我听懂了} (\emph{Wǒ tīng dǒng le}, « j'ai entendu et compris ») et découvert que le chinois colle au verbe un second élément qui en donne le \cle{résultat}. Nous allons maintenant passer de l'observation à la \textbf{construction} : comment fabrique-t-on, à coup sûr, une phrase correcte avec un complément de résultat ? C'est l'objet de cette section : un schéma général, des règles de placement strictes, puis un répertoire des paires « verbe + résultat » les plus utiles.

\courssub{Le schéma général de la phrase}

\begin{propriete}[Schéma général du complément de résultat]
La phrase affirmative au complément de résultat suit l'ordre rigide suivant :
\begin{center}
\textbf{Sujet + Verbe + Résultat (+ 了) (+ Objet)}
\end{center}
\begin{itemize}
\item Le \cle{verbe} exprime l'action : \zh{看} kàn « regarder », \zh{听} tīng « écouter », \zh{写} xiě « écrire », \zh{买} mǎi « acheter ».
\item Le \cle{résultat} (un caractère, parfois deux) exprime l'aboutissement : \zh{见} jiàn « percevoir », \zh{懂} dǒng « comprendre », \zh{完} wán « finir », \zh{到} dào « atteindre », \zh{对} duì « juste », \zh{错} cuò « faux ».
\item \zh{了} le marque l'accomplissement et se place \textbf{après tout le groupe} Verbe + Résultat.
\item L'\cle{objet} se place \textbf{en dernier}, après le groupe Verbe + Résultat (+ 了).
\end{itemize}
\end{propriete}

\begin{exemplebox}[Le schéma en action]
\zh{我做完了作业。}\\
\emph{Wǒ zuò wán le zuòyè.}\\
« J'ai fini mes devoirs. »\\[4pt]
Découpage : \zh{我} (S) + \zh{做完} (V + R, bloc unique) + \zh{了} + \zh{作业} (O).\\[6pt]
\zh{我找到了我的钥匙。}\\
\emph{Wǒ zhǎodào le wǒ de yàoshi.}\\
« J'ai retrouvé mes clés. »\\[4pt]
Découpage : \zh{我} (S) + \zh{找到} (V + R) + \zh{了} + \zh{我的钥匙} (O).
\end{exemplebox}

\begin{popfigure}
\begin{center}
\begin{tikzpicture}[
  bloc/.style={draw=popblue, fill=popblueL, rounded corners=2pt, minimum height=0.9cm, align=center, font=\small},
  blocR/.style={draw=poporange, fill=poporangeL, rounded corners=2pt, minimum height=0.9cm, align=center, font=\small},
  blocL/.style={draw=popgreen, fill=popgreenL, rounded corners=2pt, minimum height=0.9cm, align=center, font=\small},
  blocO/.style={draw=poppurple, fill=poppurpleL, rounded corners=2pt, minimum height=0.9cm, align=center, font=\small},
  lbl/.style={font=\scriptsize\itshape, text=popmuted}
]
\node[bloc, align=center] (S) at (0,0){我\\ Sujet};
\node[blocR, right=0.6cm of S, align=center] (VR){做完\\ V + R (bloc unique)};
\node[blocL, right=0.6cm of VR, align=center] (LE){了\\ accompli};
\node[blocO, right=0.6cm of LE, align=center] (O){作业\\ Objet};
\draw[-{Stealth}, thick, popdark] (S) -- (VR);
\draw[-{Stealth}, thick, popdark] (VR) -- (LE);
\draw[-{Stealth}, thick, popdark] (LE) -- (O);
\node[lbl, below=0.25cm of VR] (interdit) {position interdite};
\draw[thick, poppink] ($(VR.south west)+(0.15,-0.05)$) -- ($(VR.south west)+(-0.05,-0.25)$);
\draw[thick, poppink] ($(VR.south west)+(-0.05,-0.05)$) -- ($(VR.south west)+(0.15,-0.25)$);
\node[font=\scriptsize, text=poppink, above=0.15cm of VR] {rien ne s'insère entre 做 et 完};
\end{tikzpicture}
\end{center}
\legende{Schéma de la phrase : Sujet + [Verbe + Résultat] + 了 + Objet. La croix marque la position interdite entre le verbe et son résultat.}
\end{popfigure}

\courssub{Verbe et résultat : un bloc inséparable}

\begin{propriete}[Règle d'ordre des mots]
Le complément de résultat se \cle{colle immédiatement au verbe} : les deux forment un bloc inséparable, comme un seul mot. \textbf{Aucun mot} (ni 了, ni l'objet, ni un complément de temps ou de lieu) ne peut s'insérer entre le verbe et son résultat. L'objet vient \textbf{après tout le groupe}.
\end{propriete}

Comparons avec le français. En français, on dit « j'ai \emph{fini} mes devoirs » : un seul verbe, et l'objet juste après. Le francophone est donc tenté de reproduire cet ordre en chinois. Or le chinois exige \textbf{deux éléments soudés} : le verbe d'action \zh{做} « faire » + le résultat \zh{完} « fini ». C'est ce groupe soudé \zh{做完} qui joue le rôle du verbe français « finir ».

\begin{center}
\begin{tabularx}{\linewidth}{|X|X|X|}
\hline
\rowcolor{popblueL} \textbf{Phrase} & \textbf{Analyse} & \textbf{Verdict} \\
\hline
\zh{我做完了作业。} \emph{Wǒ zuò wán le zuòyè.} & 了 après le groupe 做完 & Correct \\
\hline
\zh{我做了完作业。} \emph{Wǒ zuò le wán zuòyè.} & 了 inséré entre V et R & Faux \\
\hline
\zh{我听懂了老师的话。} \emph{Wǒ tīng dǒng le lǎoshī de huà.} & objet après le groupe & Correct \\
\hline
\zh{我听了懂老师的话。} \emph{Wǒ tīng le dǒng lǎoshī de huà.} & 了 entre V et R & Faux \\
\hline
\zh{我完了作业。} \emph{Wǒ wán le zuòyè.} & verbe d'action 做 absent & Faux \\
\hline
\end{tabularx}
\end{center}

\begin{attention}[Erreur fréquente des francophones : traduire mot à mot]
Ne traduisez jamais « j'ai fini mes devoirs » par \zh{我完了作业}. En chinois, \zh{完} « finir » ne peut pas porter seul un objet : il faut d'abord le \textbf{verbe d'action} qui dit \emph{ce qu'on fait}, puis \zh{完} qui dit \emph{que c'est fini}. On dit donc \zh{我做完了作业} (\emph{Wǒ zuò wán le zuòyè}) : littéralement « je faire-finir-ACCOMPLI devoirs ». Même logique partout : « j'ai fini de manger » = \zh{我吃完了} (manger-finir), « j'ai fini d'écrire la lettre » = \zh{我写完了信} (écrire-finir-lettre).
\end{attention}

\begin{exemplebox}[Exemples traduits]
\zh{请你说清楚。}\\
\emph{Qǐng nǐ shuō qīngchu.}\\
« Explique-toi clairement, s'il te plaît. » (pas de 了 : c'est un impératif, l'action n'est pas encore accomplie)\\[6pt]
\zh{妈妈买到了新鲜的鱼。}\\
\emph{Māma mǎi dào le xīnxiān de yú.}\\
« Maman a réussi à acheter du poisson frais. »\\[6pt]
\zh{我看错了这个字。}\\
\emph{Wǒ kàn cuò le zhège zì.}\\
« J'ai mal lu ce caractère. »
\end{exemplebox}

Remarque utile : quand l'objet est long ou quand on insiste sur l'accomplissement global, on entend aussi un second \zh{了} en toute fin de phrase : \zh{我做完了作业了} (\emph{Wǒ zuò wán le zuòyè le}, « j'ai fini mes devoirs, c'est fait »). À ce stade, contentez-vous de la structure de base : un seul \zh{了}, placé après le groupe Verbe + Résultat.

\courssub{La construction avec 把 (rappel utile)}

Le complément de résultat s'emploie très souvent avec la structure en \zh{把} bǎ, étudiée précédemment, car \zh{把} exige que le verbe soit suivi d'un élément qui montre ce que l'objet \emph{devient}. Le groupe Verbe + Résultat remplit parfaitement ce rôle. Schéma : \textbf{Sujet + 把 + Objet + Verbe + Résultat (+ 了)}.

\begin{exemplebox}[Avec 把]
\zh{我把作业做完了。}\\
\emph{Wǒ bǎ zuòyè zuò wán le.}\\
« J'ai fini mes devoirs. » (l'objet 作业 est avancé avant le verbe)\\[6pt]
\zh{请把话说清楚。}\\
\emph{Qǐng bǎ huà shuō qīngchu.}\\
« Exprime-toi clairement, s'il te plaît. »
\end{exemplebox}

Signalons simplement ici que la négation (\zh{没(有)}), la question et la forme potentielle (\zh{看得见} / \zh{看不见}) feront l'objet de la section suivante ; retenez pour l'instant que le bloc Verbe + Résultat reste inséparable dans toutes ces constructions.

\courssub{Les paires verbe + résultat lexicalisées}

Certaines combinaisons sont si fréquentes qu'elles fonctionnent comme de véritables verbes composés : il faut les apprendre \textbf{en paires}. Le tableau croisé ci-dessous montre comment un même verbe d'action se combine avec plusieurs résultats, et inversement. Attention : toutes les combinaisons ne sont pas possibles ; seules les cases remplies correspondent à des paires réellement employées.

\begin{popfigure}
\begin{center}
\begin{tikzpicture}[font=\small]
\draw[fill=popblueL, draw=popblue] (0,0) rectangle (2.6,0.8);
\node at (1.3,0.4) {\textbf{action \textbackslash{} résultat}};
\foreach \x/\v in {3.0/见 jiàn, 4.6/懂 dǒng, 6.2/完 wán, 7.8/对 duì, 9.4/错 cuò, 11.0/到 dào}{
  \draw[fill=poporangeL, draw=poporange] (\x,0) rectangle (\x+1.4,0.8);
  \node at (\x+0.7,0.4) {\v};
}
\foreach \y/\a in {-0.9/看 kàn, -1.8/听 tīng, -2.7/写 xiě, -3.6/说 shuō, -4.5/买 mǎi, -5.4/找 zhǎo}{
  \draw[fill=popgreenL, draw=popgreen] (0,\y) rectangle (2.6,\y+0.8);
  \node at (1.3,\y+0.4) {\a};
}
\foreach \y in {-0.9,-1.8,-2.7,-3.6,-4.5,-5.4}{
  \foreach \x in {3.0,4.6,6.2,7.8,9.4,11.0}{
    \draw[draw=popline] (\x,\y) rectangle (\x+1.4,\y+0.8);
  }
}
\node at (3.7,-0.5) {看见};
\node at (5.3,-0.5) {看懂};
\node at (7.0,-0.5) {看完};
\node at (10.2,-0.5) {看错};
\node at (11.7,-0.5) {看到};
\node at (3.7,-1.4) {听见};
\node at (5.3,-1.4) {听懂};
\node at (7.0,-1.4) {听完};
\node at (11.7,-1.4) {听到};
\node at (7.0,-2.3) {写完};
\node at (8.6,-2.3) {写对};
\node at (10.2,-2.3) {写错};
\node at (7.0,-3.2) {说完};
\node at (10.2,-3.2) {说错};
\node at (11.7,-3.2) {说到};
\node at (7.0,-4.1) {买完};
\node at (8.6,-4.1) {买对};
\node at (10.2,-4.1) {买错};
\node at (11.7,-4.1) {买到};
\node at (11.7,-5.0) {找到};
\end{tikzpicture}
\end{center}
\legende{Tableau croisé : les verbes d'action (lignes) se combinent avec les résultats (colonnes). Les cases vides sont des combinaisons rares ou impossibles : on ne dit pas 说懂 ni 看对.}
\end{popfigure}

\begin{center}
\begin{tabularx}{\linewidth}{|X|X|X|X|}
\hline
\rowcolor{popblueL} \textbf{Caractères} & \textbf{Pinyin} & \textbf{Nature} & \textbf{Traduction} \\
\hline
看见 & kànjian & verbe + résultat & voir (regarder et percevoir) \\
\hline
听见 & tīngjian & verbe + résultat & entendre (écouter et percevoir) \\
\hline
买到 & mǎidào & verbe + résultat & réussir à acheter, trouver à acheter \\
\hline
找到 & zhǎodào & verbe + résultat & trouver (chercher et atteindre) \\
\hline
听懂 & tīngdǒng & verbe + résultat & comprendre en écoutant \\
\hline
看懂 & kàndǒng & verbe + résultat & comprendre en lisant \\
\hline
写完 & xiěwán & verbe + résultat & finir d'écrire \\
\hline
做对 & zuòduì & verbe + résultat & faire juste, réussir (un exercice) \\
\hline
做错 & zuòcuò & verbe + résultat & faire faux, se tromper \\
\hline
说清楚 & shuō qīngchu & verbe + résultat & expliquer clairement \\
\hline
\end{tabularx}
\end{center}

Remarquez la logique interne de chaque paire : \zh{看见} = « regarder-percevoir » donc \emph{voir} ; \zh{找到} = « chercher-atteindre » donc \emph{trouver}. Le chinois dit explicitement ce que le français laisse implicite : on ne « voit » que si l'action de regarder \emph{aboutit}. Notez aussi que dans \zh{看见} et \zh{听见}, le second caractère \zh{见} perd son ton (ton neutre : \emph{jiàn} → \emph{jian}).

\begin{methode}[Construire une phrase au complément de résultat en 5 étapes]
\begin{enumerate}
\item \textbf{Choisir le verbe d'action} : que fait le sujet ? Ex. \zh{写} xiě « écrire ».
\item \textbf{Choisir le résultat} : quel est l'aboutissement ? Ex. \zh{完} wán « fini ».
\item \textbf{Coller les deux} en un bloc inséparable : \zh{写完} xiěwán.
\item \textbf{Ajouter 了} si l'action est accomplie : \zh{写完了}. (Pas de 了 à l'impératif ou au futur.)
\item \textbf{Placer l'objet à la fin} : \zh{我写完了汉字。} \emph{Wǒ xiě wán le Hànzì.} « J'ai fini d'écrire les caractères. »
\end{enumerate}
Vérification finale : rien entre le verbe et le résultat ; 了 après le groupe ; l'objet en dernier.
\end{methode}

\begin{exoresolu}[Construire la phrase]
Énoncé : traduisez en chinois « Le professeur a expliqué clairement la leçon d'aujourd'hui, et les élèves ont tout compris en écoutant. » (verbes : 说 shuō « dire, expliquer » ; 听 tīng « écouter » ; résultats : 清楚 qīngchu « clair » ; 懂 dǒng « comprendre »).
\tcblower
Solution détaillée, étape par étape :
\begin{enumerate}
\item Première proposition : verbe d'action \zh{说}, résultat \zh{清楚} → bloc inséparable \zh{说清楚}. Action accomplie → on ajoute \zh{了} après le groupe. L'objet \zh{今天的课} « la leçon d'aujourd'hui » est long : la tournure la plus naturelle est la construction en \zh{把}, qui avance l'objet avant le verbe : \zh{老师把今天的课说清楚了。} (La tournure sans 把, \zh{老师说清楚了今天的课。}, est grammaticalement possible mais moins idiomatique.)
\item Seconde proposition : verbe \zh{听}, résultat \zh{懂} → bloc \zh{听懂} + \zh{了} ; « tout » se rend par \zh{都} dōu, placé \emph{avant} le verbe : \zh{学生们都听懂了。}
\end{enumerate}
Phrase complète : \zh{老师把今天的课说清楚了，学生们都听懂了。}\\
\emph{Lǎoshī bǎ jīntiān de kè shuō qīngchu le, xuéshengmen dōu tīng dǒng le.}\\
« Le professeur a expliqué clairement la leçon d'aujourd'hui, et les élèves ont tout compris en écoutant. »\\
Contrôle : aucun mot entre 说 et 清楚, ni entre 听 et 懂 ; 了 après chaque groupe ; dans la construction en 把, l'objet 今天的课 est bien placé entre 把 et le groupe Verbe + Résultat.
\end{exoresolu}

\begin{attention}[Piège : l'objet ne remonte jamais entre le verbe et le résultat]
Les francophones, calquant « j'ai fini \emph{mes devoirs} », écrivent parfois \zh{我做作业完了} ou \zh{我做了完作业}. Retenez la règle d'or : le groupe Verbe + Résultat est \textbf{soudé}, et l'objet attend sagement \textbf{à la fin} : \zh{我做完了作业}. Si l'on veut mettre l'objet en avant, on utilise \zh{把} : \zh{我把作业做完了} — mais jamais d'objet entre le verbe et son résultat. Même vigilance pour les compléments de lieu : \zh{在} + lieu se place \emph{avant} le verbe (\zh{我在市场买到了鱼}), jamais entre le verbe et le résultat.
\end{attention}

\begin{exercice}[À vous de construire]
Traduisez en chinois avec un complément de résultat, en suivant la méthode en cinq étapes :
\begin{enumerate}
\item J'ai fini mes devoirs. (做 + 完)
\item Elle a retrouvé son livre. (找 + 到)
\item Nous avons compris le professeur en l'écoutant. (听 + 懂)
\item Il a mal écrit ce caractère. (写 + 错)
\item Maman a réussi à acheter du poisson au marché de Douala. (买 + 到)
\end{enumerate}
\end{exercice}

\begin{corrige}[Exercice « À vous de construire »]
\begin{enumerate}
\item \zh{我做完了作业。} \emph{Wǒ zuò wán le zuòyè.}
\item \zh{她找到了她的书。} \emph{Tā zhǎodào le tā de shū.}
\item \zh{我们听懂了老师的话。} \emph{Wǒmen tīng dǒng le lǎoshī de huà.}
\item \zh{他写错了这个字。} \emph{Tā xiě cuò le zhège zì.}
\item \zh{妈妈在杜阿拉的市场买到了鱼。} \emph{Māma zài Dù'ālā de shìchǎng mǎi dào le yú.} (le complément de lieu \zh{在…} se place \emph{avant} le verbe, jamais entre verbe et résultat.)
\end{enumerate}
\end{corrige}

\begin{aretenir}[L'essentiel de la section]
\begin{itemize}
\item Schéma : \textbf{Sujet + [Verbe + Résultat] + 了 + Objet}.
\item Verbe et résultat forment un \textbf{bloc inséparable} : rien ne s'insère entre eux, pas même 了.
\item \zh{了} se place \textbf{après le groupe} ; l'objet se place \textbf{à la fin}.
\item \zh{完} ne s'emploie jamais seul avec un objet : il faut un verbe d'action devant (\zh{做完}, \zh{吃完}, \zh{写完}).
\item Avec \zh{把} : Sujet + 把 + Objet + Verbe + Résultat (+ 了).
\item Les paires fréquentes (\zh{看见}, \zh{听懂}, \zh{找到}, \zh{说清楚}…) s'apprennent comme des verbes composés ; les combinaisons impossibles (\zh{说懂}, \zh{看对}) doivent être évitées.
\end{itemize}
\end{aretenir}

\coursec{Négation, question et forme potentielle}

Dans la section précédente, nous avons construit le complément de résultat à l'affirmatif : \zh{我看懂了} (\emph{Wǒ kàndǒng le.} « J'ai lu et j'ai compris. »). Mais comment dire « je n'ai \emph{pas} compris » ? Comment demander « as-tu fini ? » ? Et surtout, comment distinguer « je n'ai pas vu » de « je n'arrive pas à voir » ? C'est l'objet de cette section : la \cle{négation}, la \cle{question} et la \cle{forme potentielle} du complément de résultat.

\courssub{Observation : deux négations qui ne disent pas la même chose}

\begin{textebox}[À la cantine du lycée de Yaoundé]
A ： 昨天的听力考试你听懂了吗？\\
\emph{Zuótiān de tīnglì kǎoshì nǐ tīngdǒng le ma ?}\\
« Hier, à l'épreuve de compréhension orale, as-tu compris ? »\\
B ： 我没听懂。录音太快了，我听不懂。\\
\emph{Wǒ méi tīngdǒng. Lùyīn tài kuài le, wǒ tīng bu dǒng.}\\
« Je n'ai pas compris. L'enregistrement était trop rapide, je n'arrivais pas à comprendre. »\\
A ： 那你看见黑板上的答案了吗？\\
\emph{Nà nǐ kànjian hēibǎn shang de dá'àn le ma ?}\\
« Alors, as-tu vu les réponses au tableau ? »\\
B ： 我坐在最后，看不见。你找到你的本子没有？\\
\emph{Wǒ zuò zài zuìhòu, kàn bu jiàn. Nǐ zhǎodào nǐ de běnzi méiyǒu ?}\\
« J'étais assis tout au fond, je n'arrivais pas à voir. As-tu retrouvé ton cahier ? »\\
A ： 找到了！在我的书包里。\\
\emph{Zhǎodào le ! Zài wǒ de shūbāo lǐ.}\\
« Retrouvé ! Il était dans mon sac. »
\end{textebox}

Observez attentivement : \zh{没听懂} et \zh{听不懂} se traduisent tous deux par une négation en français, mais ils ne veulent pas dire la même chose. De même, \zh{看不见} et \zh{没看见} sont très différents. C'est toute la logique de cette section.

\begin{center}
\begin{tabularx}{\linewidth}{|X|X|X|X|}
\hline
\rowcolor{popblueL} Caractères & Pinyin & Nature & Traduction \\
\hline
\zh{录音} & lùyīn & nom & enregistrement (audio) \\
\hline
\zh{听力} & tīnglì & nom & compréhension orale \\
\hline
\zh{答案} & dá'àn & nom & réponse, corrigé \\
\hline
\zh{本子} & běnzi & nom & cahier \\
\hline
\zh{书包} & shūbāo & nom & sac d'école, cartable \\
\hline
\zh{钥匙} & yàoshi & nom & clé \\
\hline
\zh{眼镜} & yǎnjìng & nom & lunettes \\
\hline
\zh{消息} & xiāoxi & nom & nouvelle, information \\
\hline
\end{tabularx}
\end{center}

\courssub{La négation : 没(有) + Verbe + Résultat}

\begin{propriete}[Règle de négation du complément de résultat]
Le complément de résultat se nie avec \cle{没(有)}, \textbf{jamais avec 不 seul} (sauf dans la forme potentielle, voir plus bas).\\
\textbf{Schéma :} Sujet + 没(有) + Verbe + Résultat (+ complément).\\
La négation porte sur \emph{l'obtention du résultat} : l'action a eu lieu (ou aurait dû avoir lieu), mais le résultat attendu \textbf{n'a pas été atteint}.
\end{propriete}

Pourquoi \zh{没} et non \zh{不} ? Rappelons la distinction fondamentale : \zh{不} nie un état, une habitude ou un fait présent/futur ; \zh{没(有)} nie la \emph{réalisation} d'un fait accompli. Or le complément de résultat exprime précisément un résultat \emph{obtenu ou non} : nier ce résultat, c'est constater qu'il ne s'est \emph{pas produit}. On utilise donc obligatoirement \zh{没(有)}.

\begin{exemplebox}[Exemples de négation avec 没(有)]
\zh{我没看见他。} \emph{Wǒ méi kànjian tā.} « Je ne l'ai pas vu. » (j'ai regardé, mais je ne l'ai pas aperçu)\\
\zh{作业还没做完。} \emph{Zuòyè hái méi zuòwán.} « Les devoirs ne sont pas encore finis. »\\
\zh{我没听懂，请你再说一遍。} \emph{Wǒ méi tīngdǒng, qǐng nǐ zài shuō yí biàn.} « Je n'ai pas compris, répétez s'il vous plaît. »\\
\zh{这个字我没写对。} \emph{Zhège zì wǒ méi xiěduì.} « Ce caractère, je ne l'ai pas écrit correctement. »\\
\zh{他没找到他的钥匙。} \emph{Tā méi zhǎodào tā de yàoshi.} « Il n'a pas retrouvé ses clés. »
\end{exemplebox}

Notez l'adverbe \zh{还} (\emph{hái}, « encore ») dans \zh{还没做完} : il signifie que le résultat n'est pas atteint \emph{pour l'instant}, mais qu'il pourrait l'être. C'est l'équivalent du français « pas encore ».

\begin{attention}[Erreur fréquente : 我不看见 ✗]
Les francophones traduisent mot à mot « je ne vois pas » par \zh{我不看见}, ce qui est \textbf{incorrect}. Le français « je ne vois pas » a deux sens possibles :
\begin{itemize}
\item \emph{constat} (j'ai cherché, mais je n'ai pas aperçu) → \zh{我没看见。} \emph{Wǒ méi kànjian.}
\item \emph{incapacité} (obstacle, distance, obscurité) → \zh{我看不见。} \emph{Wǒ kàn bu jiàn.}
\end{itemize}
Jamais \zh{不} directement devant « Verbe + Résultat » hors forme potentielle !
\end{attention}

\courssub{La question affirmative-négative}

Pour demander si un résultat a été atteint, le chinois utilise la question affirmative-négative. Deux formes coexistent :

\begin{propriete}[Question avec le complément de résultat]
\textbf{Forme complète :} Verbe + Résultat + 没 + Verbe + Résultat\\
\textbf{Forme abrégée (la plus courante) :} Verbe + Résultat + 没有 \emph{ou} Verbe + 没 + Verbe + Résultat\\
On peut aussi utiliser la particule \zh{吗} avec la forme affirmative : V + R + 了 + 吗？
\end{propriete}

\begin{exemplebox}[Poser la question]
\zh{你看没看见？} \emph{Nǐ kànjian méi kànjian ?} « As-tu vu ou pas ? » (forme complète)\\
\zh{你找到没有？} \emph{Nǐ zhǎodào méiyǒu ?} « As-tu trouvé ? » (forme abrégée)\\
\zh{作业做完了没有？} \emph{Zuòyè zuòwán le méiyǒu ?} « As-tu fini tes devoirs ? »\\
\zh{你听懂了吗？} \emph{Nǐ tīngdǒng le ma ?} « As-tu compris ? » (avec 吗)
\end{exemplebox}

Réponse courte : à l'affirmatif, on reprend V + R (+ \zh{了}) : \zh{找到了。} « Trouvé. » ; au négatif : \zh{没找到。} « Pas trouvé. »

\courssub{La forme potentielle : 得 et 不 entre le verbe et le résultat}

C'est la grande nouveauté de cette section. En insérant \zh{得} ou \zh{不} \emph{entre} le verbe et son complément de résultat, on exprime non plus un fait, mais une \cle{capacité} ou une \cle{possibilité} d'obtenir le résultat.

\begin{propriete}[Forme potentielle du complément de résultat]
\textbf{Potentiel affirmatif :} Verbe + 得 + Résultat = « capable d'obtenir le résultat », « pouvoir + V ».\\
\textbf{Potentiel négatif :} Verbe + 不 + Résultat = « impossible d'obtenir le résultat », « ne pas pouvoir / ne pas arriver à + V ».\\
Attention : dans la forme potentielle, on n'emploie \textbf{jamais} \zh{了} ni \zh{没}.
\end{propriete}

\begin{exemplebox}[Potentiel affirmatif : V + 得 + R]
\zh{我看得见。} \emph{Wǒ kàn de jiàn.} « J'arrive à voir. »\\
\zh{这个字我写得对。} \emph{Zhège zì wǒ xiě de duì.} « Ce caractère, je suis capable de l'écrire correctement. »\\
\zh{今天的作业我做得完。} \emph{Jīntiān de zuòyè wǒ zuò de wán.} « Les devoirs d'aujourd'hui, je peux les finir. »\\
\zh{老师说的话你听得懂吗？} \emph{Lǎoshī shuō de huà nǐ tīng de dǒng ma ?} « Comprends-tu ce que dit le professeur ? »
\end{exemplebox}

\begin{exemplebox}[Potentiel négatif : V + 不 + R]
\zh{太远了，我看不见。} \emph{Tài yuǎn le, wǒ kàn bu jiàn.} « C'est trop loin, je n'arrive pas à voir. »\\
\zh{听不懂。} \emph{Tīng bu dǒng.} « (Je) n'arrive pas à comprendre. »\\
\zh{没有字典，我翻译不完这篇文章。} \emph{Méiyǒu zìdiǎn, wǒ fānyì bu wán zhè piān wénzhāng.} « Sans dictionnaire, je ne peux pas finir de traduire cet article. »\\
\zh{这个词我记不住。} \emph{Zhège cí wǒ jì bu zhù.} « Ce mot, je n'arrive pas à le retenir. »
\end{exemplebox}

Remarque phonétique importante : au potentiel négatif, le complément de résultat prend généralement le \emph{ton léger} : \emph{kàn bu jiàn}, \emph{tīng bu dǒng}. Au potentiel affirmatif, \zh{得} est aussi atone : \emph{kàn de jiàn}.

\courssub{Le contraste capital : 没看见 ≠ 看不见}

\begin{methode}[Distinguer 没+V+R et V+不+R]
Pour choisir la bonne forme, posez-vous la question suivante :
\begin{enumerate}
\item Parlez-vous d'un \textbf{fait passé précis} dont le résultat ne s'est pas produit ? → \zh{没} + V + R : \zh{我昨天没看见他。} « Hier, je ne l'ai pas vu. »
\item Parlez-vous d'une \textbf{incapacité} (générale ou actuelle) à obtenir le résultat ? → V + \zh{不} + R : \zh{我看不见。} « Je n'arrive pas à voir. »
\item Testez avec le français : si « ne pas \emph{pouvoir} / ne pas \emph{arriver à} » convient, c'est le potentiel (V+不+R) ; si c'est un simple constat au passé composé (« je n'ai pas vu »), c'est \zh{没}+V+R.
\end{enumerate}
\end{methode}

\begin{center}
\begin{tabularx}{\linewidth}{|X|X|X|}
\hline
\rowcolor{popblueL} Forme & Exemple (caractères + pinyin) & Traduction \\
\hline
Affirmation (résultat obtenu) & \zh{我看见了。} \emph{Wǒ kànjian le.} & J'ai vu. \\
\hline
Négation (résultat non obtenu) & \zh{我没看见。} \emph{Wǒ méi kànjian.} & Je n'ai pas vu. \\
\hline
Question & \zh{你看见没有？} \emph{Nǐ kànjian méiyǒu ?} & As-tu vu ? \\
\hline
Potentiel affirmatif & \zh{我看得见。} \emph{Wǒ kàn de jiàn.} & J'arrive à voir. \\
\hline
Potentiel négatif & \zh{我看不见。} \emph{Wǒ kàn bu jiàn.} & Je n'arrive pas à voir. \\
\hline
\end{tabularx}
\end{center}

\begin{popfigure}
\begin{tikzpicture}[
  node distance=0.9cm,
  question/.style={diamond, draw=popdark, fill=popgoldL, align=center, inner sep=1pt, aspect=2.2, font=\small},
  result/.style={rectangle, rounded corners, draw=popdark, fill=popblueL, align=center, font=\small, inner sep=4pt},
  pot/.style={rectangle, rounded corners, draw=popdark, fill=popgreenL, align=center, font=\small, inner sep=4pt},
  fleche/.style={-{Stealth[length=2.5mm]}, thick, popdark}
]
\node[question, align=center] (q1){Le résultat\\est-il obtenu ?};
\node[result, below left=1cm and 0.4cm of q1, align=center] (oui){Oui : V + R + 了\\\zh{我看懂了。}};
\node[result, below right=1cm and 0.4cm of q1, align=center] (non){Non : 没 + V + R\\\zh{我没看懂。}};
\node[question, below=2.6cm of q1, align=center] (q2){Capacité à obtenir\\le résultat ?};
\node[pot, below left=1cm and 0.4cm of q2, align=center] (cap){Possible : V + 得 + R\\\zh{我看得懂。}};
\node[pot, below right=1cm and 0.4cm of q2, align=center] (incap){Impossible : V + 不 + R\\\zh{我看不懂。}};
\draw[fleche] (q1) -- node[above left, font=\footnotesize]{oui} (oui);
\draw[fleche] (q1) -- node[above right, font=\footnotesize]{non} (non);
\draw[fleche] (q2) -- node[above left, font=\footnotesize]{oui} (cap);
\draw[fleche] (q2) -- node[above right, font=\footnotesize]{non} (incap);
\end{tikzpicture}
\legende{Organigramme de choix : fait constaté (haut) ou capacité (bas) ?}
\end{popfigure}

\begin{exoresolu}[Traduire en chinois]
Énoncé : traduisez les phrases suivantes en choisissant la bonne forme. 1) « Je n'ai pas fini mes devoirs. » 2) « C'est trop difficile, je ne peux pas finir mes devoirs. » 3) « As-tu entendu cette nouvelle ? » 4) « Je n'arrive pas à retenir ces caractères. »
\tcblower
Solution détaillée :
\begin{enumerate}
\item Fait passé, résultat non atteint → \zh{没} + V + R : \zh{我没做完作业。} \emph{Wǒ méi zuòwán zuòyè.}
\item Incapacité (cause : difficulté) → V + \zh{不} + R : \zh{太难了，我做不完作业。} \emph{Tài nán le, wǒ zuò bu wán zuòyè.} La forme \zh{我没做完作业} serait grammaticalement correcte mais changerait le sens (« je ne les ai pas finis », simple constat).
\item Question sur un résultat → forme abrégée : \zh{你听到这个消息没有？} \emph{Nǐ tīngdào zhège xiāoxi méiyǒu ?}
\item Incapacité générale → V + \zh{不} + R : \zh{这些汉字我记不住。} \emph{Zhèxiē Hànzì wǒ jì bu zhù.}
\end{enumerate}
\end{exoresolu}

\begin{attention}[Pièges à éviter]
\begin{itemize}
\item Ne dites jamais \zh{我不做完} ni \zh{我不找到} : hors potentiel, la négation du résultat exige \zh{没(有)}.
\item Au potentiel, pas de \zh{了} : \zh{我看得见了} ✗ → \zh{我看得见。} ✓
\item Ne confondez pas le \zh{得} du potentiel avec le \zh{得} du complément de degré (\zh{他说得很好} « il parle bien ») : au potentiel, \zh{得} se place entre le verbe et un \emph{résultat} (\zh{说得完} « pouvoir finir de dire »).
\item \zh{看不见} (je n'arrive pas à voir) ≠ \zh{没看见} (je n'ai pas vu) : intervertir les deux change complètement le sens de la phrase.
\end{itemize}
\end{attention}

\begin{exercice}[À vous de jouer]
1) Traduisez : « Hier soir, je n'ai pas compris le texte. » ; « Sans lunettes, je n'arrive pas à voir le tableau. » ; « As-tu retrouvé ton stylo ? — Non, pas retrouvé. »\\
2) Transformez au potentiel négatif : \zh{我听懂了。} / \zh{她写完了。}\\
3) Choisissez la bonne forme (\zh{没} ou \zh{不}) : 我今天看\_\_\_见黑板，因为我坐在最后。
\end{exercice}

\begin{corrige}[Exercice « À vous de jouer »]
1) \zh{昨天晚上我没看懂课文。} \emph{Zuótiān wǎnshang wǒ méi kàndǒng kèwén.} ; \zh{不戴眼镜，我看不见黑板。} \emph{Bù dài yǎnjìng, wǒ kàn bu jiàn hēibǎn.} ; \zh{你找到你的笔没有？——没找到。} \emph{Nǐ zhǎodào nǐ de bǐ méiyǒu ? — Méi zhǎodào.}\\
2) \zh{我听不懂。} \emph{Wǒ tīng bu dǒng.} ; \zh{她写不完。} \emph{Tā xiě bu wán.}\\
3) \zh{不} : \zh{我今天看不见黑板，因为我坐在最后。} — il s'agit d'une incapacité actuelle (cause : la distance), donc potentiel négatif, et non d'un constat passé.
\end{corrige}

\begin{aretenir}[L'essentiel de la section]
\begin{itemize}
\item Négation du résultat : \zh{没(有)} + V + R, jamais \zh{不} seul : \zh{我没听懂。}
\item Question : V + R + 没有 (abrégé) ou V + R + 没 + V + R (complet) : \zh{你找到没有？}
\item Potentiel affirmatif : V + \zh{得} + R (capable) : \zh{我看得见。}
\item Potentiel négatif : V + \zh{不} + R (impossible) : \zh{我看不见。}
\item \zh{没看见} = constat passé ; \zh{看不见} = incapacité. Deux mondes différents !
\end{itemize}
\end{aretenir}

\coursec{Les fractions avec 分之}

Après avoir maîtrisé la négation \zh{没(有) + V + R} et la forme potentielle \zh{V + 得/不 + R} du complément de résultat, nous quittons le domaine des verbes pour celui des nombres. Dans le module « Environnement et santé », on lit souvent des phrases comme « un tiers des élèves », « les trois quarts de la planète ». Comment dit-on une fraction en chinois ? C'est l'objet de cette section : la structure avec \zh{分之}, dont l'ordre est \textbf{inversé} par rapport au français.

\courssub{Observation : un ordre à l'envers}

Lisons d'abord ces deux phrases tirées d'un exposé sur l'environnement :

\begin{exemplebox}[Observer avant de comprendre]
\zh{四分之三的地球表面是水。}\\
\emph{Sì fēn zhī sān de dìqiú biǎomiàn shì shuǐ.}\\
« Les trois quarts de la surface de la Terre sont de l'eau. »

\medskip
\zh{这个班五分之二的学生是女生。}\\
\emph{Zhège bān wǔ fēn zhī èr de xuésheng shì nǚshēng.}\\
« Deux cinquièmes des élèves de cette classe sont des filles. »
\end{exemplebox}

Regardons \zh{四分之三} de près. En français, « trois quarts » dit d'abord le \textbf{numérateur} (trois), puis le \textbf{dénominateur} (quarts). Le chinois fait exactement le contraire : \zh{四} (quatre) vient d'abord, puis \zh{分之}, puis \zh{三} (trois). Le chinois dit littéralement : « sur quatre parts, trois parts ».

\begin{definition}[Les termes d'une fraction]
\begin{itemize}
\item \zh{分母} \emph{fēnmǔ} : le dénominateur (le nombre de parts totales) ;
\item \zh{分子} \emph{fēnzǐ} : le numérateur (le nombre de parts prises) ;
\item \zh{分之} \emph{fēn zhī} : « parmi … parts », charnière obligatoire entre les deux.
\end{itemize}
\end{definition}

\courssub{La règle : dénominateur + 分之 + numérateur}

\begin{propriete}[Construction de la fraction chinoise]
\textbf{Schéma :} \cle{dénominateur} + \zh{分} + \zh{之} + \cle{numérateur}, c'est-à-dire \zh{分母} + \zh{分之} + \zh{分子}.

\medskip
En chinois, on dit d'abord le \textbf{DÉNOMINATEUR} (en combien de parts on divise), puis \zh{分之}, puis le \textbf{NUMÉRATEUR} (combien de parts on prend). C'est l'ordre \textbf{inverse} du français.

\medskip
Lecture mentale conseillée : « sur \emph{n} parts, \emph{m} parts ».
\begin{itemize}
\item \zh{三分之一} \emph{sān fēn zhī yī} : « sur 3 parts, 1 » = un tiers ;
\item \zh{四分之三} \emph{sì fēn zhī sān} : « sur 4 parts, 3 » = trois quarts ;
\item \zh{五分之二} \emph{wǔ fēn zhī èr} : « sur 5 parts, 2 » = deux cinquièmes.
\end{itemize}
\end{propriete}

\begin{center}
\begin{tabularx}{\linewidth}{|X|X|X|X|}
\hline
\rowcolor{popblueL} Fraction & Chinois & Pinyin & Lecture littérale \\
\hline
un demi & 二分之一 & èr fēn zhī yī & sur 2 parts, 1 \\
\hline
un tiers & 三分之一 & sān fēn zhī yī & sur 3 parts, 1 \\
\hline
deux tiers & 三分之二 & sān fēn zhī èr & sur 3 parts, 2 \\
\hline
un quart & 四分之一 & sì fēn zhī yī & sur 4 parts, 1 \\
\hline
trois quarts & 四分之三 & sì fēn zhī sān & sur 4 parts, 3 \\
\hline
deux cinquièmes & 五分之二 & wǔ fēn zhī èr & sur 5 parts, 2 \\
\hline
trois dixièmes & 十分之三 & shí fēn zhī sān & sur 10 parts, 3 \\
\hline
\end{tabularx}
\end{center}

\begin{popfigure}
\begin{tikzpicture}[font=\small]
\node[draw=popblue, fill=popblueL, rounded corners, align=center, minimum width=4.2cm, minimum height=1.1cm] (fr) at (0,2.6) {\textbf{Français : trois quarts}\\ numérateur $\rightarrow$ dénominateur};
\node[draw=poporange, fill=poporangeL, rounded corners, align=center, minimum width=4.2cm, minimum height=1.1cm] (zh) at (0,0) {\textbf{Chinois : 四分之三}\\ dénominateur $\rightarrow$ numérateur};
\node[draw=popdark, fill=white, rounded corners, align=center, minimum width=2.2cm] (num) at (6,2.6) {\textbf{三} (3)\\ numérateur};
\node[draw=popdark, fill=white, rounded corners, align=center, minimum width=2.2cm] (den) at (6,0) {\textbf{四} (4)\\ dénominateur};
\draw[-{Stealth[length=3mm]}, very thick, popblue] (fr.east) -- (num.west);
\draw[-{Stealth[length=3mm]}, very thick, popblue, dashed] (fr.south east) .. controls (3.2,1.6) .. (den.north west);
\draw[-{Stealth[length=3mm]}, very thick, poporange] (zh.east) -- (den.west);
\draw[-{Stealth[length=3mm]}, very thick, poporange, dashed] (zh.north east) .. controls (3.2,1.0) .. (num.south west);
\node[align=center, poppurple] at (3,1.3) {ordre\\ croisé !};
\end{tikzpicture}
\legende{Le français va du numérateur vers le dénominateur ; le chinois fait le chemin inverse : 四 (dénominateur) + 分之 + 三 (numérateur).}
\end{popfigure}

\courssub{Le cas particulier de « un demi » : 半}

Pour « un demi », la fraction régulière \zh{二分之一} \emph{èr fēn zhī yī} existe et est correcte, mais dans la langue courante on préfère \zh{一半} \emph{yíbàn} « une moitié » ou simplement \zh{半} \emph{bàn} « demi ». Retenez bien ce mot : il est très fréquent à l'examen.

\begin{exemplebox}[半 dans la vie quotidienne]
\zh{半个苹果} \emph{bàn ge píngguǒ} : « une demi-pomme »\\
\zh{半小时} \emph{bàn xiǎoshí} : « une demi-heure »\\
\zh{一半的学生喜欢运动。} \emph{Yíbàn de xuésheng xǐhuan yùndòng.} : « La moitié des élèves aiment le sport. »\\
\zh{我们班一半的同学每天喝水。} \emph{Wǒmen bān yíbàn de tóngxué měitiān hē shuǐ.} : « La moitié de nos camarades boivent de l'eau chaque jour. »
\end{exemplebox}

\begin{attention}[Place de 半]
\zh{半} se place \textbf{avant} le classificateur (\zh{半个苹果}, \zh{半碗饭} \emph{bàn wǎn fàn}), jamais après. On ne dit pas ✗ \zh{一个半苹果} pour « une demi-pomme » (cela signifierait « une pomme et demie »).
\end{attention}

\courssub{La fraction dans la phrase : deux positions}

\begin{propriete}[Emploi de la fraction dans la phrase]
\textbf{Position 1 — devant le nom, avec \zh{的} :} fraction + \zh{的} + nom. La fraction fonctionne comme un déterminant.

\zh{三分之一的学生病了。} \emph{Sān fēn zhī yī de xuésheng bìng le.} « Un tiers des élèves sont malades. »

\medskip
\textbf{Position 2 — après le verbe \zh{占} (zhàn, « occuper, représenter ») :} sujet + \zh{占} + fraction. Très utile pour présenter des statistiques.

\zh{女生占三分之二。} \emph{Nǚshēng zhàn sān fēn zhī èr.} « Les filles représentent les deux tiers. »
\end{propriete}

\begin{exemplebox}[Contexte santé et environnement]
\zh{三分之一的学生每天喝干净的水。}\\
\emph{Sān fēn zhī yī de xuésheng měitiān hē gānjìng de shuǐ.}\\
« Un tiers des élèves boivent de l'eau propre chaque jour. »

\medskip
\zh{四分之三的地球表面是水。}\\
\emph{Sì fēn zhī sān de dìqiú biǎomiàn shì shuǐ.}\\
« Les trois quarts de la surface de la Terre sont de l'eau. »

\medskip
\zh{我们学校五分之二的学生参加健康俱乐部。}\\
\emph{Wǒmen xuéxiào wǔ fēn zhī èr de xuésheng cānjiā jiànkāng jùlèbù.}\\
« Deux cinquièmes des élèves de notre école participent au club santé. »

\medskip
\zh{在杜阿拉，一半的人坐摩托车上班。}\\
\emph{Zài Dù'ālā, yíbàn de rén zuò mótuōchē shàngbān.}\\
« À Douala, la moitié des gens vont au travail en moto-taxi. »

\medskip
\zh{这个医院的病人占全市病人的四分之一。}\\
\emph{Zhège yīyuàn de bìngrén zhàn quán shì bìngrén de sì fēn zhī yī.}\\
« Les malades de cet hôpital représentent un quart des malades de toute la ville. »
\end{exemplebox}

\begin{popfigure}
\begin{tikzpicture}
\fill[popblueL] (0,0) circle (2cm);
\fill[poporange] (0,0) -- (0:2) arc (0:270:2) -- cycle;
\draw[popdark, thick] (0,0) circle (2cm);
\draw[popdark] (0,0) -- (90:2);
\draw[popdark] (0,0) -- (180:2);
\draw[popdark] (0,0) -- (270:2);
\node[white, font=\bfseries] at (45:1.2) {水};
\node[white, font=\bfseries] at (135:1.2) {水};
\node[white, font=\bfseries] at (270:1.2) {水};
\node[popdark, font=\bfseries] at (225:1.2) {陆地};
\node[align=center] at (0,-2.8) {四分之三的地球表面是水。\\ \emph{Sì fēn zhī sān de dìqiú biǎomiàn shì shuǐ.}};
\end{tikzpicture}
\legende{Camembert : les trois quarts (四分之三) de la surface de la Terre sont couverts d'eau ; un quart (四分之一) est de la terre ferme.}
\end{popfigure}

\courssub{Les pourcentages avec 百分之}

\begin{propriete}[Construction du pourcentage]
\textbf{Schéma :} \zh{百分之} + \cle{numérateur} (le dénominateur 100 est figé dans \zh{百分之}).

\medskip
\zh{百分之} \emph{bǎi fēn zhī} signifie littéralement « sur cent parts ». On n'énonce que le numérateur.
\begin{itemize}
\item \zh{百分之五十} \emph{bǎi fēn zhī wǔshí} : 50 \% (cinquante pour cent) ;
\item \zh{百分之百} \emph{bǎi fēn zhī bǎi} : 100 \% (cent pour cent) ;
\item \zh{百分之三十} \emph{bǎi fēn zhī sānshí} : 30 \% ;
\item \zh{百分之五} \emph{bǎi fēn zhī wǔ} : 5 \%.
\end{itemize}
\end{propriete}

\begin{exemplebox}[Pourcentages en contexte santé/environnement]
\zh{百分之七十的地球表面是水。}\\
\emph{Bǎi fēn zhī qīshí de dìqiú biǎomiàn shì shuǐ.}\\
« 70 \% de la surface de la Terre sont de l'eau. »

\medskip
\zh{这个医院百分之九十的病人康复了。}\\
\emph{Zhège yīyuàn bǎi fēn zhī jiǔshí de bìngrén kāngfù le.}\\
« 90 \% des malades de cet hôpital sont guéris. »

\medskip
\zh{在雅温得，百分之六十的人走路上班。}\\
\emph{Zài Yāwēndé, bǎi fēn zhī liùshí de rén zǒu lù shàngbān.}\\
« À Yaoundé, 60 \% des gens vont au travail à pied. »
\end{exemplebox}

\begin{attention}[Ne pas confondre fraction et pourcentage]
\zh{十分之三} = 3/10 = 30 \% = \zh{百分之三十}.
L'ordre reste : \textbf{base (100) + 分之 + chiffre}. Ne jamais dire ✗ \zh{三十百分之}.
\end{attention}

\courssub{Comparaison avec le français : tableau de synthèse}

\begin{center}
\begin{tabularx}{\linewidth}{|X|X|X|}
\hline
\rowcolor{popblueL} Notion & Français & Chinois \\
\hline
Ordre (fraction) & numérateur + dénominateur (trois quarts) & dénominateur + 分之 + numérateur (四分之三) \\
\hline
Ordre (pourcentage) & chiffre + « pour cent » & 百分之 + chiffre \\
\hline
Mot-charnière & « sur » (3 sur 4), souvent implicite & 分之, toujours obligatoire \\
\hline
un demi & « la moitié » & 一半 (courant) ou 二分之一 \\
\hline
devant le nom & « un tiers \emph{des} élèves » & 三分之一\emph{的}学生 / 百分之三十\emph{的}学生 \\
\hline
statistique & « représenter les deux tiers » & 占三分之二 / 占百分之六十 \\
\hline
\end{tabularx}
\end{center}

\begin{attention}[L'erreur n° 1 des francophones : garder l'ordre français]
Pour dire « trois quarts », beaucoup d'élèves écrivent ✗ \zh{三分之三} en calquant « \textbf{trois} quarts ». C'est faux : \zh{三分之三} signifie « trois tiers », c'est-à-dire le tout !

\medskip
\textbf{Méthode anti-erreur :} demandez-vous toujours « en combien de parts je coupe ? ». Ce nombre vient \textbf{en premier}. « Trois quarts » = je coupe en \textbf{4}, je prends \textbf{3} → \zh{四分之三}.

\medskip
Autre piège : ne pas oublier \zh{的} entre la fraction et le nom : ✗ \zh{三分之一学生} → ✓ \zh{三分之一的学生}. Et n'oubliez jamais \zh{之} : ✗ \zh{四分三} → ✓ \zh{四分之三}.
\end{attention}

\begin{exoresolu}[Traduire et construire des fractions et pourcentages]
\textbf{Énoncé.} 1) Traduisez en chinois : « Les deux tiers des élèves de Première A sont malades aujourd'hui. » 2) Traduisez en français : \zh{我们学校四分之三的学生吃得很健康。} 3) Écrivez en chinois la fraction « trois cinquièmes » et le pourcentage « 80 \% ».
\tcblower
\textbf{Solution détaillée.}

1) « Deux tiers » : je coupe en \textbf{3}, je prends \textbf{2} → \zh{三分之二}. Devant le nom \zh{学生}, il faut \zh{的}. Phrase complète :\\
\zh{今天，高一A班三分之二的学生病了。}\\
\emph{Jīntiān, Gāoyī A bān sān fēn zhī èr de xuésheng bìng le.}

2) \zh{四分之三} = trois quarts ; \zh{吃得很健康} = « mangent très sainement » (complément de degré avec 得). Traduction : « Les trois quarts des élèves de notre école mangent très sainement. »

3) « Trois cinquièmes » : je coupe en \textbf{5}, je prends \textbf{3} → \zh{五分之三} \emph{wǔ fēn zhī sān}. (Le piège était d'écrire ✗ \zh{三分之五}.)\\
« 80 \% » → \zh{百分之八十} \emph{bǎi fēn zhī bāshí}.
\end{exoresolu}

\begin{exercice}[À vous de jouer]
1) Écrivez en chinois (caractères + pinyin) : un quart ; quatre cinquièmes ; un dixième ; la moitié (forme courante) ; 25 \% ; 100 \%.\\
2) Traduisez en chinois : « Un quart des élèves ne boit pas assez d'eau. »\\
3) Traduisez en français : \zh{这个学校一半的老师是女老师。}\\
4) Transformez avec \zh{占} : \zh{三分之一的学生是女生。}\\
5) Transformez avec \zh{占} : \zh{百分之二十的同学喜欢吃蔬菜。}
\end{exercice}

\begin{corrige}[Exercice « À vous de jouer »]
1) \zh{四分之一} \emph{sì fēn zhī yī} ; \zh{五分之四} \emph{wǔ fēn zhī sì} ; \zh{十分之一} \emph{shí fēn zhī yī} ; \zh{一半} \emph{yíbàn} ; \zh{百分之二十五} \emph{bǎi fēn zhī èrshíwǔ} ; \zh{百分之百} \emph{bǎi fēn zhī bǎi}.\\
2) \zh{四分之一的学生喝水喝得不够。} \emph{Sì fēn zhī yī de xuésheng hē shuǐ hē de bù gòu.} (Notez le complément de degré \zh{喝得不够} : belle révision !)\\
3) « La moitié des professeurs de cette école sont des femmes. »\\
4) \zh{女生占三分之一。} \emph{Nǚshēng zhàn sān fēn zhī yī.} « Les filles représentent un tiers. »\\
5) \zh{喜欢吃蔬菜的同学占百分之二十。} \emph{Xǐhuan chī shūcài de tóngxué zhàn bǎi fēn zhī èrshí.} « Les camarades qui aiment manger des légumes représentent 20 \%. »
\end{corrige}

\begin{savaistu}[Pourquoi 分之 ?]
Le caractère \zh{分} (fēn) signifie « diviser, partager » : on le retrouve dans \zh{分钟} (minute, une « part » d'heure), \zh{分开} (séparer) et \zh{十分} (très, littéralement « dix parts » = à 100 \%). Quant à \zh{之} (zhī), c'est un mot-outil hérité du chinois littéraire classique, équivalent ancien de \zh{的}. La fraction \zh{三分之一} se lit donc littéralement : « de trois parts divisées, une ». En Chine comme au Cameroun, les fractions servent partout : au marché de Mokolo, on négocie « la moitié » (\zh{一半}) du prix ; dans les bulletins météo chinois, on annonce quelle fraction du pays sera couverte de pluie. Retenir \zh{分} = « diviser » vous aidera à ne plus jamais inverser l'ordre : on annonce d'abord en combien de parts on \emph{divise} !
\end{savaistu}

\begin{aretenir}[L'essentiel de la section]
\begin{itemize}
\item Fraction chinoise = \cle{dénominateur} + \zh{分之} + \cle{numérateur} : ordre inverse du français.
\item \zh{三分之一} = un tiers ; \zh{四分之三} = trois quarts ; \zh{五分之二} = deux cinquièmes.
\item Pourcentage = \zh{百分之} + chiffre (dénominateur 100 implicite).
\item « Un demi » : \zh{二分之一} ou, plus courant, \zh{一半} ; \zh{半} se place avant le classificateur (\zh{半小时}).
\item Devant un nom : fraction/pourcentage + \zh{的} + nom ; dans une statistique : sujet + \zh{占} + fraction/pourcentage.
\item Réflexe : « en combien de parts je coupe ? » → ce nombre d'abord.
\end{itemize}
\end{aretenir}

\coursec{Les pourcentages avec 百分之}

Dans la section précédente, nous avons appris à construire les fractions chinoises avec la structure \zh{分母 + 分之 + 分子} (dénominateur + 分之 + numérateur) : \zh{三分之一} \emph{(sān fēn zhī yī)} « un tiers », \zh{四分之三} \emph{(sì fēn zhī sān)} « trois quarts ». Nous allons maintenant étudier un cas particulier extrêmement fréquent : la fraction dont le dénominateur est 100, c'est-à-dire le \cle{pourcentage}. En chinois, il suffit de remplacer le dénominateur par \zh{百} (bǎi, « cent ») : \zh{百分之七十} \emph{(bǎi fēn zhī qīshí)} = « 70 \% ». Cette structure est indispensable pour lire et produire des statistiques sur la santé et l'environnement, le thème de notre module.

\courssub{Observation : lire les pourcentages dans un texte}

\begin{textebox}[L'eau et la santé dans notre ville]
\zh{在我们城市，百分之七十的疾病来自脏水。百分之六十的人每天喝不干净的水。但是，情况在变化：去年，百分之四十的家庭有干净的水，今年已经有百分之五十五了。医生说，如果百分之一百的家庭都能喝干净的水，百分之八十的疾病会消失。}\\
\emph{Zài wǒmen chéngshì, bǎi fēn zhī qīshí de jíbìng láizì zāng shuǐ. Bǎi fēn zhī liùshí de rén měitiān hē bù gānjìng de shuǐ. Dànshì, qíngkuàng zài biànhuà : qùnián, bǎi fēn zhī sìshí de jiātíng yǒu gānjìng de shuǐ, jīnnián yǐjīng yǒu bǎi fēn zhī wǔshíwǔ le. Yīshēng shuō, rúguǒ bǎi fēn zhī yībǎi de jiātíng dōu néng hē gānjìng de shuǐ, bǎi fēn zhī bāshí de jíbìng huì xiāoshī.}\\
« Dans notre ville, 70 \% des maladies viennent de l'eau sale. 60 \% des gens boivent chaque jour de l'eau non potable. Mais la situation change : l'année dernière, 40 \% des familles avaient de l'eau propre ; cette année, on en est déjà à 55 \%. Le médecin dit que si 100 \% des familles pouvaient boire de l'eau potable, 80 \% des maladies disparaîtraient. »
\end{textebox}

\textbf{Questions d'observation :}
\begin{enumerate}
\item Repérez tous les groupes qui commencent par \zh{百分之}. Quel nombre suit chaque groupe ?
\item Dans \zh{百分之七十的疾病}, quel mot-outil relie le pourcentage au nom \zh{疾病} ?
\item Comparez avec le français : dans « soixante-dix pour cent », où se trouve le nombre ? Et dans \zh{百分之七十} ?
\end{enumerate}

On observe que le pourcentage chinois se construit exactement comme une fraction, avec \zh{百} en position de dénominateur : \zh{百 + 分 + 之 + nombre}. Le nombre (le numérateur) vient \textbf{en dernier}, alors qu'en français il vient en premier (« soixante-dix pour cent »). C'est la différence essentielle à retenir.

\courssub{La règle : 百分之 + nombre}

\begin{propriete}[Règle du pourcentage]
\textbf{Schéma :} \zh{百 + 分 + 之 + [nombre]} = « [nombre] pour cent ».\\
Un pourcentage est une fraction sur 100 : on remplace simplement le dénominateur par \zh{百} (bǎi). Les trois caractères \zh{百分之} forment un bloc inséparable, et le nombre se place \textbf{après}, jamais avant.\\
\zh{百分之七十} \emph{(bǎi fēn zhī qīshí)} = 70 \% ; \zh{百分之五} \emph{(bǎi fēn zhī wǔ)} = 5 \%.
\end{propriete}

\begin{exemplebox}[Lecture des cas particuliers]
\begin{itemize}
\item \zh{百分之一} \emph{(bǎi fēn zhī yī)} = 1 \%
\item \zh{百分之五} \emph{(bǎi fēn zhī wǔ)} = 5 \%
\item \zh{百分之十} \emph{(bǎi fēn zhī shí)} = 10 \%
\item \zh{百分之二十五} \emph{(bǎi fēn zhī èrshíwǔ)} = 25 \%
\item \zh{百分之五十} \emph{(bǎi fēn zhī wǔshí)} = 50 \% (= \zh{一半} yībàn, « la moitié »)
\item \zh{百分之九十九} \emph{(bǎi fēn zhī jiǔshíjiǔ)} = 99 \%
\item \zh{百分之一百} \emph{(bǎi fēn zhī yībǎi)} = 100 \%
\end{itemize}
\end{exemplebox}

\begin{attention}[Le cas de 100 \%]
Pour dire « 100 \% », on ne dit jamais \zh{百分之百百} ✗ : le nombre complet est \zh{一百} (yībǎi), donc la forme correcte est \zh{百分之一百} \emph{(bǎi fēn zhī yībǎi)} ✓. De même, « 0 \% » se dit \zh{百分之零} \emph{(bǎi fēn zhī líng)}.
\end{attention}

\begin{popfigure}
\begin{tikzpicture}
\node[draw=popblue, fill=popblueL, rounded corners, minimum width=1.1cm, minimum height=0.9cm, font=\Large] (bai) at (0,0) {百};
\node[draw=popblue, fill=popblueL, rounded corners, minimum width=1.1cm, minimum height=0.9cm, font=\Large] (fen) at (1.5,0) {分};
\node[draw=popblue, fill=popblueL, rounded corners, minimum width=1.1cm, minimum height=0.9cm, font=\Large] (zhi) at (3,0) {之};
\node[draw=poporange, fill=poporangeL, rounded corners, minimum width=1.6cm, minimum height=0.9cm, font=\Large] (x) at (5,0) {X};
\node[below=0.15cm of bai, align=center, font=\small] {« sur\\ cent »};
\node[below=0.15cm of fen, align=center, font=\small] {part};
\node[below=0.15cm of zhi, align=center, font=\small] {de};
\node[below=0.15cm of x, align=center, font=\small, text=poporange] {le nombre\\ = X \%};
\draw[-{Stealth}, thick, popdark] (6.3,0) -- (7.3,0);
\node[draw=popgreen, fill=popgreenL, rounded corners, align=center, font=\small] at (9.6,0) {百分之七十\\ = 70 \%};
% barre de progression
\draw[fill=popcream, draw=popline] (0,-2.2) rectangle (10,-1.7);
\foreach \i in {0,...,10} {
  \draw[popline] (\i,-2.2) -- (\i,-1.7);
}
\fill[popblue!50] (0,-2.2) rectangle (7,-1.7);
\node[below, font=\scriptsize] at (1,-2.25) {百分之十};
\node[below, font=\scriptsize] at (5,-2.25) {百分之五十};
\node[below, font=\scriptsize] at (10,-2.25) {百分之一百};
\node[above, font=\small\bfseries, text=popblue] at (3.5,-1.65) {百分之七十 = 70 \%};
\end{tikzpicture}
\legende{Schéma de la structure 百分之 + X et barre graduée de 10 \% à 100 \% (la zone colorée représente 70 \%).}
\end{popfigure}

\courssub{Emplois du pourcentage dans la phrase}

\textbf{1. Devant un nom, avec le mot-outil \zh{的}.} Le pourcentage qualifie le nom comme un adjectif : \zh{百分之 + X + 的 + nom}.

\begin{exemplebox}[Pourcentage + 的 + nom]
\zh{百分之八十的病人已经好了。} \emph{(Bǎi fēn zhī bāshí de bìngrén yǐjīng hǎo le.)} « 80 \% des malades sont déjà guéris. »\\
\zh{这个城市百分之六十的水不干净。} \emph{(Zhège chéngshì bǎi fēn zhī liùshí de shuǐ bù gānjìng.)} « 60 \% de l'eau de cette ville n'est pas potable. »\\
\zh{我们班百分之九十的学生通过了考试。} \emph{(Wǒmen bān bǎi fēn zhī jiǔshí de xuésheng tōngguò le kǎoshì.)} « 90 \% des élèves de notre classe ont réussi l'examen. »
\end{exemplebox}

\textbf{2. Avec le verbe \zh{占} (zhàn, « représenter, occuper »).} Très fréquent dans les statistiques : Sujet + \zh{占} + (complément) + pourcentage.

\begin{exemplebox}[Avec 占]
\zh{农村人口占全国人口的百分之五十五。} \emph{(Nóngcūn rénkǒu zhàn quánguó rénkǒu de bǎi fēn zhī wǔshíwǔ.)} « La population rurale représente 55 \% de la population du pays. »\\
\zh{水占人体的百分之七十。} \emph{(Shuǐ zhàn réntǐ de bǎi fēn zhī qīshí.)} « L'eau représente 70 \% du corps humain. »
\end{exemplebox}

\textbf{3. Après les verbes d'augmentation et de diminution.} Avec \zh{增加} (zēngjiā, « augmenter ») et \zh{减少} (jiǎnshǎo, « diminuer »), le pourcentage suit directement le verbe, sans \zh{的}, car il n'est pas devant un nom.

\begin{exemplebox}[Augmenter et diminuer de X \%]
\zh{今年，种树的面积增加了百分之二十。} \emph{(Jīnnián, zhòng shù de miànjī zēngjiā le bǎi fēn zhī èrshí.)} « Cette année, la surface plantée d'arbres a augmenté de 20 \%. »\\
\zh{上个月，这个城市的病人减少了百分之十。} \emph{(Shàng ge yuè, zhège chéngshì de bìngrén jiǎnshǎo le bǎi fēn zhī shí.)} « Le mois dernier, le nombre de malades dans cette ville a diminué de 10 \%. »
\end{exemplebox}

\courssub{Tableau récapitulatif et comparaison avec le français}

\begin{center}
\begin{tabularx}{\linewidth}{|X|X|X|X|}
\hline
\rowcolor{popblueL} Caractères & Pinyin & Nature & Traduction \\
\hline
百分之 & bǎi fēn zhī & bloc fixe & pour cent (de) \\
\hline
百分之七十 & bǎi fēn zhī qīshí & expression numérale & 70 \% \\
\hline
百分之一百 & bǎi fēn zhī yībǎi & expression numérale & 100 \% \\
\hline
占 & zhàn & verbe & représenter, occuper \\
\hline
增加 & zēngjiā & verbe & augmenter \\
\hline
减少 & jiǎnshǎo & verbe & diminuer \\
\hline
疾病 & jíbìng & nom & maladie \\
\hline
人口 & rénkǒu & nom & population \\
\hline
来自 & láizì & verbe & venir de, provenir de \\
\hline
面积 & miànjī & nom & surface, superficie \\
\hline
\end{tabularx}
\end{center}

\begin{center}
\begin{tabularx}{\linewidth}{|X|X|X|}
\hline
\rowcolor{popblueL} Structure & Exemple (caractères + pinyin) & Traduction \\
\hline
百分之 + nombre + 的 + nom & 百分之八十的学生 \emph{bǎi fēn zhī bāshí de xuésheng} & 80 \% des élèves \\
\hline
Sujet + 占 + 百分之 + nombre & 水占百分之七十 \emph{shuǐ zhàn bǎi fēn zhī qīshí} & l'eau représente 70 \% \\
\hline
verbe + 了 + 百分之 + nombre & 增加了百分之二十 \emph{zēngjiā le bǎi fēn zhī èrshí} & a augmenté de 20 \% \\
\hline
\end{tabularx}
\end{center}

\textbf{Comparaison des deux langues :} en français, « 70 \% » se lit « \emph{soixante-dix} pour cent » : le nombre vient \textbf{d'abord}. En chinois, c'est l'inverse : \zh{百分之} vient d'abord, le nombre \textbf{ensuite} : \zh{百分之七十}. Le chinois suit l'ordre logique de la fraction (« sur cent parts, soixante-dix »), exactement comme \zh{十分之七} « sept dixièmes ». Retenez : \textbf{le nombre est toujours en dernière position}.

\courssub{Exercice résolu et pièges à éviter}

\begin{exoresolu}[Traduisez en chinois]
Énoncé : Traduisez en chinois : a) 90 \% des élèves de notre classe ont réussi l'examen. b) 30 \% des maladies viennent de l'air sale. c) La population de la ville a augmenté de 15 \%.
\tcblower
Solution détaillée :
\begin{enumerate}
\item a) Pourcentage + \zh{的} + nom, en tête de phrase : \zh{我们班百分之九十的学生通过了考试。} \emph{(Wǒmen bān bǎi fēn zhī jiǔshí de xuésheng tōngguò le kǎoshì.)} On place d'abord le cadre \zh{我们班}, puis le pourcentage, puis \zh{的}, puis le nom.
\item b) \zh{百分之三十的疾病来自脏空气。} \emph{(Bǎi fēn zhī sānshí de jíbìng láizì zāng kōngqì.)} Le verbe \zh{来自} (láizì) signifie « venir de ».
\item c) \zh{城市的人口增加了百分之十五。} \emph{(Chéngshì de rénkǒu zēngjiā le bǎi fēn zhī shíwǔ.)} Le pourcentage suit directement le verbe \zh{增加} + \zh{了}, sans \zh{的}, car il n'est pas devant un nom.
\end{enumerate}
\end{exoresolu}

\begin{attention}[Erreurs fréquentes des francophones]
\begin{itemize}
\item \textbf{Inverser l'ordre} sous l'influence du français « soixante-dix pour cent » : \zh{七十百分之} ✗. La bonne forme : \zh{百分之七十} ✓. Le bloc \zh{百分之} vient toujours en premier.
\item \textbf{Oublier \zh{之}} : \zh{百分七十} ✗. Les trois caractères \zh{百分之} sont obligatoires et inséparables devant le nombre.
\item \textbf{Oublier \zh{的}} devant le nom : \zh{百分之八十学生} ✗ ; il faut \zh{百分之八十的学生} ✓.
\item \textbf{Confondre avec les fractions} : \zh{三分之七十} ✗ n'existe pas ; pour 70 \%, le dénominateur est toujours \zh{百}.
\item \textbf{Mettre \zh{的} après un verbe} : \zh{增加了的百分之二十} ✗ ; après \zh{增加}/\zh{减少}, le pourcentage se place directement, sans \zh{的}.
\end{itemize}
\end{attention}

\begin{savaistu}[Pourcentages et soldes en Chine]
Dans les statistiques officielles chinoises, on utilise aussi le symbole \%, qui se lit toujours \zh{百分之} + nombre : « 70\% » se lit \zh{百分之七十}. Mais dans les magasins, les soldes s'expriment autrement : on dit \zh{打八折} \emph{(dǎ bā zhé)}, littéralement « faire huit dixièmes », c'est-à-dire payer 80 \% du prix, donc bénéficier de 20 \% de réduction ! Attention au piège : \zh{打八折} ne veut pas dire « moins 80 \% ». C'est une autre façon, très concrète, de parler de proportions dans la vie quotidienne chinoise.
\end{savaistu}

\begin{exercice}[À vous de jouer]
\begin{enumerate}
\item Écrivez en chinois : 5 \% ; 25 \% ; 50 \% ; 99 \% ; 100 \%.
\item Traduisez : \zh{百分之六十的人每天喝干净的水。}
\item Traduisez en chinois : « 40 \% des familles de ce village ont de l'eau potable. »
\item Corrigez la phrase fautive : \zh{我们班九十百分之的学生喜欢汉语。}
\item Traduisez en chinois : « Cette année, le nombre d'élèves de notre école a augmenté de 10 \%. »
\end{enumerate}
\end{exercice}

\begin{corrige}[Exercice « À vous de jouer »]
\begin{enumerate}
\item \zh{百分之五} (bǎi fēn zhī wǔ) ; \zh{百分之二十五} (bǎi fēn zhī èrshíwǔ) ; \zh{百分之五十} (bǎi fēn zhī wǔshí) ; \zh{百分之九十九} (bǎi fēn zhī jiǔshíjiǔ) ; \zh{百分之一百} (bǎi fēn zhī yībǎi).
\item « 60 \% des gens boivent chaque jour de l'eau potable. »
\item \zh{这个村子百分之四十的家庭有干净的水。} \emph{(Zhège cūnzi bǎi fēn zhī sìshí de jiātíng yǒu gānjìng de shuǐ.)}
\item \zh{我们班百分之九十的学生喜欢汉语。} \emph{(Wǒmen bān bǎi fēn zhī jiǔshí de xuésheng xǐhuan Hànyǔ.)} — le bloc \zh{百分之} doit précéder le nombre \zh{九十}.
\item \zh{今年，我们学校的学生增加了百分之十。} \emph{(Jīnnián, wǒmen xuéxiào de xuésheng zēngjiā le bǎi fēn zhī shí.)} — après \zh{增加} + \zh{了}, pas de \zh{的}.
\end{enumerate}
\end{corrige}

\begin{aretenir}[L'essentiel de la section]
Le pourcentage chinois suit le schéma \cle{百分之 + nombre} : le bloc \zh{百分之} (« sur cent parts ») vient d'abord, le nombre ensuite — ordre inverse du français. Devant un nom, on ajoute \zh{的} : \zh{百分之八十的学生} « 80 \% des élèves ». Avec les statistiques, on utilise souvent le verbe \zh{占} : \zh{水占人体的百分之七十}. Après \zh{增加}/\zh{减少} + \zh{了}, le pourcentage se place directement : \zh{增加了百分之二十} « a augmenté de 20 \% ». Cas limites : \zh{百分之一} (1 \%) et \zh{百分之一百} (100 \%).
\end{aretenir}

\coursec{Comparaison avec le français et synthèse}

Nous venons de voir que les pourcentages se construisent avec \zh{百分之} suivi du chiffre : \zh{百分之七十} (\emph{bǎi fēn zhī qīshí}, « soixante-dix pour cent »). Cette construction a une particularité remarquable : son ordre est exactement \cle{inversé} par rapport au français. Cette dernière section fait le point : nous allons comparer systématiquement le français et le chinois pour les deux grandes structures de la leçon (le complément de résultat et les fractions), récapituler les pièges qui guettent les francophones, puis rassembler tout cela dans une carte mentale de synthèse.

\courssub{Le complément de résultat : deux verbes en français, un seul groupe en chinois}

\textbf{Observation.} Comparez les paires de phrases suivantes :

\begin{exemplebox}[Observer avant de conclure]
\zh{我做完了作业。} \emph{Wǒ zuòwán le zuòyè.} « J'ai \textbf{fini de faire} mes devoirs. »\\
\zh{我看见了王老师。} \emph{Wǒ kànjian le Wáng lǎoshī.} « J'ai \textbf{aperçu} le professeur Wang. »\\
\zh{他买到了火车票。} \emph{Tā mǎidào le huǒchēpiào.} « Il a \textbf{réussi à acheter} son billet de train. »
\end{exemplebox}

Que remarquez-vous ? En français, il faut souvent \textbf{deux verbes} (« finir de faire », « réussir à acheter ») ou un verbe plus précis (« apercevoir » au lieu de « voir »). En chinois, on garde le verbe de base (\zh{做} faire, \zh{看} regarder, \zh{买} acheter) et on lui \textbf{colle} un petit mot de résultat : \zh{完} (terminé), \zh{见} (perçu), \zh{到} (obtenu). Le français \emph{analyse} l'action en deux temps ; le chinois \emph{soude} le résultat au verbe.

\begin{propriete}[Principe de comparaison]
En français, le résultat d'une action est exprimé par une périphrase (\emph{finir de + infinitif}, \emph{réussir à + infinitif}) ou par un adverbe (\emph{bien}, \emph{complètement}). En chinois, le résultat forme avec le verbe un \cle{groupe inséparable} : \textbf{V + complément de résultat}. Rien ne peut se glisser entre les deux, et la particule \zh{了} se place \emph{après} tout le groupe, jamais entre le verbe et son résultat.
\end{propriete}

\begin{center}
\begin{tabularx}{\linewidth}{|X|X|X|}
\hline
\rowcolor{popblueL} \textbf{Français} & \textbf{Chinois (caractères + pinyin)} & \textbf{Analyse} \\
\hline
finir de faire ses devoirs & \zh{做完作业} \emph{zuòwán zuòyè} & faire + terminé \\
\hline
réussir à acheter & \zh{买到} \emph{mǎidào} & acheter + obtenu \\
\hline
apercevoir (voir et percevoir) & \zh{看见} \emph{kànjian} & regarder + perçu \\
\hline
comprendre en écoutant & \zh{听懂} \emph{tīngdǒng} & écouter + compris \\
\hline
manger complètement, tout finir & \zh{吃完} \emph{chīwán} & manger + terminé \\
\hline
écrire correctement, sans faute & \zh{写对} \emph{xiěduì} & écrire + juste \\
\hline
\end{tabularx}
\end{center}

\begin{attention}[L'adverbe français ne devient pas un adverbe chinois]
Le français dit « il a \emph{complètement} mangé » : l'adverbe est un mot séparé. En chinois, on ne dit pas \zh{他完全吃了} dans ce sens ; on dit \zh{他吃完了} (\emph{tā chīwán le}) : le résultat \zh{完} est \textbf{collé} au verbe. De même, « elle a bien entendu » (au sens de « elle a saisi ce qui a été dit ») se dit \zh{她听见了} (\emph{tā tīngjian le}), sans adverbe séparé.
\end{attention}

\courssub{Fractions et pourcentages : l'ordre inversé}

\textbf{Observation.} Lisez à voix haute :

\begin{exemplebox}[Le miroir inversé]
\zh{四分之三} \emph{sì fēn zhī sān} « trois quarts » (litt. « de quatre parts, trois »)\\
\zh{百分之七十} \emph{bǎi fēn zhī qīshí} « soixante-dix pour cent » (litt. « de cent parts, soixante-dix »)
\end{exemplebox}

\begin{propriete}[Règle de l'ordre inversé]
Le français énonce d'abord le \textbf{numérateur}, puis le \textbf{dénominateur} : « \emph{trois} quarts », « \emph{soixante-dix} pour \emph{cent} ». Le chinois fait l'inverse : \textbf{dénominateur + 分之 + numérateur}. Le mot-outil \zh{之} est obligatoire : c'est lui qui signifie « parmi, de ». Schéma : \cle{B 分之 A} = « A sur B ». Cas particulier : « 100 \% » se dit \zh{百分之百} (\emph{bǎi fēn zhī bǎi}), et « la moitié » se dit couramment \zh{一半} (\emph{yíbàn}) plutôt que \zh{二分之一} dans la langue parlée.
\end{propriete}

\begin{center}
\begin{tabularx}{\linewidth}{|X|X|X|}
\hline
\rowcolor{popblueL} \textbf{Français (numérateur/dénominateur)} & \textbf{Chinois (dénominateur 分之 numérateur)} & \textbf{Pinyin} \\
\hline
un demi (1/2) & \zh{二分之一} & \emph{èr fēn zhī yī} \\
\hline
un tiers (1/3) & \zh{三分之一} & \emph{sān fēn zhī yī} \\
\hline
trois quarts (3/4) & \zh{四分之三} & \emph{sì fēn zhī sān} \\
\hline
deux cinquièmes (2/5) & \zh{五分之二} & \emph{wǔ fēn zhī èr} \\
\hline
25 \% & \zh{百分之二十五} & \emph{bǎi fēn zhī èrshíwǔ} \\
\hline
90 \% & \zh{百分之九十} & \emph{bǎi fēn zhī jiǔshí} \\
\hline
100 \% & \zh{百分之百} & \emph{bǎi fēn zhī bǎi} \\
\hline
\end{tabularx}
\end{center}

\begin{aretenir}[L'astuce du pont]
Imaginez un pont : \zh{之} est le pont entre les deux nombres. Le français traverse le pont du petit nombre vers le grand (« trois \emph{sur} quatre ») ; le chinois le traverse en sens inverse, du grand vers le petit (\zh{四分之三} : « de quatre parts, trois »). Retenez : \textbf{le chinois commence toujours par le dénominateur}.
\end{aretenir}

\courssub{Bilan des pièges : les quatre erreurs fréquentes des francophones}

\begin{attention}[Récapitulatif des erreurs à bannir]
\begin{enumerate}
\item \textbf{\zh{了} mal placé} : \zh{我做了完作业} ✗ → \zh{我做完了作业} ✓ (\emph{Wǒ zuòwán le zuòyè.}). La particule \zh{了} vient \emph{après} le groupe V + résultat, jamais entre les deux.
\item \textbf{\zh{不} au lieu de \zh{没}} : \zh{我不看见} ✗ → \zh{我没看见} ✓ (\emph{Wǒ méi kànjian.}, « je n'ai pas vu »). Au passé, la négation du résultat se fait avec \zh{没(有)}. Attention à la nuance : \zh{我看不见} (\emph{wǒ kànbujian}) signifie « je n'arrive pas à voir » (forme potentielle), ce n'est pas la même chose.
\item \textbf{Fraction inversée} : \zh{三分之四} ✗ pour « trois quarts » → \zh{四分之三} ✓. Le dénominateur d'abord !
\item \textbf{Pourcentage inversé} : \zh{七十百分之} ✗ → \zh{百分之七十} ✓. \zh{百分之} est un bloc fixe qui se place \emph{avant} le chiffre.
\end{enumerate}
\end{attention}

\begin{methode}[Vérifier sa phrase en quatre questions]
Avant de rendre une phrase ou une traduction, posez-vous ces quatre contrôles :
\begin{enumerate}
\item Le résultat (\zh{完}, \zh{到}, \zh{见}, \zh{懂}…) est-il \textbf{collé} au verbe, sans rien entre les deux ?
\item La particule \zh{了} est-elle placée \textbf{après} tout le groupe verbe + résultat ?
\item La négation au passé est-elle bien \zh{没(有)}, et non \zh{不} ?
\item La fraction ou le pourcentage commence-t-il par le \textbf{dénominateur} (\zh{四分之三}, \zh{百分之七十}) ?
\end{enumerate}
Si les quatre réponses sont « oui », votre phrase est correcte.
\end{methode}

\begin{exoresolu}[Traduction bilan]
\textbf{Énoncé.} Traduisez en chinois : « Quatre-vingt-dix pour cent des élèves de Première A ont fini leurs devoirs de chinois. »
\tcblower
\textbf{Solution détaillée.} Étape 1 : le pourcentage. « 90 \% » = dénominateur \zh{百} + \zh{分之} + numérateur \zh{九十} : \zh{百分之九十}. Étape 2 : « des élèves de Première A » : \zh{高一A班的学生} (\emph{gāoyī A bān de xuésheng}). Étape 3 : « ont fini leurs devoirs » : verbe \zh{做} + résultat \zh{完} + \zh{了} après le groupe : \zh{做完了作业}. Phrase complète : \zh{百分之九十的高一A班学生做完了中文作业。} \emph{Bǎi fēn zhī jiǔshí de gāoyī A bān xuésheng zuòwán le Zhōngwén zuòyè.} Vérification : résultat collé au verbe ✓ ; \zh{了} après le groupe ✓ ; pourcentage dans le bon ordre ✓.
\end{exoresolu}

\begin{exemplebox}[Phrase-bilan combinant les deux structures]
\zh{百分之七十五的学生做完了练习。}\\
\emph{Bǎi fēn zhī qīshíwǔ de xuésheng zuòwán le liànxí.}\\
« 75 \% des élèves ont fini les exercices. »\\
Analyse : \zh{百分之七十五} (pourcentage, ordre inversé) + \zh{的学生} (des élèves) + \zh{做完} (verbe + résultat) + \zh{了} (accompli, après le groupe) + \zh{练习} (les exercices).
\end{exemplebox}

\begin{exercice}[À vous de jouer]
Traduisez en chinois : 1) « Un tiers des élèves a compris la leçon. » 2) « Je n'ai pas fini mes devoirs. » 3) « 60 \% des élèves de notre école aiment le football. » Puis corrigez la phrase fautive : \zh{我买了到票。}
\end{exercice}

\begin{corrige}[Exercice « À vous de jouer »]
1) \zh{三分之一的学生听懂了这一课。} \emph{Sān fēn zhī yī de xuésheng tīngdǒng le zhè yī kè.} 2) \zh{我没做完作业。} \emph{Wǒ méi zuòwán zuòyè.} (négation au passé : \zh{没}). 3) \zh{我们学校百分之六十的学生喜欢足球。} \emph{Wǒmen xuéxiào bǎi fēn zhī liùshí de xuésheng xǐhuan zúqiú.} Correction : \zh{我买了到票} ✗ → \zh{我买到票了} ✓ : le résultat \zh{到} doit être collé au verbe, et \zh{了} se place après le groupe.
\end{corrige}

\courssub{Carte mentale de synthèse}

\begin{popfigure}
\begin{tikzpicture}[mindmap, grow cyclic, every node/.style={concept, align=center, font=\small},
root concept/.append style={concept color=popblue, font=\small\bfseries},
level 1/.append style={level distance=3.4cm, sibling angle=120},
level 2/.append style={level distance=2.6cm, sibling angle=50}]
\node[root concept]{Leçon 14}
  child[concept color=poporange]{ node[align=center]{V + résultat\\ \zh{做完} \zh{看见} \zh{买到}}
    child[concept color=poporangeL]{ node[align=center]{Schéma\\ V + résultat\\ (+ \zh{了})} }
    child[concept color=poporangeL]{ node[align=center]{Négation\\ \zh{没} + V + résultat} }
    child[concept color=poporangeL]{ node[align=center]{Potentiel\\ \zh{看得见} / \zh{看不见}} }
  }
  child[concept color=popgreen]{ node[align=center]{Fractions\\ \zh{分之}}
    child[concept color=popgreenL]{ node[align=center]{Dénominateur\\ d'abord\\ \zh{四分之三}} }
    child[concept color=popgreenL]{ node[align=center]{Pourcentage\\ \zh{百分之} + chiffre\\ \zh{百分之七十}} }
  };
\end{tikzpicture}
\legende{Carte mentale récapitulative : le complément de résultat et les fractions avec \zh{分之}.}
\end{popfigure}

\begin{aretenir}[L'essentiel de la leçon]
\begin{itemize}
\item Le chinois \textbf{colle} le résultat au verbe : \zh{做完}, \zh{看见}, \zh{买到}, \zh{听懂}. Le français utilise deux verbes ou un adverbe.
\item \zh{了} se place après le groupe entier ; la négation au passé est \zh{没(有)} ; la forme potentielle s'obtient avec \zh{得} (possible) et \zh{不} (impossible) : \zh{看得见} / \zh{看不见}.
\item Fractions et pourcentages : \textbf{ordre inversé} par rapport au français — dénominateur + \zh{之} + numérateur : \zh{四分之三}, \zh{百分之九十}.
\end{itemize}
\end{aretenir}

\begin{savaistu}[Les pourcentages, un langage de santé publique]
Au Cameroun comme en Chine, les campagnes de santé publique s'expriment en pourcentages : taux de vaccination, couverture sanitaire, réussite scolaire. Savoir lire \zh{百分之九十的儿童打了疫苗} (\emph{bǎi fēn zhī jiǔshí de értóng dǎ le yìmiáo}, « 90 \% des enfants ont été vaccinés »), c'est pouvoir comprendre, à Yaoundé comme à Beijing, les affiches du ministère de la Santé… en chinois ! Les chiffres deviennent ainsi un vrai outil de citoyenneté, et pas seulement un exercice de grammaire.
\end{savaistu}

\newpage

\coursec{Activité d'intégration}

\courssub{Objectifs de l'activité}

À la fin de cette activité d'intégration, tu seras capable de :
\begin{itemize}
\item poser des questions fermées avec le \cle{complément de résultat} (\zh{你做完了作业没有？}) pour recueillir des informations réelles dans la classe ;
\item employer la \cle{forme potentielle} (\zh{听得懂} / \zh{听不懂}) pour dire si un résultat est possible ou impossible ;
\item compter les réponses et exprimer les résultats avec les \cle{fractions} (\zh{三分之一}, \zh{四分之三}) et les \cle{pourcentages} (\zh{百分之八十}) ;
\item rédiger un court \cle{rapport d'enquête} de 6 à 8 phrases, correct dans l'ordre des mots et les mots-outils ;
\item présenter oralement ce rapport avec une prononciation claire des tons et réagir aux questions des camarades ;
\item évaluer la clarté et la correction grammaticale d'une production d'autrui (esprit critique linguistique).
\end{itemize}

\courssub{Situation}

Le club de santé scolaire de ton lycée de Yaoundé prépare la « Semaine de la santé ». La présidente du club, qui apprend aussi le chinois, veut publier dans le journal bilingue du lycée un petit article chinois intitulé « Les chiffres parlent chinois ». Elle demande à chaque groupe de la classe de Première A de mener un mini-sondage sur les habitudes des élèves : devoirs, compréhension des leçons, eau propre à la maison, petit-déjeuner, sport.

Voici la fiche de mission remise par la présidente du club :

\begin{textebox}[Fiche de mission du club de santé — 健康俱乐部任务单]
健康俱乐部任务单 \\
\emph{Jiànkāng jùlèbù rènwudān.} \\
« Fiche de mission du club de santé. » \\[4pt]
同学们好！为了健康周，请你们做一个小调查。 \\
\emph{Tóngxuémen hǎo! Wèile jiànkāng zhōu, qǐng nǐmen zuò yí ge xiǎo diàochá.} \\
« Bonjour à tous ! Pour la Semaine de la santé, faites un petit sondage. » \\[4pt]
问题一：你昨天晚上做完了作业没有？ \\
\emph{Wèntí yī : Nǐ zuótiān wǎnshang zuò wán le zuòyè méiyou?} \\
« Question 1 : As-tu fini tes devoirs hier soir ? » \\[4pt]
问题二：今天的语法，你听得懂吗？ \\
\emph{Wèntí èr : Jīntiān de yǔfǎ, nǐ tīng de dǒng ma?} \\
« Question 2 : La grammaire d'aujourd'hui, peux-tu la comprendre en écoutant ? » \\[4pt]
问题三：你在家喝干净的水吗？ \\
\emph{Wèntí sān : Nǐ zài jiā hē gānjìng de shuǐ ma?} \\
« Question 3 : Bois-tu de l'eau propre à la maison ? » \\[4pt]
请数好答案，用分数和百分数写报告，六到八句话。 \\
\emph{Qǐng shǔ hǎo dá'àn, yòng fēnshù hé bǎifēnshù xiě bàogào, liù dào bā jù huà.} \\
« Comptez bien les réponses, rédigez un rapport avec des fractions et des pourcentages, six à huit phrases. »
\end{textebox}

\begin{attention}[Lecture guidée de la fiche de mission]
Observe les deux premières questions : elles n'utilisent pas la même structure.
\begin{itemize}
\item \zh{你做完了作业没有？} : question au \emph{passé accompli} avec le complément de résultat \zh{完} — on demande si l'action a \emph{abouti}.
\item \zh{你听得懂吗？} : question avec la \emph{forme potentielle} \zh{听得懂} — on demande si le résultat est \emph{possible}. C'est la deuxième grande structure de la leçon.
\end{itemize}
Repère aussi \zh{数好答案} : \zh{好} est ici un complément de résultat (« compter \emph{bien}, jusqu'au bon résultat »).
\end{attention}

\begin{savaistu}[Santé et eau : Cameroun et Chine]
Au Cameroun comme en Chine, l'accès à l'eau potable est un grand sujet de santé publique. En Chine, on boit traditionnellement de l'eau \cle{bouillie} (\zh{开水} \emph{kāishuǐ}, littéralement « eau ouverte », c'est-à-dire portée à ébullition) plutôt que de l'eau froide du robinet : c'est pourquoi les écoles chinoises ont souvent des distributeurs d'eau chaude. Au Cameroun, beaucoup de familles filtrent ou font bouillir l'eau avant de la boire. Parler de ces habitudes en chinois, avec des chiffres, rend la leçon de grammaire concrète et utile.
\end{savaistu}

\courssub{Consignes}

Travail par groupes de 4. Répartissez les rôles : un(e) enquêteur/trice qui pose les questions, un(e) secrétaire qui note, un(e) comptable qui calcule les fractions et pourcentages, un(e) rapporteur/euse qui présente.

\begin{enumerate}
\item \textbf{Préparer les questions (10 min).} Avec ton groupe, formule à l'écrit 4 questions fermées (réponse oui/non) : deux avec le complément de résultat à la forme interrogative (\zh{Verbe + complément + 了 + 没有？}) et deux questions simples avec \zh{吗} (dont une à la forme potentielle : \zh{你听得懂吗？}). Vérifie l'ordre des mots avant de commencer.
\item \textbf{Mener le sondage (10 min).} Pose tes questions à au moins 12 élèves de la classe. Note les réponses « oui » et « non » dans un tableau. Attention : réponds aux questions des autres groupes toi aussi !
\item \textbf{Compter et convertir (10 min).} Compte les réponses. Exprime chaque résultat sous deux formes : une fraction (\zh{三分之一}, \zh{四分之三}…) et un pourcentage (\zh{百分之…}). Rappel : en chinois, on dit d'abord le dénominateur, puis le numérateur : \zh{三分之一} = littéralement « de trois parts, une ».
\item \textbf{Rédiger le rapport (15 min).} Écris un rapport de 6 à 8 phrases commençant par \zh{我们班…}. Utilise : au moins \textbf{deux fractions}, au moins \textbf{deux pourcentages}, au moins \textbf{deux compléments de résultat} (forme affirmative ou négative : \zh{做完了} / \zh{没听懂}).
\item \textbf{Présenter à l'oral (10 min).} Le rapporteur de chaque groupe lit le rapport à la classe, en soignant les tons. Les autres groupes écoutent et notent une phrase correcte et une erreur éventuelle.
\item \textbf{Voter et justifier (5 min).} La classe vote pour le rapport le plus clair et le plus correct. Justifie ton vote en français : ordre des mots, place de \zh{没}, bonne formation des fractions.
\end{enumerate}

\courssub{Ressources à mobiliser}

\begin{propriete}[Rappel 1 : le complément de résultat — schéma et emplois]
Le \cle{complément de résultat} est un mot placé \emph{immédiatement après le verbe} qui indique que l'action \emph{aboutit} à un résultat : \zh{完} « fini », \zh{懂} « compris », \zh{见} « aperçu », \zh{到} « atteint », \zh{对} « juste », \zh{干净} « propre », \zh{好} « bien, achevé ».\\[4pt]
\textbf{Schéma affirmatif :} Sujet + Verbe + complément + \zh{了} (+ complément d'objet).\\
\zh{我做完了作业。} \emph{Wǒ zuò wán le zuòyè.} « J'ai fini mes devoirs. »\\[4pt]
\textbf{Schéma négatif :} Sujet + \zh{没(有)} + Verbe + complément (+ objet). La négation du résultat exige \zh{没}, jamais \zh{不}.\\
\zh{我没听懂。} \emph{Wǒ méi tīng dǒng.} « Je n'ai pas compris (en écoutant). »\\[4pt]
\textbf{Schéma interrogatif :} Sujet + Verbe + complément + \zh{了} + \zh{没有} ?\\
\zh{你听懂了没有？} \emph{Nǐ tīng dǒng le méiyou?} « As-tu compris ? »\\[4pt]
\textbf{Exception importante :} si le complément d'objet est long ou mis en avant, on peut le placer en tête de phrase : \zh{作业，我做完了。} « Les devoirs, je les ai finis. »
\end{propriete}

\begin{propriete}[Rappel 2 : la forme potentielle du complément de résultat]
Pour dire qu'un résultat est \emph{possible} ou \emph{impossible}, on insère \zh{得} (possible) ou \zh{不} (impossible) \emph{entre le verbe et le complément}.\\[4pt]
\textbf{Schéma :} Sujet + Verbe + \zh{得}/\zh{不} + complément (+ objet).\\
\zh{我听得懂。} \emph{Wǒ tīng de dǒng.} « Je peux comprendre (en écoutant). »\\
\zh{我听不懂。} \emph{Wǒ tīng bu dǒng.} « Je ne peux pas comprendre (c'est trop difficile, trop rapide). »\\
\zh{这个字你看得见吗？} \emph{Zhè ge zì nǐ kàn de jiàn ma?} « Ce caractère, arrives-tu à le voir ? »\\[4pt]
\textbf{Attention à la nuance :} \zh{我没听懂} = « je n'ai pas compris » (constat d'un fait passé) ; \zh{我听不懂} = « je n'arrive pas à comprendre » (impossibilité, présent ou général). C'est ici que \zh{不} apparaît légitimement : à la forme potentielle seulement.
\end{propriete}

\begin{propriete}[Rappel 3 : fractions et pourcentages avec 分之]
\textbf{Fraction :} dénominateur + \zh{分之} + numérateur. L'ordre est \emph{inverse} du français.\\
\zh{三分之一} \emph{sān fēn zhī yī} « un tiers » (litt. « de trois parts, une ») ; \zh{四分之三} \emph{sì fēn zhī sān} « trois quarts ».\\[4pt]
\textbf{Pourcentage :} \zh{百分之} + nombre (le dénominateur 100 est déjà dit par \zh{百}).\\
\zh{百分之八十} \emph{bǎi fēn zhī bāshí} « 80 \% » ; \zh{百分之百} \emph{bǎi fēn zhī bǎi} « 100 \% ».\\[4pt]
\textbf{Place dans la phrase :} fraction/pourcentage + \zh{的} + nom. Le mot-outil \zh{的} est obligatoire.\\
\zh{我们班百分之八十的同学做完了作业。} \emph{Wǒmen bān bǎi fēn zhī bāshí de tóngxué zuò wán le zuòyè.} « 80 \% des élèves de notre classe ont fini leurs devoirs. »
\end{propriete}

\begin{methode}[Convertir un résultat de sondage en chinois, pas à pas]
\begin{enumerate}
\item \textbf{Compte} le total des réponses et le nombre de « oui » : par exemple 9 « oui » sur 12 élèves.
\item \textbf{Simplifie} la fraction si possible : 9/12 = 3/4.
\item \textbf{Forme la fraction chinoise} : dénominateur d'abord, numérateur ensuite : \zh{四分之三}.
\item \textbf{Forme le pourcentage} : 3/4 = 75 \% → \zh{百分之七十五}.
\item \textbf{Construis la phrase} : fraction/pourcentage + \zh{的} + \zh{同学} + verbe + complément : \zh{四分之三的同学做完了作业。}
\item \textbf{Vérifie} : ordre \zh{分之} respecté ? \zh{的} présent ? négation avec \zh{没} ?
\end{enumerate}
\end{methode}

\begin{center}
\begin{tabularx}{\linewidth}{|X|X|X|X|}
\hline
\rowcolor{popblueL} Caractères & Pinyin & Nature & Traduction \\
\hline
调查 & diàochá & n./v. & enquête, sondage ; enquêter \\
\hline
报告 & bàogào & n. & rapport \\
\hline
分数 & fēnshù & n. & fraction ; note \\
\hline
百分数 & bǎifēnshù & n. & pourcentage \\
\hline
同学 & tóngxué & n. & camarade de classe \\
\hline
干净 & gānjìng & adj. & propre \\
\hline
做完 & zuò wán & v. + compl. & finir de faire \\
\hline
听懂 & tīng dǒng & v. + compl. & comprendre en écoutant \\
\hline
看见 & kàn jiàn & v. + compl. & voir (apercevoir) \\
\hline
喝 & hē & v. & boire \\
\hline
问题 & wèntí & n. & question, problème \\
\hline
答案 & dá'àn & n. & réponse (à une question) \\
\hline
\end{tabularx}
\end{center}

\begin{attention}[Erreurs fréquentes des francophones]
\begin{itemize}
\item Ne dis pas \zh{不听懂} pour « je n'ai pas compris » : la négation du résultat constaté exige \zh{没} : \zh{我没听懂。} Réserve \zh{不} à la forme potentielle : \zh{我听不懂。} « je n'arrive pas à comprendre ».
\item Ne place pas l'objet entre le verbe et son complément : dis \zh{我做完了作业}, jamais \zh{我做作业完了}.
\item Dans \zh{三分之一}, l'ordre est inverse du français : le chinois dit « de trois parts, une ». Ne dis jamais \zh{一分之三}.
\item \zh{百分之} sert pour les pourcentages seulement ; pour « un tiers », dis \zh{三分之一}, pas \zh{百分之三}.
\item Après la fraction ou le pourcentage, n'oublie pas \zh{的} avant le nom : \zh{四分之三的同学}.
\item À la forme potentielle, \zh{得} et \zh{不} se placent \emph{entre} le verbe et le complément : \zh{听得懂}, \zh{听不懂} — jamais \zh{听懂得} ni \zh{不听懂}.
\end{itemize}
\end{attention}

\courssub{Proposition de production et corrigé commenté}

\begin{exoresolu}[Rapport d'enquête modèle — Groupe 3, Lycée de Yaoundé]
Rédige le rapport de ton groupe (6 à 8 phrases, 2 fractions, 2 pourcentages, 2 compléments de résultat).
\tcblower
\textbf{Production modèle :}\\[4pt]
我们班的小调查 \\
\emph{Wǒmen bān de xiǎo diàochá.} « Le petit sondage de notre classe. » \\[4pt]
上个星期二，我们组问了十二个同学三个问题。 \\
\emph{Shàng ge xīngqī'èr, wǒmen zǔ wèn le shí'èr ge tóngxué sān ge wèntí.} \\
« Mardi dernier, notre groupe a posé trois questions à douze camarades. » \\[4pt]
百分之七十五的同学做完了作业，四分之一的同学没做完。 \\
\emph{Bǎi fēn zhī qīshíwǔ de tóngxué zuò wán le zuòyè, sì fēn zhī yī de tóngxué méi zuò wán.} \\
« 75 \% des camarades ont fini leurs devoirs, un quart ne les a pas finis. » \\[4pt]
三分之二的同学听懂了语法，但是三分之一的同学没听懂。 \\
\emph{Sān fēn zhī èr de tóngxué tīng dǒng le yǔfǎ, dànshì sān fēn zhī yī de tóngxué méi tīng dǒng.} \\
« Les deux tiers ont compris la grammaire, mais un tiers n'a pas compris. » \\[4pt]
十二分之十一的同学在家喝干净的水。 \\
\emph{Shí'èr fēn zhī shíyī de tóngxué zài jiā hē gānjìng de shuǐ.} \\
« Onze douzièmes des camarades boivent de l'eau propre à la maison. » \\[4pt]
我们的调查做完了，谢谢大家！ \\
\emph{Wǒmen de diàochá zuò wán le, xièxie dàjiā!} \\
« Notre enquête est terminée, merci à tous ! » \\[6pt]
\textbf{Commentaire des choix de langue :}
\begin{itemize}
\item \zh{做完了} et \zh{听懂} : compléments de résultat correctement placés juste après le verbe ; la négation utilise bien \zh{没} (\zh{没做完}, \zh{没听懂}), jamais \zh{不}.
\item \zh{四分之一}, \zh{三分之二}, \zh{三分之一}, \zh{十二分之十一} : dénominateur avant \zh{分之}, numérateur après — ordre chinois respecté. Les chiffres sont cohérents avec le sondage : 9/12 = 3/4 = 75 \%, 8/12 = 2/3, 11/12.
\item \zh{百分之七十五} : pourcentage formé avec \zh{百分之 + nombre}, suivi de \zh{的 + 同学}.
\item La phrase finale réutilise \zh{做完了} avec humour (« l'enquête est finie ») : belle réutilisation de la structure dans un autre contexte.
\end{itemize}
\end{exoresolu}

\begin{exercice}[Entraînement rapide avant le sondage]
\begin{enumerate}
\item Mets à la forme négative : \zh{我看見了老师。} → \ldots
\item Mets à la forme potentielle (possible / impossible) : \zh{做完} → \ldots / \ldots
\item Traduis en chinois : « Les trois quarts des élèves ont compris. »
\item Traduis en chinois : « 60 \% des camarades boivent de l'eau propre. »
\item Corrige la faute : \zh{一分之二的同学没做完作业。}
\end{enumerate}
\end{exercice}

\begin{corrige}[Entraînement rapide]
\begin{enumerate}
\item \zh{我没看见老师。} \emph{Wǒ méi kàn jiàn lǎoshī.} « Je n'ai pas vu le professeur. » (négation avec \zh{没}, sans \zh{了}).
\item \zh{做得完} \emph{zuò de wán} « possible de finir » ; \zh{做不完} \emph{zuò bu wán} « impossible de finir ».
\item \zh{四分之三的同学听懂了。} \emph{Sì fēn zhī sān de tóngxué tīng dǒng le.}
\item \zh{百分之六十的同学喝干净的水。} \emph{Bǎi fēn zhī liùshí de tóngxué hē gānjìng de shuǐ.}
\item \zh{二分之一的同学没做完作业。} — le dénominateur (\zh{二}) précède \zh{分之}, le numérateur (\zh{一}) suit.
\end{enumerate}
\end{corrige}

\courssub{Bilan de l'activité}

Cette activité t'a permis de mobiliser en situation réelle les deux grammaires de la leçon : le \cle{complément de résultat} pour poser des questions et rapporter des faits accomplis (\zh{做完了}, \zh{听懂了}, \zh{没听懂}), sa \cle{forme potentielle} pour exprimer la possibilité (\zh{听得懂}, \zh{听不懂}), et les \cle{fractions et pourcentages} pour exprimer des résultats chiffrés (\zh{三分之一}, \zh{百分之八十}). Tu as également pratiqué les quatre compétences : compréhension (écouter les réponses), expression orale (sondage et présentation), expression écrite (rapport) et calcul en chinois. Pour t'auto-évaluer, vérifie que ton rapport contient : au moins deux fractions bien formées (dénominateur + \zh{分之} + numérateur), deux pourcentages avec \zh{百分之}, deux compléments de résultat dont un à la forme négative avec \zh{没}, et le mot-outil \zh{的} après chaque fraction ou pourcentage. Ces structures te seront très utiles pour commenter des statistiques dans les leçons du Module 2 sur l'environnement et la santé.

\newpage

\coursec{Méthodes et exercices résolus}

\courssub{Observation : découvrir le complément de résultat et les fractions}

\begin{textebox}[À l'infirmerie du lycée]
\zh{老师：你怎么了？}\\
\emph{Lǎoshī : Nǐ zěnme le?}\\
« Professeur : Qu'est-ce que tu as ? »\\
\zh{学生：老师，我没听清楚您的问题，请再说一遍。}\\
\emph{Xuésheng : Lǎoshī, wǒ méi tīng qīngchu nín de wèntí, qǐng zài shuō yí biàn.}\\
« Élève : Professeur, je n'ai pas bien entendu votre question, répétez s'il vous plaît. »\\
\zh{老师：昨天晚上的作业，你做完了吗？}\\
\emph{Lǎoshī : Zuótiān wǎnshang de zuòyè, nǐ zuòwán le ma?}\\
« Professeur : As-tu fini les devoirs d'hier soir ? »\\
\zh{学生：做完了。我还看见了一个通知：从明天开始，百分之八十的学生要参加体育活动。}\\
\emph{Xuésheng : Zuòwán le. Wǒ hái kànjian le yí ge tōngzhī : cóng míngtiān kāishǐ, bǎi fēn zhī bāshí de xuésheng yào cānjiā tǐyù huódòng.}\\
« Élève : Je les ai finis. J'ai aussi vu une annonce : à partir de demain, 80 \% des élèves doivent participer aux activités sportives. »\\
\zh{老师：很好！我们班三分之二的同学都喜欢运动，这很健康。}\\
\emph{Lǎoshī : Hěn hǎo! Wǒmen bān sān fēn zhī èr de tóngxué dōu xǐhuan yùndòng, zhè hěn jiànkāng.}\\
« Professeur : Très bien ! Les deux tiers des camarades de notre classe aiment le sport, c'est très sain. »
\end{textebox}

\textbf{Questions d'observation.}
\begin{enumerate}
\item Dans \zh{听清楚}, \zh{做完}, \zh{看见}, quel mot indique l'action et quel mot indique son résultat ?
\item Où se place le mot de résultat par rapport au verbe ? Peut-on mettre quelque chose entre les deux ?
\item Dans \zh{百分之八十} et \zh{三分之二}, quel nombre vient en premier ? Comparez avec le français « 80 \% » et « deux tiers ».
\end{enumerate}

\begin{definition}[Le complément de résultat]
Le \cle{complément de résultat} (结果补语) est un mot placé \cle{immédiatement après le verbe} pour indiquer que l'action \textbf{aboutit} et quel est son \textbf{résultat}. Le français rend souvent cette idée en changeant de verbe : \zh{看} kàn « regarder » (action) → \zh{看见} kànjian « voir » (résultat) ; \zh{听} tīng « écouter » → \zh{听懂} tīngdǒng « comprendre en écoutant ».
\end{definition}

\begin{center}
\begin{tabularx}{\linewidth}{|X|X|X|X|}
\hline
\rowcolor{popblueL} Caractères & Pinyin & Nature & Traduction \\
\hline
看见 & kànjian & verbe + résultat & voir (résultat de regarder) \\
\hline
听见 & tīngjian & verbe + résultat & entendre (résultat d'écouter) \\
\hline
听懂 & tīngdǒng & verbe + résultat & comprendre en écoutant \\
\hline
看懂 & kàndǒng & verbe + résultat & comprendre en lisant \\
\hline
做完 & zuòwán & verbe + résultat & finir de faire \\
\hline
写完 & xiěwán & verbe + résultat & finir d'écrire \\
\hline
买到 & mǎidào & verbe + résultat & réussir à acheter \\
\hline
找到 & zhǎodào & verbe + résultat & trouver (résultat de chercher) \\
\hline
说清楚 & shuō qīngchu & verbe + résultat & expliquer clairement \\
\hline
洗干净 & xǐ gānjìng & verbe + résultat & laver (jusqu'à ce que ce soit propre) \\
\hline
做对 & zuòduì & verbe + résultat & faire correctement \\
\hline
写错 & xiěcuò & verbe + résultat & mal écrire, se tromper en écrivant \\
\hline
\end{tabularx}
\end{center}

\courssub{Savoir-faire 1 : Construire le complément de résultat}

\begin{methode}[Construire une phrase avec un complément de résultat]
Le complément de résultat indique le \cle{résultat} ou l'\cle{aboutissement} de l'action : on ne dit pas seulement « écouter » mais « entendre » (résultat de l'écoute), pas seulement « regarder » mais « voir ».
\begin{enumerate}
\item Identifier le verbe d'action : \zh{看} kàn (regarder), \zh{听} tīng (écouter), \zh{买} mǎi (acheter), \zh{写} xiě (écrire), \zh{做} zuò (faire).
\item Choisir le complément de résultat adapté : \zh{见} jiàn (percevoir), \zh{到} dào (atteindre), \zh{完} wán (terminer), \zh{懂} dǒng (comprendre), \zh{对} duì (correct), \zh{错} cuò (faux), \zh{清楚} qīngchu (clair), \zh{干净} gānjìng (propre).
\item Coller le complément directement \cle{après le verbe}, sans mot entre les deux : \zh{看见}, \zh{听懂}, \zh{做完}.
\item Placer le complément d'objet \cle{après} le groupe verbe + résultat : \zh{我做完作业了。}
\item Ajouter \zh{了} en fin de phrase si l'action est accomplie.
\end{enumerate}
Schéma à retenir : \textbf{Sujet + Verbe + Résultat + (Objet) + 了}
\end{methode}

\begin{propriete}[Place des éléments dans la phrase]
\begin{enumerate}
\item Le complément de résultat forme avec le verbe un \cle{bloc inséparable} : rien ne peut s'intercaler entre eux (sauf \zh{得} ou \zh{不} dans la forme potentielle, voir Savoir-faire 2).
\item L'objet suit toujours le bloc : \zh{我看完这本书了。} \emph{Wǒ kànwán zhè běn shū le.} « J'ai fini ce livre. »
\item \zh{了} se place après l'objet (ou juste après le bloc s'il n'y a pas d'objet) : \zh{他找到了工作。} \emph{Tā zhǎodào le gōngzuò.} « Il a trouvé un travail. »
\end{enumerate}
\end{propriete}

\begin{exoresolu}[Application : compléter avec le bon complément de résultat]
\textbf{Énoncé.} Complétez les phrases avec le complément de résultat qui convient (\zh{见}, \zh{懂}, \zh{完}, \zh{到}, \zh{清楚}), puis traduisez.
\begin{enumerate}
\item \zh{我没听\_\_\_\_老师说的话。}
\item \zh{作业太多，我晚上十点才做\_\_\_\_。}
\item \zh{在杜阿拉机场，我看\_\_\_\_了我的朋友。}
\item \zh{请你说\_\_\_\_，你的名字怎么写？}
\item \zh{我在市场上买\_\_\_\_了新鲜的水果。}
\end{enumerate}
\tcblower
\textbf{Solution détaillée.}
\begin{enumerate}
\item \zh{没听懂} méi tīngdǒng : écouter sans comprendre le contenu appelle \zh{懂} (comprendre). Phrase : \zh{我没听懂老师说的话。} \emph{Wǒ méi tīngdǒng lǎoshī shuō de huà.} « Je n'ai pas compris ce que le professeur a dit. »
\item \zh{做完} zuòwán : l'idée est la fin du travail, donc \zh{完} (terminer). \zh{作业太多，我晚上十点才做完。} \emph{Zuòyè tài duō, wǒ wǎnshang shí diǎn cái zuòwán.} « Il y avait trop de devoirs, je ne les ai terminés qu'à dix heures du soir. »
\item \zh{看见} kànjian : regarder avec le résultat « voir » exige \zh{见}. \zh{在杜阿拉机场，我看见了我的朋友。} \emph{Zài Dù'ālā jīchǎng, wǒ kànjian le wǒ de péngyou.} « À l'aéroport de Douala, j'ai vu mon ami. »
\item \zh{说清楚} shuō qīngchu : parler de façon claire, résultat \zh{清楚}. \zh{请你说清楚，你的名字怎么写？} \emph{Qǐng nǐ shuō qīngchu, nǐ de míngzi zěnme xiě?} « Parle clairement, s'il te plaît : comment s'écrit ton nom ? »
\item \zh{买到} mǎidào : l'achat aboutit, on obtient l'objet, donc \zh{到}. \zh{我在市场上买到了新鲜的水果。} \emph{Wǒ zài shìchǎng shang mǎidào le xīnxiān de shuǐguǒ.} « J'ai réussi à acheter des fruits frais au marché. »
\end{enumerate}
\textbf{Conclusion.} Le complément se choisit selon le sens du résultat et se colle toujours au verbe ; l'objet vient après le bloc verbe + résultat.
\end{exoresolu}

\courssub{Savoir-faire 2 : Nier, interroger et utiliser la forme potentielle}

\begin{methode}[Négation, question et potentiel du complément de résultat]
\begin{enumerate}
\item \textbf{Négation} : on n'utilise jamais \zh{不} seul devant le verbe ; la négation d'une action non aboutie se fait avec \zh{没（有）} devant le groupe verbe + résultat, et \zh{了} disparaît : \zh{我没看见他。} \emph{Wǒ méi kànjian tā.} « Je ne l'ai pas vu. »
\item \textbf{Question} : question avec \zh{吗} en gardant \zh{了} : \zh{你听懂了吗？} \emph{Nǐ tīngdǒng le ma?} « As-tu compris ? » ; ou forme affirmative-négative : \zh{你看见了没有？} \emph{Nǐ kànjian le méiyǒu?} « L'as-tu vu ou non ? »
\item \textbf{Forme potentielle affirmative} : insérer \zh{得} (de, ton neutre) entre le verbe et le résultat : \zh{听得懂} \emph{tīng de dǒng} « capable de comprendre en écoutant ».
\item \textbf{Forme potentielle négative} : insérer \zh{不} : \zh{听不懂} \emph{tīng bu dǒng} « incapable de comprendre en écoutant ».
\item Attention : la forme potentielle ne prend \cle{jamais} \zh{了} ; elle exprime la capacité, pas un fait accompli.
\end{enumerate}
Schémas : \textbf{没 + V + Résultat} (négation) ; \textbf{V + 得 + Résultat} (potentiel +) ; \textbf{V + 不 + Résultat} (potentiel −).
\end{methode}

\begin{propriete}[Résumé des quatre formes]
\begin{center}
\begin{tabularx}{\linewidth}{|X|X|X|}
\hline
\rowcolor{popblueL} Forme & Exemple en caractères + pinyin & Traduction \\
\hline
Affirmative (accompli) & \zh{我听懂了。} \emph{Wǒ tīngdǒng le.} & J'ai compris (en écoutant). \\
\hline
Négative & \zh{我没听懂。} \emph{Wǒ méi tīngdǒng.} & Je n'ai pas compris. \\
\hline
Potentielle + & \zh{我听得懂。} \emph{Wǒ tīng de dǒng.} & Je peux comprendre (c'est à ma portée). \\
\hline
Potentielle − & \zh{我听不懂。} \emph{Wǒ tīng bu dǒng.} & Je ne peux pas comprendre (trop difficile). \\
\hline
\end{tabularx}
\end{center}
Nuance essentielle : \zh{没听懂} constate un fait (« je n'ai pas compris cette fois ») ; \zh{听不懂} exprime une incapacité (« c'est au-dessus de mon niveau »).
\end{propriete}

\begin{exoresolu}[Application : transformer les phrases]
\textbf{Énoncé.} (a) Mettez à la forme négative : \zh{我看完了这本书。} (b) Transformez en question avec \zh{吗} : \zh{他听懂了老师的问题。} (c) Exprimez la capacité : « Ce texte est facile, je peux le comprendre en le lisant » avec \zh{看得懂}.
\tcblower
\textbf{Solution détaillée.}
\begin{enumerate}
\item[(a)] La négation d'un résultat accompli exige \zh{没} et la suppression de \zh{了} : \zh{我没看完这本书。} \emph{Wǒ méi kànwán zhè běn shū.} « Je n'ai pas fini ce livre. » Erreur à éviter : \zh{我不看完了} est impossible.
\item[(b)] On garde \zh{了} (fait accompli) et on ajoute \zh{吗} : \zh{他听懂老师的问题了吗？} \emph{Tā tīngdǒng lǎoshī de wèntí le ma?} « A-t-il compris la question du professeur ? »
\item[(c)] Potentiel affirmatif avec \zh{得} : \zh{这篇课文很容易，我看得懂。} \emph{Zhè piān kèwén hěn róngyì, wǒ kàn de dǒng.} « Ce texte est très facile, je peux le comprendre en le lisant. » Pas de \zh{了} car il s'agit d'une capacité, non d'un fait.
\end{enumerate}
\textbf{Conclusion.} Trois réflexes : négation avec \zh{没} sans \zh{了} ; question avec \zh{吗} en gardant \zh{了} ; potentiel avec \zh{得}/\zh{不} inséré, jamais de \zh{了}.
\end{exoresolu}

\courssub{Savoir-faire 3 : Exprimer les fractions avec 分之}

\begin{methode}[Lire et écrire une fraction avec 分之]
En chinois, la fraction se lit \cle{à l'envers} du français : on dit d'abord le \textbf{dénominateur}, puis \zh{分之}, puis le \textbf{numérateur}.
\begin{enumerate}
\item Identifier le dénominateur (le nombre de parts) : \zh{三} sān (3), \zh{四} sì (4), \zh{五} wǔ (5).
\item Ajouter \zh{分之} fēn zhī (littéralement « parmi les parts »).
\item Ajouter le numérateur : \zh{一} yī, \zh{二} èr, \zh{三} sān.
\item Schéma : \textbf{dénominateur + 分之 + numérateur}. Ainsi « un tiers » = \zh{三分之一} \emph{sān fēn zhī yī} ; « deux cinquièmes » = \zh{五分之二} \emph{wǔ fēn zhī èr}.
\item Devant un nom, ajouter \zh{的} : \zh{三分之一的学生} \emph{sān fēn zhī yī de xuésheng} « un tiers des élèves ».
\end{enumerate}
Réflexe anti-calque : le français dit « un \textbf{sur} trois » (numérateur d'abord) ; le chinois dit « trois parts \textbf{parmi lesquelles} une » (dénominateur d'abord).
\end{methode}

\begin{propriete}[Correspondances français -- chinois]
\begin{center}
\begin{tabularx}{\linewidth}{|X|X|X|}
\hline
\rowcolor{popblueL} Français & Chinois + pinyin & Lecture littérale \\
\hline
un demi (1/2) & \zh{二分之一} \emph{èr fēn zhī yī} & sur deux parts, une \\
\hline
un tiers (1/3) & \zh{三分之一} \emph{sān fēn zhī yī} & sur trois parts, une \\
\hline
un quart (1/4) & \zh{四分之一} \emph{sì fēn zhī yī} & sur quatre parts, une \\
\hline
deux tiers (2/3) & \zh{三分之二} \emph{sān fēn zhī èr} & sur trois parts, deux \\
\hline
trois cinquièmes (3/5) & \zh{五分之三} \emph{wǔ fēn zhī sān} & sur cinq parts, trois \\
\hline
\end{tabularx}
\end{center}
Remarque : pour « un demi » dans la vie quotidienne, on utilise le plus souvent \zh{一半} \emph{yíbàn} « une moitié » ; \zh{二分之一} est la forme mathématique et écrite.
\end{propriete}

\begin{exoresolu}[Application : fractions dans un texte sur la santé]
\textbf{Énoncé.} Traduisez en chinois : (1) « Un quart des élèves de la classe fait du sport le matin. » (2) « Les deux tiers des élèves de notre classe aiment le football. » (3) Lisez et traduisez : \zh{我们班五分之三的同学住在雅温得。}
\tcblower
\textbf{Solution détaillée.}
\begin{enumerate}
\item « Un quart » : dénominateur 4 d'abord, numérateur 1 ensuite : \zh{四分之一}. Phrase : \zh{班上四分之一的学生早上做运动。} \emph{Bān shang sì fēn zhī yī de xuésheng zǎoshang zuò yùndòng.} « Un quart des élèves de la classe fait du sport le matin. »
\item « Les deux tiers » : \zh{三分之二} (3 d'abord, 2 ensuite). Phrase : \zh{我们班三分之二的学生喜欢足球。} \emph{Wǒmen bān sān fēn zhī èr de xuésheng xǐhuan zúqiú.} « Les deux tiers des élèves de notre classe aiment le football. »
\item \zh{五分之三} = 3/5 (dénominateur 5, numérateur 3). Traduction : « Trois cinquièmes des camarades de notre classe habitent à Yaoundé. » \emph{Wǒmen bān wǔ fēn zhī sān de tóngxué zhù zài Yǎwēndé.}
\end{enumerate}
\textbf{Conclusion.} Toujours vérifier l'ordre : le premier nombre chinois est le dénominateur. Inverser donnerait 1/4 au lieu de 3/4, une erreur grave de sens.
\end{exoresolu}

\courssub{Savoir-faire 4 : Exprimer les pourcentages avec 百分之}

\begin{methode}[Lire et écrire un pourcentage avec 百分之]
Un pourcentage est une fraction sur cent : on applique le même schéma avec \zh{百} bǎi (100) comme dénominateur.
\begin{enumerate}
\item Écrire \zh{百分之} bǎi fēn zhī (« sur cent parts »).
\item Ajouter le chiffre du pourcentage : \zh{百分之二十} \emph{bǎi fēn zhī èrshí} = 20 \%.
\item Pour les nombres composés, suivre la numération normale : \zh{百分之七十五} \emph{bǎi fēn zhī qīshíwǔ} = 75 \%.
\item Devant un nom, utiliser \zh{的} : \zh{百分之九十的学生} \emph{bǎi fēn zhī jiǔshí de xuésheng} « 90 \% des élèves ».
\item Pour « 100 \% », on dit \zh{百分之百} \emph{bǎi fēn zhī bǎi} ; dans la vie quotidienne, on préfère souvent l'adverbe \zh{都} : \zh{学生都来了。} \emph{Xuésheng dōu lái le.} « Tous les élèves sont venus. »
\end{enumerate}
Schéma : \textbf{百分之 + nombre (+ 的 + nom)}. Le symbole \% ne s'écrit pas dans une phrase chinoise rédigée : on écrit le nombre en toutes lettres.
\end{methode}

\begin{exoresolu}[Application : pourcentages dans un mini-texte sur l'environnement]
\textbf{Énoncé.} Lisez le texte, puis répondez en français et complétez la dernière phrase.\\
\zh{在我们学校，百分之八十的学生每天喝水，百分之三十的学生早上跑步，只有百分之五的学生抽烟。}\\
\emph{Zài wǒmen xuéxiào, bǎi fēn zhī bāshí de xuésheng měitiān hē shuǐ, bǎi fēn zhī sānshí de xuésheng zǎoshang pǎobù, zhǐyǒu bǎi fēn zhī wǔ de xuésheng chōuyān.}\\
Questions : (1) Quelle proportion d'élèves boit de l'eau chaque jour ? (2) Quelle proportion d'élèves court le matin ? (3) Traduisez la dernière partie. (4) Complétez : « 60 \% des élèves aiment le ndolé » : \zh{百分之\_\_\_\_的学生喜欢恩多莱。}
\tcblower
\textbf{Solution détaillée.}
\begin{enumerate}
\item \zh{百分之八十} = 80 \% (dénominateur 100, numérateur 80). Réponse : 80 \% des élèves boivent de l'eau chaque jour.
\item \zh{百分之三十} = 30 \%. Réponse : 30 \% des élèves courent le matin.
\item \zh{只有百分之五的学生抽烟。} \emph{Zhǐyǒu bǎi fēn zhī wǔ de xuésheng chōuyān.} « Seulement 5 \% des élèves fument. » Noter \zh{只有} zhǐyǒu « seulement », qui souligne la faiblesse du chiffre.
\item 60 = \zh{六十} liùshí : \zh{百分之六十的学生喜欢恩多莱。} \emph{Bǎi fēn zhī liùshí de xuésheng xǐhuan Ēnduōlái.} « 60 \% des élèves aiment le ndolé. »
\end{enumerate}
\textbf{Conclusion.} Le pourcentage n'est qu'un cas particulier de la fraction : maîtriser \zh{分之} suffit pour lire tous les chiffres d'un texte d'environnement ou de santé.
\end{exoresolu}

\courssub{Savoir-faire 5 : Rédiger un court texte combinant résultat et fractions}

\begin{methode}[Produire un paragraphe avec complément de résultat et fractions]
\begin{enumerate}
\item Choisir un sujet concret (santé, sport, environnement à Yaoundé ou Douala).
\item Construire d'abord les blocs verbe + résultat sur brouillon : \zh{看见}, \zh{做完}, \zh{听懂}, \zh{洗干净} xǐ gānjìng (laver propre).
\item Ajouter une ou deux données chiffrées avec \zh{分之} ou \zh{百分之} pour étayer le propos.
\item Vérifier l'ordre : Sujet + V + Résultat + Objet + 了 ; dénominateur + 分之 + numérateur.
\item Relire en traquant : \zh{了} interdit après un potentiel, \zh{没} (et non \zh{不}) pour nier un résultat accompli, \zh{的} entre la fraction et le nom.
\end{enumerate}
\end{methode}

\begin{exoresolu}[Application : rédaction guidée]
\textbf{Énoncé.} Rédigez un paragraphe de 4 à 5 phrases en chinois sur le thème « Le sport et la santé dans mon lycée », en utilisant au moins deux compléments de résultat (\zh{做完}, \zh{看见}, \zh{听懂}…) et une fraction ou un pourcentage.
\tcblower
\textbf{Solution détaillée (modèle de rédaction).}\\
\zh{在我们学校，百分之七十的学生喜欢运动。早上，我看见很多同学在操场上跑步。老师说，运动对身体很好，我们都听懂了。昨天我做完了作业，就和朋友们去踢足球。踢完球以后，我觉得很舒服，也很健康。}\\
\emph{Zài wǒmen xuéxiào, bǎi fēn zhī qīshí de xuésheng xǐhuan yùndòng. Zǎoshang, wǒ kànjian hěn duō tóngxué zài cāochǎng shang pǎobù. Lǎoshī shuō, yùndòng duì shēntǐ hěn hǎo, wǒmen dōu tīngdǒng le. Zuótiān wǒ zuòwán le zuòyè, jiù hé péngyoumen qù tī zúqiú. Tī wán qiú yǐhòu, wǒ juéde hěn shūfu, yě hěn jiànkāng.}\\
« Dans notre lycée, 70 \% des élèves aiment le sport. Le matin, je vois beaucoup de camarades courir sur le terrain. Le professeur dit que le sport est bon pour la santé, et nous avons tous compris. Hier, j'ai fini mes devoirs, puis je suis allé jouer au football avec mes amis. Après avoir joué, je me suis senti bien et en bonne santé. »\\
\textbf{Justifications.} \zh{看见} : résultat de la perception ; \zh{听懂} : compréhension aboutie, d'où \zh{了} ; \zh{做完} et \zh{踢完} : actions menées à leur terme ; \zh{百分之七十} respecte l'ordre dénominateur + numérateur, suivi de \zh{的} devant \zh{学生}.
\end{exoresolu}

\begin{attention}[Les 5 erreurs qui coûtent des points à l'examen]
\begin{enumerate}
\item \textbf{Inverser la fraction} : écrire \zh{一分之三} pour « un tiers » est faux ; le dénominateur vient \cle{toujours} en premier : \zh{三分之一}. Calque direct du français « un sur trois ».
\item \textbf{Séparer le verbe de son complément de résultat} : on ne peut pas dire \zh{看作业完} ; le résultat se colle au verbe et l'objet suit : \zh{做完作业}. Jamais de mot entre verbe et résultat.
\item \textbf{Nier avec 不 au lieu de 没} : \zh{我不看见他} est incorrect pour un fait passé ; la négation du résultat accompli exige \zh{没（有）} et la chute de \zh{了} : \zh{我没看见他。}
\item \textbf{Ajouter 了 après une forme potentielle} : \zh{我看得懂了} est fautif ; le potentiel (\zh{看得懂} / \zh{看不懂}) exprime une capacité et refuse \zh{了}.
\item \textbf{Oublier 的 entre la fraction et le nom, ou écrire le symbole \%} : il faut \zh{百分之二十的学生} (avec \zh{的}) et non un collage direct ; dans une phrase rédigée, on écrit le nombre en caractères, jamais « 20\% ». Attention aussi aux tons : \zh{分之} fēn zhī (1\textsuperscript{er} ton + 1\textsuperscript{er} ton), \zh{百} bǎi (3\textsuperscript{e} ton).
\end{enumerate}
\end{attention}

\begin{aretenir}[Bilan de la leçon]
\begin{enumerate}
\item Le complément de résultat se colle au verbe : \textbf{Sujet + V + Résultat + (Objet) + 了}. Exemple : \zh{我做完作业了。}
\item Négation : \zh{没} + V + Résultat, sans \zh{了} : \zh{我没做完作业。}
\item Potentiel : V + \zh{得}/\zh{不} + Résultat, jamais de \zh{了} : \zh{我做得完。} / \zh{我做不完。}
\item Fraction : \textbf{dénominateur + 分之 + numérateur} : \zh{三分之一} « un tiers ».
\item Pourcentage : \zh{百分之} + nombre : \zh{百分之八十} « 80 \% » ; \zh{的} devant le nom.
\end{enumerate}
\end{aretenir}

\begin{exercice}[Pour s'entraîner]
\begin{enumerate}
\item Traduisez en chinois : (a) « J'ai fini mes devoirs. » (b) « Je n'ai pas vu le professeur. » (c) « Ce caractère est trop difficile, je ne peux pas l'écrire correctement. » (utiliser \zh{写对})
\item Mettez à la forme négative puis à la forme potentielle négative : \zh{他找到了他的书。}
\item Écrivez en chinois : 1/2 (forme mathématique), 3/4, 2/5, 45 \%, 99 \%.
\item Traduisez en français : \zh{在杜阿拉，百分之六十的学生每天走路上学，四分之一的学生坐出租车。}
\item Rédigez 4 phrases sur « La santé dans mon quartier » avec un complément de résultat et un pourcentage.
\end{enumerate}
\end{exercice}

\begin{corrige}[Exercice « Pour s'entraîner »]
\begin{enumerate}
\item (a) \zh{我做完作业了。} \emph{Wǒ zuòwán zuòyè le.} (b) \zh{我没看见老师。} \emph{Wǒ méi kànjian lǎoshī.} (c) \zh{这个字太难，我写不对。} \emph{Zhè ge zì tài nán, wǒ xiě bu duì.}
\item Négation : \zh{他没找到他的书。} \emph{Tā méi zhǎodào tā de shū.} « Il n'a pas trouvé son livre. » Potentiel négatif : \zh{他找不到他的书。} \emph{Tā zhǎo bu dào tā de shū.} « Il n'arrive pas à trouver son livre. »
\item \zh{二分之一} èr fēn zhī yī ; \zh{四分之三} sì fēn zhī sān ; \zh{五分之二} wǔ fēn zhī èr ; \zh{百分之四十五} bǎi fēn zhī sìshíwǔ ; \zh{百分之九十九} bǎi fēn zhī jiǔshíjiǔ.
\item « À Douala, 60 \% des élèves vont à l'école à pied chaque jour, un quart des élèves prend le taxi. »
\item Modèle : \zh{在我们区，百分之五十的人早上做运动。我看见很多老人在公园里走路。医生说的话，我都听懂了：多喝水，少吃糖。昨天我洗完了衣服，觉得身体很舒服。} \emph{Zài wǒmen qū, bǎi fēn zhī wǔshí de rén zǎoshang zuò yùndòng. Wǒ kànjian hěn duō lǎorén zài gōngyuán li zǒulù. Yīshēng shuō de huà, wǒ dōu tīngdǒng le : duō hē shuǐ, shǎo chī táng. Zuótiān wǒ xǐwán le yīfu, juéde shēntǐ hěn shūfu.}
\end{enumerate}
\end{corrige}

\newpage

\coursec{Exercices}

\courssub{Je vérifie mes connaissances}

\begin{exercice}[QCM : le bon complément de résultat]
Pour chaque phrase, choisis la bonne réponse (A, B ou C).
\begin{enumerate}
\item 我听\underline{\hspace{1cm}}了老师的话，现在明白了。\quad A. 完\quad B. 懂\quad C. 错
\item 我的书找不\underline{\hspace{1cm}}了，可能在教室里。\quad A. 见\quad B. 完\quad C. 到
\item 请你写\underline{\hspace{1cm}}一点儿，我看不清楚。\quad A. 清楚\quad B. 懂\quad C. 见
\item 昨天的考试很难，我做\underline{\hspace{1cm}}了两道题。\quad A. 到\quad B. 错\quad C. 完
\end{enumerate}
\end{exercice}

\begin{exercice}[Vrai ou faux ? Justifie ta réponse]
Dis si chaque phrase est correcte (V) ou fautive (F), puis justifie en français en citant la règle de grammaire.
\begin{enumerate}
\item 我做完了作业。(Wǒ zuòwán le zuòyè.)
\item 我不看见他。(Wǒ bù kànjiàn tā.)
\item 我们班五分之三的学生是女生。(Wǒmen bān wǔ fēn zhī sān de xuésheng shì nǚshēng.)
\item 七十百分之的水很干净。(Qīshí bǎi fēn zhī de shuǐ hěn gānjìng.)
\end{enumerate}
\end{exercice}

\begin{exercice}[Vocabulaire : compléter le tableau]
Complète le tableau en écrivant les éléments manquants (caractères, pinyin ou traduction).
\begin{center}
\begin{tabularx}{\linewidth}{|X|X|X|X|}\hline
\rowcolor{popblueL} Caractères & Pinyin & Nature & Traduction \\ \hline
完成 & & verbe & \\ \hline
 & tīngjiàn & verbe & \\ \hline
看见 & & verbe & \\ \hline
 & qīngchu & adjectif & \\ \hline
百分之 & & expression & \\ \hline
\end{tabularx}
\end{center}
\end{exercice}

\begin{exercice}[Associer pinyin et caractères]
Associe chaque pinyin (colonne A) au caractère ou mot qui convient (colonne B).
\begin{center}
\begin{tabularx}{\linewidth}{|X|X|}\hline
\rowcolor{popblueL} Colonne A (pinyin) & Colonne B (caractères) \\ \hline
1. zuòwán & a. 听懂 \\ \hline
2. tīngdǒng & b. 做完 \\ \hline
3. zhǎodào & c. 看错 \\ \hline
4. kàncuò & d. 找到 \\ \hline
5. shuōqīngchu & e. 说清楚 \\ \hline
\end{tabularx}
\end{center}
\end{exercice}

\courssub{Je m'entraîne}

\begin{exercice}[Traduire en français]
Traduis les phrases suivantes en français.
\begin{enumerate}
\item 我看懂了这篇文章。(Wǒ kàndǒng le zhè piān wénzhāng.)
\item 他没做完作业，所以不能去玩。(Tā méi zuòwán zuòyè, suǒyǐ bù néng qù wán.)
\item 我们班三分之二的学生喜欢足球。(Wǒmen bān sān fēn zhī èr de xuésheng xǐhuan zúqiú.)
\item 百分之九十的病人已经看过医生了。(Bǎi fēn zhī jiǔshí de bìngrén yǐjīng kànguò yīshēng le.)
\end{enumerate}
\end{exercice}

\begin{exercice}[Transformer : négation et forme potentielle]
Pour chaque phrase, donne : (a) la forme négative avec 没(有) ; (b) la forme potentielle affirmative avec 得 ; (c) la forme potentielle négative avec 不.
\begin{enumerate}
\item 我做完了作业。(Wǒ zuòwán le zuòyè.)
\item 他听懂了老师的话。(Tā tīngdǒng le lǎoshī de huà.)
\item 我们找到了医院。(Wǒmen zhǎodào le yīyuàn.)
\end{enumerate}
\end{exercice}

\begin{exercice}[Texte à trous : compléter avec le bon complément]
Complète le texte « À la clinique de Yaoundé » ci-dessous avec le complément de résultat qui convient parmi : 完、到、见、懂、错、清楚 (chaque complément est utilisé une seule fois).
\end{exercice}

\begin{textebox}[À la clinique de Yaoundé]
昨天我不舒服，妈妈带我去医院。我们很快就找\underline{\hspace{1cm}}了医院。医生说得很慢，我都听\underline{\hspace{1cm}}了。他让我吃\underline{\hspace{1cm}}药，可是我还没看\underline{\hspace{1cm}}说明书。妈妈对我说：「请你说\underline{\hspace{1cm}}一点儿，哪里不舒服？」我说\underline{\hspace{1cm}}了，不是头疼，是肚子疼。\\[4pt]
Zuótiān wǒ bù shūfu, māma dài wǒ qù yīyuàn. Wǒmen hěn kuài jiù zhǎo\underline{\hspace{0.8cm}} le yīyuàn. Yīshēng shuō de hěn màn, wǒ dōu tīng\underline{\hspace{0.8cm}} le. Tā ràng wǒ chī\underline{\hspace{0.8cm}} yào, kěshì wǒ hái méi kàn\underline{\hspace{0.8cm}} shuōmíngshū. Māma duì wǒ shuō : « Qǐng nǐ shuō\underline{\hspace{0.8cm}} yìdiǎnr, nǎli bù shūfu ? » Wǒ shuō\underline{\hspace{0.8cm}} le, bù shì tóu téng, shì dùzi téng.\\[4pt]
\emph{Traduction : Hier je ne me sentais pas bien, maman m'a emmené à l'hôpital. Nous avons vite trouvé l'hôpital. Le médecin parlait lentement, j'ai tout compris. Il m'a demandé de finir de prendre le médicament, mais je n'avais pas encore fini de lire la notice. Maman m'a dit : « Parle un peu plus clairement, où as-tu mal ? » Je m'étais trompé : ce n'était pas la tête, c'était le ventre.}
\end{textebox}

\begin{exercice}[Remettre les mots dans l'ordre]
Remets les mots dans le bon ordre pour former des phrases correctes, puis traduis-les en français.
\begin{enumerate}
\item 学生 / 五分之三 / 的 / 我们班 / 是 / 女生
\item 百分之 / 干净 / 水 / 八十 / 的 / 很
\item 作业 / 做完了 / 我 / 今天的
\item 听得懂 / 你 / 吗 / 老师的话
\end{enumerate}
\end{exercice}

\courssub{Je me prépare au Bac}

\begin{textebox}[L'eau potable dans notre ville]
在我们城市，百分之六十的家庭有干净的水。五分之二的家庭还没有干净的水，所以很多人常常生病。上个星期，医生来我们学校讲课。他说得很清楚，我们都听懂了。他告诉我们：喝水以前，一定要看清楚水干净不干净。听完课以后，我们班四分之三的学生说：「我们找对了方法！」现在，百分之九十的学生喝干净的水，生病的人越来越少了。\\[4pt]
Zài wǒmen chéngshì, bǎi fēn zhī liùshí de jiātíng yǒu gānjìng de shuǐ. Wǔ fēn zhī èr de jiātíng hái méiyǒu gānjìng de shuǐ, suǒyǐ hěn duō rén chángcháng shēngbìng. Shàng ge xīngqī, yīshēng lái wǒmen xuéxiào jiǎngkè. Tā shuō de hěn qīngchu, wǒmen dōu tīngdǒng le. Tā gàosu wǒmen : hē shuǐ yǐqián, yídìng yào kànqīngchu shuǐ gānjìng bù gānjìng. Tīngwán kè yǐhòu, wǒmen bān sì fēn zhī sān de xuésheng shuō : « Wǒmen zhǎoduì le fāngfǎ ! » Xiànzài, bǎi fēn zhī jiǔshí de xuésheng hē gānjìng de shuǐ, shēngbìng de rén yuè lái yuè shǎo le.\\[4pt]
\emph{Traduction : Dans notre ville, 60 \% des familles ont de l'eau propre. Deux cinquièmes des familles n'ont pas encore d'eau propre, donc beaucoup de gens tombent souvent malades. La semaine dernière, un médecin est venu faire un exposé dans notre école. Il parlait clairement et nous avons tout compris. Il nous a dit : avant de boire l'eau, il faut bien vérifier si elle est propre ou non. Après le cours, les trois quarts des élèves de notre classe ont dit : « Nous avons trouvé la bonne méthode ! » Maintenant, 90 \% des élèves boivent de l'eau propre et les malades sont de moins en moins nombreux.}
\end{textebox}

\begin{exercice}[Texte de compréhension : la santé au Cameroun]
Lis le texte « L'eau potable dans notre ville » ci-dessus, puis réponds aux questions.

\textbf{Questions de compréhension} (réponds en français)
\begin{enumerate}
\item Quelle proportion des familles de la ville a de l'eau potable ? Écris cette fraction en chinois.
\item Pourquoi beaucoup de personnes tombent-elles malades ?
\item Que conseille le médecin avant de boire l'eau ?
\end{enumerate}

\textbf{Questions de langue}
\begin{enumerate}
\item Releve dans le texte deux compléments de résultat différents et explique leur construction (verbe + complément).
\item Mets à la forme négative : 我们都听懂了。
\item Écris en chinois la fraction « un quart » et le pourcentage « 75 \% ».
\end{enumerate}
\end{exercice}

\begin{exercice}[Thème et version]
\textbf{A. Thème.} Traduis en chinois (caractères + pinyin) :
\begin{enumerate}
\item Les trois quarts des élèves ont fini leurs devoirs.
\item Je n'ai pas compris, je n'arrive pas à voir le tableau.
\item 65 \% des malades ont vu le médecin hier.
\end{enumerate}

\textbf{B. Version.} Traduis en français :
\begin{enumerate}
\item 我们班五分之四的学生做完了练习。(Wǒmen bān wǔ fēn zhī sì de xuésheng zuòwán le liànxí.)
\item 他看错了字，所以写错了。(Tā kàncuò le zì, suǒyǐ xiěcuò le.)
\item 百分之百的医生都说这个药很好。(Bǎi fēn zhī bǎi de yīshēng dōu shuō zhège yào hěn hǎo.)
\end{enumerate}
\end{exercice}

\begin{exercice}[Expression écrite guidée : la santé à Douala]
Voici un tableau de statistiques sanitaires (fictives) sur un quartier de Douala :
\begin{center}
\begin{tabularx}{\linewidth}{|X|X|}\hline
\rowcolor{popblueL} Situation & Proportion \\ \hline
Familles ayant de l'eau potable & 70 \% \\ \hline
Enfants vaccinés & 五分之四 \\ \hline
Malades ayant vu un médecin & 60 \% \\ \hline
Élèves ayant compris l'exposé du médecin & 四分之三 \\ \hline
Habitants ayant trouvé la clinique & 90 \% \\ \hline
\end{tabularx}
\end{center}
À partir de ce tableau, rédige un court texte de \textbf{5 à 8 phrases} (au moins 60 caractères chinois) sur la santé dans ce quartier. Tu dois utiliser :
\begin{itemize}
\item au moins \textbf{deux pourcentages} avec 百分之 ;
\item au moins \textbf{une fraction} avec 分之 ;
\item au moins \textbf{deux compléments de résultat} différents (par exemple 完、到、懂、见、错、清楚).
\end{itemize}
\end{exercice}

\begin{methode}[Réussir ta production écrite]
\begin{enumerate}
\item \textbf{Repère} dans le tableau les données que tu vas utiliser et traduis mentalement chaque proportion en chinois : 70 \% = 百分之七十 ; 五分之四 = 4/5.
\item \textbf{Planifie} : phrase d'introduction (où ? quoi ?), puis une phrase par donnée, puis une conclusion (avis ou conseil).
\item \textbf{Rédige} en respectant les structures : Sujet + verbe + complément de résultat + 了 ; dénominateur + 分之 + numérateur + 的 + nom ; 百分之 + nombre + 的 + nom.
\item \textbf{Vérifie} : ordre des mots, présence de 了 après le complément de résultat à l'affirmatif, négation avec 没(有), place de 的 avant le nom.
\end{enumerate}
\end{methode}

\begin{exemplebox}[Modèle de production]
在杜阿拉这个区，百分之七十的家庭有干净的水。五分之四的孩子打了疫苗，他们很健康。百分之六十的病人看了医生，医生说得清楚，他们都听懂了。四分之三的学生听完了医生的课。百分之九十的人找到了诊所。现在，这个区的人生病越来越少了，大家都很高兴！\\[4pt]
Zài Dù'ālā zhège qū, bǎi fēn zhī qīshí de jiātíng yǒu gānjìng de shuǐ. Wǔ fēn zhī sì de háizi dǎ le yìmiáo, tāmen hěn jiànkāng. Bǎi fēn zhī liùshí de bìngrén kàn le yīshēng, yīshēng shuō de qīngchu, tāmen dōu tīngdǒng le. Sì fēn zhī sān de xuésheng tīngwán le yīshēng de kè. Bǎi fēn zhī jiǔshí de rén zhǎodào le zhěnsuǒ. Xiànzài, zhège qū de rén shēngbìng yuè lái yuè shǎo le, dàjiā dōu hěn gāoxìng !\\[4pt]
\emph{Traduction : Dans ce quartier de Douala, 70 \% des familles ont de l'eau propre. Quatre cinquièmes des enfants sont vaccinés et sont en bonne santé. 60 \% des malades ont vu un médecin ; le médecin parlait clairement et ils ont tous compris. Les trois quarts des élèves ont écouté l'exposé du médecin jusqu'au bout. 90 \% des habitants ont trouvé la clinique. Maintenant, les habitants de ce quartier tombent de moins en moins malades, et tout le monde est content !}
\end{exemplebox}

\begin{attention}[Erreurs fréquentes des francophones]
\begin{itemize}
\item \textbf{Ordre du pourcentage} : le français dit « soixante-dix pour cent », le chinois dit 百分之七十 (litt. « cent parts — soixante »). Ne dis jamais 七十百分之.
\item \textbf{Ordre de la fraction} : en chinois, le \textbf{dénominateur} vient d'abord : 三分之二 = « sur trois parts, deux » = 2/3. Ne traduis pas mot à mot « deux tiers ».
\item \textbf{Négation du complément de résultat} : au passé, utilise 没(有) sans 了 : 我没做完。 et non 我不做完了。
\item \textbf{Forme potentielle} : 得 / 不 se placent \textbf{entre} le verbe et le complément, sans 了 : 听得懂、看不见、做不完。
\end{itemize}
\end{attention}

\newpage

\coursec{Corrigés des exercices}

\begin{corrige}[Exercice 1 — QCM : le bon complément de résultat]
\begin{enumerate}
\item \textbf{B. 懂} : 我听懂了老师的话，现在明白了。(Wǒ tīngdǒng le lǎoshī de huà, xiànzài míngbai le.) « J'ai compris les paroles du professeur, maintenant je comprends. » 听懂 = écouter + comprendre : le complément 懂 indique que l'écoute a abouti à la compréhension.
\item \textbf{C. 到} : 我的书找不到了，可能在教室里。(Wǒ de shū zhǎo bù dào le, kěnéng zài jiàoshì lǐ.) « Je n'arrive pas à trouver mon livre, il est peut-être dans la salle de classe. » Forme potentielle négative : verbe + 不 + complément (找 + 不 + 到).
\item \textbf{A. 清楚} : 请你写清楚一点儿，我看不清楚。(Qǐng nǐ xiě qīngchu yìdiǎnr, wǒ kàn bù qīngchu.) « Écris un peu plus lisiblement, s'il te plaît, je n'arrive pas à lire clairement. » 清楚 « clair » s'accorde avec 看 « regarder » dans la phrase qui suit.
\item \textbf{C. 完} : 昨天的考试很难，我做完了两道题。(Zuótiān de kǎoshì hěn nán, wǒ zuòwán le liǎng dào tí.) « L'examen d'hier était difficile, j'ai terminé deux exercices. » 完 exprime l'accomplissement total de l'action ; 道 est le classificateur des questions et exercices.
\end{enumerate}
\textbf{Points de vigilance} : 完 (terminer), 到 (atteindre, obtenir), 懂 (comprendre), 错 (se tromper), 清楚 (clairement) ne sont pas interchangeables ; le choix dépend du sens du verbe et du résultat logique attendu.
\end{corrige}

\begin{corrige}[Exercice 2 — Vrai ou faux ? Justifie ta réponse]
\begin{enumerate}
\item \textbf{V.} Structure correcte : Sujet + verbe + complément de résultat (完) + 了 + complément d'objet. « J'ai fini mes devoirs. »
\item \textbf{F.} La négation d'une action achevée se fait avec 没(有), jamais avec 不 devant un verbe + complément de résultat au passé. Correction : 我没看见他。(Wǒ méi kànjiàn tā.) « Je ne l'ai pas vu. » Rappel : 不 ne s'emploie devant ce type de verbe que dans la forme potentielle insérée : 我看不见他。 « Je n'arrive pas à le voir. »
\item \textbf{V.} Fraction correcte : dénominateur + 分之 + numérateur : 五分之三 = 3/5. « Les trois cinquièmes des élèves de notre classe sont des filles. »
\item \textbf{F.} Ordre fautif : le pourcentage chinois se construit 百分之 + nombre (littéralement « sur cent parts, X »), contrairement au français « soixante-dix pour cent ». Correction : 百分之七十的水很干净。(Bǎi fēn zhī qīshí de shuǐ hěn gānjìng.) « 70 \% de l'eau est propre. »
\end{enumerate}
\end{corrige}

\begin{corrige}[Exercice 3 — Vocabulaire : compléter le tableau]
\begin{center}
\begin{tabularx}{\linewidth}{|X|X|X|X|}\hline
\rowcolor{popblueL} Caractères & Pinyin & Nature & Traduction \\ \hline
完成 & wánchéng & verbe & terminer, achever \\ \hline
听见 & tīngjiàn & verbe & entendre \\ \hline
看见 & kànjiàn & verbe & voir, apercevoir \\ \hline
清楚 & qīngchu & adjectif & clair, net \\ \hline
百分之 & bǎi fēn zhī & expression & pour cent (\%) \\ \hline
\end{tabularx}
\end{center}
\textbf{Points de vigilance} : attention aux tons de 见 (jiàn, 4\up{e} ton) et à la voyelle de 清楚 (qīngchu, le second caractère est atone). 听见 et 看见 partagent le même complément 见 « percevoir », l'un pour l'ouïe, l'autre pour la vue.
\end{corrige}

\begin{corrige}[Exercice 4 — Associer pinyin et caractères]
1 -- b (zuòwán 做完, finir de faire) ; 2 -- a (tīngdǒng 听懂, comprendre en écoutant) ; 3 -- d (zhǎodào 找到, trouver) ; 4 -- c (kàncuò 看错, mal lire, se tromper en regardant) ; 5 -- e (shuōqīngchu 说清楚, expliquer clairement).

\textbf{Points de vigilance} : ne pas confondre 到 (dào, atteindre) et 懂 (dǒng, comprendre) ; vérifier le ton de 错 (cuò, 4\up{e} ton).
\end{corrige}

\begin{corrige}[Exercice 5 — Traduire en français]
\begin{enumerate}
\item « J'ai compris cet article (en le lisant). » 看懂 = lire + comprendre ; 篇 (piān) est le classificateur des articles et des textes.
\item « Il n'a pas fini ses devoirs, donc il ne peut pas aller jouer. » Négation du complément de résultat : 没 + verbe + complément, sans 了.
\item « Les deux tiers des élèves de notre classe aiment le football. » 三分之二 = 2/3 (dénominateur d'abord).
\item « 90 \% des malades ont déjà vu le médecin. » 百分之九十 = 90 \% ; 过 (kànguò) marque l'expérience vécue.
\end{enumerate}
\end{corrige}

\begin{corrige}[Exercice 6 — Transformer : négation et forme potentielle]
\begin{enumerate}
\item (a) 我没(有)做完作业。(Wǒ méi zuòwán zuòyè.) (b) 我做得完作业。(Wǒ zuò de wán zuòyè.) (c) 我做不完作业。(Wǒ zuò bù wán zuòyè.) « Je n'arrive pas à finir mes devoirs. »
\item (a) 他没(有)听懂老师的话。(b) 他听得懂老师的话。(c) 他听不懂老师的话。(Tā tīng bu dǒng lǎoshī de huà.) « Il ne comprend pas ce que dit le professeur. »
\item (a) 我们没(有)找到医院。(b) 我们找得到医院。(c) 我们找不到医院。(Wǒmen zhǎo bu dào yīyuàn.) « Nous n'arrivons pas à trouver l'hôpital. »
\end{enumerate}
\textbf{Rappel de la règle} : la forme potentielle s'obtient en insérant 得 (capacité) ou 不 (incapacité) \textbf{entre} le verbe et le complément de résultat ; on ne met jamais 了 dans la forme potentielle. La négation simple du passé utilise 没(有) et supprime 了.
\end{corrige}

\begin{corrige}[Exercice 7 — Texte à trous : compléter avec le bon complément]
Réponses dans l'ordre : 找到 (trouver) ; 听懂 (comprendre en écoutant) ; 吃完 (finir de prendre le médicament) ; 看完 (finir de lire) ; 说清楚 (parler clairement) ; 说错 (se tromper en parlant).

Texte complet : 昨天我不舒服，妈妈带我去医院。我们很快就找到了医院。医生说得很慢，我都听懂了。他让我吃完药，可是我还没看完说明书。妈妈对我说：「请你说清楚一点儿，哪里不舒服？」我说错了，不是头疼，是肚子疼。

(Zuótiān wǒ bù shūfu, māma dài wǒ qù yīyuàn. Wǒmen hěn kuài jiù zhǎodào le yīyuàn. Yīshēng shuō de hěn màn, wǒ dōu tīngdǒng le. Tā ràng wǒ chīwán yào, kěshì wǒ hái méi kànwán shuōmíngshū. Māma duì wǒ shuō : « Qǐng nǐ shuō qīngchu yìdiǎnr, nǎli bù shūfu ? » Wǒ shuōcuò le, bù shì tóu téng, shì dùzi téng.)

« Hier je ne me sentais pas bien, maman m'a emmené à l'hôpital. Nous avons vite trouvé l'hôpital. Le médecin parlait lentement, j'ai tout compris. Il m'a demandé de finir de prendre le médicament, mais je n'avais pas encore fini de lire la notice. Maman m'a dit : "Parle un peu plus clairement, où as-tu mal ?" Je m'étais trompé : ce n'était pas la tête, c'était le ventre. »

\textbf{Points de vigilance} : chaque complément n'est utilisé qu'une fois ; vérifier la cohérence sémantique (on « finit » de lire une notice : 看完 ; on se « trompe » en parlant : 说错).
\end{corrige}

\begin{corrige}[Exercice 8 — Remettre les mots dans l'ordre]
\begin{enumerate}
\item 我们班五分之三的学生是女生。(Wǒmen bān wǔ fēn zhī sān de xuésheng shì nǚshēng.) « Les trois cinquièmes des élèves de notre classe sont des filles. » La fraction + 的 se place devant le nom 学生.
\item 百分之八十的水很干净。(Bǎi fēn zhī bāshí de shuǐ hěn gānjìng.) « 80 \% de l'eau est propre. » 百分之 + nombre + 的 + nom.
\item 我做完了今天的作业。(Wǒ zuòwán le jīntiān de zuòyè.) « J'ai fini les devoirs d'aujourd'hui. » 了 se place juste après le complément de résultat, avant le complément d'objet.
\item 你听得懂老师的话吗？(Nǐ tīng de dǒng lǎoshī de huà ma ?) « Comprends-tu ce que dit le professeur ? » Forme potentielle interrogative avec 吗 final.
\end{enumerate}
\end{corrige}

\begin{corrige}[Exercice 9 — Texte de compréhension : la santé au Cameroun]
\textbf{Questions de compréhension}
\begin{enumerate}
\item 60 \% des familles de la ville ont de l'eau potable ; en chinois : 百分之六十 (bǎi fēn zhī liùshí).
\item Parce que deux cinquièmes des familles n'ont pas encore d'eau propre (五分之二的家庭还没有干净的水), donc beaucoup de gens tombent souvent malades.
\item Le médecin conseille de bien vérifier si l'eau est propre avant de la boire (喝水以前，一定要看清楚水干净不干净).
\end{enumerate}
\textbf{Questions de langue}
\begin{enumerate}
\item Exemples relevés : 听懂 (tīng + dǒng : écouter + comprendre) ; 看清楚 (kàn + qīngchu : regarder + clair) ; 听完 (tīng + wán : finir d'écouter) ; 找对 (zhǎo + duì : trouver juste). Construction : verbe d'action + complément indiquant le résultat de l'action.
\item 我们都没(有)听懂。(Wǒmen dōu méi tīngdǒng.) « Nous n'avons pas compris. » (没 sans 了.)
\item Un quart : 四分之一 (sì fēn zhī yī) ; 75 \% : 百分之七十五 (bǎi fēn zhī qīshíwǔ).
\end{enumerate}
\textbf{Barème indicatif (sur 10)} : compréhension 4 pts (1 + 1,5 + 1,5) ; langue 6 pts (2 + 2 + 2). \textbf{Points de vigilance} : exiger la justification grammaticale complète (verbe + complément) et non la simple traduction ; sanctionner l'inversion de la fraction (四分之三 et non 三分之四).
\end{corrige}

\begin{corrige}[Exercice 10 — Thème et version]
\textbf{A. Thème}
\begin{enumerate}
\item 四分之三的学生做完了作业。(Sì fēn zhī sān de xuésheng zuòwán le zuòyè.)
\item 我没听懂，我看不见黑板。(Wǒ méi tīngdǒng, wǒ kàn bú jiàn hēibǎn.) Première proposition : négation au passé avec 没 ; seconde : forme potentielle négative 看不见 « ne pas arriver à voir ».
\item 百分之六十五的病人昨天看了医生。(Bǎi fēn zhī liùshíwǔ de bìngrén zuótiān kàn le yīshēng.)
\end{enumerate}
\textbf{B. Version}
\begin{enumerate}
\item « Les quatre cinquièmes des élèves de notre classe ont fini les exercices. »
\item « Il a mal lu le caractère, donc il l'a mal écrit. » (看错 : se tromper en lisant ; 写错 : se tromper en écrivant.)
\item « 100 \% des médecins disent que ce médicament est très bon. » (百分之百 = la totalité, 100 \%.)
\end{enumerate}
\textbf{Barème indicatif (sur 10)} : thème 5 pts, version 5 pts ; 0,5 pt retiré par faute de structure (ordre de la fraction, place de 了, négation). \textbf{Points de vigilance} : ne pas traduire « je n'arrive pas à voir » par 我不看见 (fautif) mais par la forme potentielle 我看不见 ; placer le complément de temps 昨天 avant le verbe.
\end{corrige}

\begin{corrige}[Exercice 11 — Expression écrite guidée : la santé à Douala]
\textbf{Barème indicatif (sur 20)} : respect des contraintes (2 pourcentages, 1 fraction, 2 compléments de résultat différents) : 6 pts ; correction grammaticale (ordre des mots, 了, 没, 的) : 6 pts ; richesse du vocabulaire et orthographe des caractères : 4 pts ; cohérence et organisation du texte (introduction, développement, conclusion) : 4 pts.

\textbf{Proposition de rédaction modèle} :

在杜阿拉这个区，百分之七十的家庭有干净的水。五分之四的孩子打了疫苗，他们很健康。百分之六十的病人看了医生，医生说得很清楚，他们都听懂了。四分之三的学生听完了医生的课。百分之九十的人找到了诊所。现在，这个区的人生病越来越少了，大家都很高兴！

(Zài Dù'ālā zhège qū, bǎi fēn zhī qīshí de jiātíng yǒu gānjìng de shuǐ. Wǔ fēn zhī sì de háizi dǎ le yìmiáo, tāmen hěn jiànkāng. Bǎi fēn zhī liùshí de bìngrén kàn le yīshēng, yīshēng shuō de hěn qīngchu, tāmen dōu tīngdǒng le. Sì fēn zhī sān de xuésheng tīngwán le yīshēng de kè. Bǎi fēn zhī jiǔshí de rén zhǎodào le zhěnsuǒ. Xiànzài, zhège qū de rén shēngbìng yuè lái yuè shǎo le, dàjiā dōu hěn gāoxìng !)

\emph{Traduction : Dans ce quartier de Douala, 70 \% des familles ont de l'eau propre. Quatre cinquièmes des enfants sont vaccinés et sont en bonne santé. 60 \% des malades ont vu un médecin ; le médecin parlait clairement et ils ont tous compris. Les trois quarts des élèves ont écouté l'exposé du médecin jusqu'au bout. 90 \% des habitants ont trouvé la clinique. Maintenant, les habitants de ce quartier tombent de moins en moins malades, et tout le monde est content !}

\textbf{Points de vigilance} :
\begin{itemize}
\item Vérifier l'ordre du pourcentage : 百分之七十 et jamais 七十百分之.
\item Vérifier l'ordre de la fraction : dénominateur + 分之 + numérateur (五分之四 = 4/5).
\item Placer 了 juste après le complément de résultat à l'affirmatif (找到了、听懂了、听完了) et utiliser 没(有) sans 了 à la négation.
\item Ne pas oublier 的 entre la proportion et le nom : 百分之七十的家庭.
\item Compter au moins 60 caractères et 5 phrases ; conclure par un avis ou un conseil.
\end{itemize}
\end{corrige}

\newpage

\coursec{Fiche bilan}

\begin{popfigure}
\begin{tikzpicture}[mindmap, grow cyclic, every node/.style={concept, font=\small}, concept color=popblue, root concept/.append style={concept color=popblue, font=\small\bfseries}, level 1/.append style={level distance=3.2cm, sibling angle=51.4, font=\footnotesize}, level 2/.append style={level distance=2.2cm, font=\scriptsize}]
\node[root concept, align=center]{Leçon 14\\ Résultat et fractions}
  child[concept color=poporange]{ node[align=center]{Complément de résultat\\ V + R}
    child { node[align=center]{看懂 kàndǒng\\ 听懂 tīngdǒng} }
    child { node[align=center]{做完 zuòwán\\ 洗干净 xǐ gānjìng} }
  }
  child[concept color=popgreen]{ node[align=center]{Négation\\ 没(有) + V + R}
    child { node[align=center]{没听懂\\ méi tīngdǒng} }
  }
  child[concept color=poppurple]{ node[align=center]{Forme potentielle\\ V + 得/不 + R}
    child { node {听得懂 / 听不懂} }
    child { node {做不完 zuò bu wán} }
  }
  child[concept color=poppink]{ node[align=center]{Fractions\\ 分母 + 分之 + 分子}
    child { node[align=center]{三分之一\\ 1/3} }
    child { node[align=center]{五分之二\\ 2/5} }
  }
  child[concept color=popgold]{ node[align=center]{Pourcentages\\ 百分之 + X}
    child { node[align=center]{百分之八十\\ 80 \%} }
  }
  child[concept color=popgreenL]{ node {Emplois}
    child { node[align=center]{Santé, environnement\\ statistiques} }
  }
  child[concept color=poporangeL]{ node {Pièges}
    child { node[align=center]{Pas de 了 après R\\ ordre inversé du français} }
  };
\end{tikzpicture}
\legende{Carte mentale de la leçon : le complément de résultat et l'expression des fractions.}
\end{popfigure}

\begin{bilan}
\textbf{1. Lexique clé à connaître par cœur}
\begin{center}
\begin{tabularx}{\linewidth}{|X|X|X|X|}
\hline
\rowcolor{popblueL} Caractères & Pinyin & Nature & Traduction \\
\hline
懂 & dǒng & verbe (résultat) & comprendre \\
\hline
完 & wán & verbe (résultat) & finir, terminer \\
\hline
见 & jiàn & verbe (résultat) & percevoir (voir/entendre) \\
\hline
干净 & gānjìng & adjectif (résultat) & propre \\
\hline
清楚 & qīngchu & adjectif (résultat) & clair \\
\hline
错 & cuò & adjectif (résultat) & faux, erroné \\
\hline
分之 & fēn zhī & mot-outil & « sur » (fractions) \\
\hline
百分之 & bǎi fēn zhī & mot-outil & pour cent \\
\hline
污染 & wūrǎn & nom / verbe & pollution / polluer \\
\hline
病人 & bìngrén & nom & malade, patient \\
\hline
\end{tabularx}
\end{center}

\textbf{2. Règles de grammaire à connaître par cœur}
\begin{center}
\begin{tabularx}{\linewidth}{|X|X|X|}
\hline
\rowcolor{popblueL} Structure & Exemple (caractères + pinyin) & Traduction \\
\hline
Sujet + V + R (+ complément) & 我看懂了。Wǒ kàndǒng le. & J'ai compris (en lisant). \\
\hline
Négation : 没(有) + V + R & 我没听懂。Wǒ méi tīngdǒng. & Je n'ai pas compris (en écoutant). \\
\hline
Potentiel + : V + 得 + R & 我听得懂。Wǒ tīng de dǒng. & Je peux comprendre. \\
\hline
Potentiel -- : V + 不 + R & 我做不完。Wǒ zuò bu wán. & Je ne peux pas finir. \\
\hline
Fraction : 分母 + 分之 + 分子 & 三分之一 sān fēn zhī yī & un tiers (1/3) \\
\hline
Pourcentage : 百分之 + X & 百分之六十 bǎi fēn zhī liùshí & 60 \% \\
\hline
\end{tabularx}
\end{center}

\textbf{3. Méthodes express}
\begin{itemize}
\item Pour traduire « réussir à faire » : choisir le bon résultat (懂 comprendre, 完 finir, 见 percevoir, 对 juste, 错 faux) et le coller juste après le verbe, sans 了 entre les deux.
\item Pour une fraction : lire d'abord le \cle{dénominateur}, puis 分之, puis le \cle{numérateur} — ordre inverse du français (« deux cinquièmes » = 五分之二).
\item Pour un pourcentage : toujours 百分之 + le nombre, jamais 分 + nombre.
\item Au potentiel, 得 = capacité positive, 不 = impossibilité ; la négation simple du résultat se fait avec 没(有).
\end{itemize}
\end{bilan}

\begin{aretenir}[Checklist avant le Bac]
\begin{itemize}
\item[$\square$] Je sais former un complément de résultat : verbe + résultat (看懂, 做完, 洗干净).
\item[$\square$] Je sais nier avec 没(有) + V + R : 我没看见。Wǒ méi kànjiàn.
\item[$\square$] Je sais former le potentiel affirmatif avec 得 : 听得懂 tīng de dǒng.
\item[$\square$] Je sais former le potentiel négatif avec 不 : 做不完 zuò bu wán.
\item[$\square$] Je ne mets jamais 了 entre le verbe et son complément de résultat.
\item[$\square$] Je place le complément d'objet après le groupe V + R : 我看完了这本书。
\item[$\square$] Je lis une fraction dans le bon ordre : dénominateur + 分之 + numérateur (四分之三 = 3/4).
\item[$\square$] Je sais dire les pourcentages : 百分之九十 = 90 \%.
\item[$\square$] Je sais utiliser ces structures dans une phrase sur la santé ou l'environnement : 百分之七十的病人看懂了医生的话。
\item[$\square$] Je connais les résultats fréquents : 懂, 完, 见, 对, 错, 干净, 清楚.
\item[$\square$] Je sais comparer avec le français : le chinois dit le moyen (看) puis le résultat (懂), le français souvent l'inverse (« comprendre en lisant »).
\item[$\square$] Je vérifie toujours l'ordre des mots avant de rendre ma copie.
\end{itemize}
\end{aretenir}

\findecours

\end{document}
